% Toolbox appendix aligned with Rectangle_CLT_Main-reworking-v3.tex,
% Updated 30 October 2026. The standalone source embeds this body; this file can
% also be included after Appendix F of the main manuscript with
%   \def\toolboxinpaper{}
%   % Toolbox appendix aligned with Rectangle_CLT_Main-reworking-v3.tex,
% Updated 30 October 2026. The standalone source embeds this body; this file can
% also be included after Appendix F of the main manuscript with
%   \def\toolboxinpaper{}
%   % Toolbox appendix aligned with Rectangle_CLT_Main-reworking-v3.tex,
% Updated 30 October 2026. The standalone source embeds this body; this file can
% also be included after Appendix F of the main manuscript with
%   \def\toolboxinpaper{}
%   % Toolbox appendix aligned with Rectangle_CLT_Main-reworking-v3.tex,
% Updated 30 October 2026. The standalone source embeds this body; this file can
% also be included after Appendix F of the main manuscript with
%   \def\toolboxinpaper{}
%   \input{Rectangle_CLT_Toolbox_appendix-reworking-v3}
% Fallback paper numbers below are checked against the v3 auxiliary file.

\makeatletter
\providecommand{\E}{\mathbb E}
\providecommand{\Pp}{\mathbb P}
\providecommand{\R}{\mathbb R}
\providecommand{\tr}{\operatorname{tr}}
\providecommand{\op}{\mathrm{op}}
\providecommand{\norm}[1]{\left\lVert#1\right\rVert}
\providecommand{\ind}{\mathbf 1}
\providecommand{\dd}{\,d}
\providecommand{\Hh}{H_2}
\providecommand{\surf}{\mathcal J}
\providecommand{\relint}{\operatorname{relint}}
\providecommand{\aff}{\operatorname{aff}}
\providecommand{\cone}{\operatorname{cone}}
\providecommand{\dist}{\operatorname{dist}}
\providecommand{\prj}{\operatorname{pr}}
\providecommand{\Haus}{\mathcal H}
\providecommand{\TBc}[1]{[#1]}
% \PN{label}{C.5} prints a number of the paper (a \ref in the paper);
% \PR{Lemma}{label}{C.5} prints "Lemma C.5"; \PE{label}{C.21} an equation.
\providecommand{\PN}[2]{\ifdefined\toolboxinpaper\ref{#1}\else#2\fi}
\providecommand{\PR}[3]{#1~\PN{#2}{#3}}
\providecommand{\PE}[2]{(\PN{#1}{#2})}
\providecommand{\usedin}[1]{\par\smallskip\noindent{\small\emph{Used in:} #1}\par\smallskip}
\@ifundefined{c@fact}{\newtheorem{fact}{Fact}[section]}{}
\makeatother

\section{Toolbox: facts from analysis and geometry}\label{app:toolbox}

The late interpolation estimate reduces high derivatives to two objects:
a Gaussian matrix weighted by the logistic approximation, and a signed
Gaussian integral over a polyhedral boundary. This appendix supplies the
analysis and geometry behind those reductions. The calculus facts also
justify the early estimate and the smoothing argument in
Appendices~\PN{appx-sec:smoothing}{B} and~\PN{appx-sec:support}{C}.
Each fact is stated with the assumptions needed for its use, followed by
a proof or a standard reference and a precise pointer to the manuscript.

We first treat differentiation, Gaussian integration by parts, and
covariance contraction. Log-concavity then controls the Gaussian law
tilted by a logistic weight. The geometric facts explain how to pair the
two contributions at each ridge and compare the resulting boundary
integral with the mass of short outward tubes. Because the faces are
flat, the required area and tube formulas follow from orthogonal
coordinates and Fubini's theorem. Finally, the first-variation formula
connects the logistic mixture to the surface estimate, including the
exceptional levels at which the polyhedron loses dimension.
Subsection~\ref{tb:sec:index} records these dependencies in one table.

\subsection{Notation}\label{tb:sec:notation}

The smoothing and interpolation arguments take place in the covariance
support \(\R^r\). They use a fixed list of \(m=2p\) unit face normals
\(u_1,\ldots,u_m\), which may repeat, and the threshold sets
\[
 A_a\coloneqq\{x\in\R^r:u_\alpha\cdot x\le a_\alpha,
                          \ \alpha=1,\ldots,m\},\qquad a\in\R^m.
\]
All thresholds are finite; inconsistent inequalities give the empty set.
We keep the full list when defining the logistic approximation, even if
some inequalities are redundant. We delete redundant inequalities only
when describing the facets of a fixed full-dimensional polyhedron.
For a specified system of unit-normal inequalities, \(Q^s\) increases
each threshold by \(s\); this notation depends on the system and does
not denote a Euclidean neighborhood. We write \(L\coloneqq\log(ep)\).

In the general facts below, \(\R^n\) carries the inner product
\(x\cdot y\) and Euclidean norm \(\norm x\). This dimension is \(r\)
in the support, or \(r-1\) when we work within a facet. We write \(dx\) for Lebesgue measure, \(\varphi_n\) and
\(\Phi_n\) for the standard Gaussian density and law on \(\R^n\), and \(G\)
for a standard Gaussian vector. For a linear subspace \(L\), \(\varphi_L\)
is the standard Gaussian density on \(L\) with respect to Lebesgue measure
on \(L\). In dimension zero, Lebesgue and Gaussian measure both assign
mass one to the origin, and \(\varphi_0(0)\coloneqq1\). This convention
includes the tangential calculation at an endpoint when \(r=1\).
For a nonempty set \(S\), \(\ind_S\) is its indicator function, \(\aff S\)
is the smallest affine subspace containing it, the dimension of \(S\) is
that of \(\aff S\), and \(\relint S\) is the interior of \(S\) relative to
\(\aff S\). The distance from a point \(y\) to a set \(S\) is
\(\dist(y,S)\coloneqq \inf_{x\in S}\norm{y-x}\).

For an integer \(q\ge1\) and \(f\in C^q(\R^n)\), the derivative of order \(q\) in the directions
\(z_1,\ldots,z_q\in\R^n\) is
\[
 D^qf(x)[z_1,\ldots,z_q]
 \coloneqq \partial_{s_1}\cdots\partial_{s_q}f(x+s_1z_1+\cdots+s_qz_q)\Big|_{s=0}
 =\sum_{i_1,\ldots,i_q}\partial_{i_1\cdots i_q}f(x)\,
   z_{1,i_1}\cdots z_{q,i_q},
\]
which is symmetric and multilinear in the directions. We write
\(D^qf(x)[z^{\otimes q}]\) when all directions equal \(z\), with
\(D^0f\coloneqq f\). The first derivative satisfies
\(Df(x)[z]=\nabla f(x)\cdot z\), and we identify the second derivative
with its Hessian matrix. The Loewner order \(M\preceq N\) means that
\(N-M\) is positive semidefinite, \(\norm\Xi_\op\) is the operator norm
of a matrix \(\Xi\), and \(\norm M_*\) is the \emph{nuclear norm}, which for
a symmetric matrix \(M\) with eigenvalues \(\mu_i\) equals \(\sum_i|\mu_i|\).

\subsection{Integration and calculus}\label{tb:sec:calc}

We repeatedly condition on an independent component and then average the
resulting bound. The following theorem specifies when these operations,
and the corresponding exchanges of integrals or sums, are valid.

\begin{fact}[Fubini and Tonelli]\label{tb:fubini}
Let \((S_1,\mu_1)\) and \((S_2,\mu_2)\) be \(\sigma\)-finite measure
spaces and let \(f\) be a measurable function on \(S_1\times S_2\). If
\(f\ge0\), then \(s_1\mapsto\int f(s_1,s_2)\,\mu_2(ds_2)\) is measurable
and
\[
 \int\Bigl(\int f\,d\mu_2\Bigr)d\mu_1
 =\int\Bigl(\int f\,d\mu_1\Bigr)d\mu_2
 =\int f\,d(\mu_1\otimes\mu_2)\in[0,\infty].
\]
If \(\int|f|\,d(\mu_1\otimes\mu_2)<\infty\), the same equalities hold, and
the inner integrals are finite almost everywhere.
\end{fact}

Three special cases occur in the proofs. If \(X\) and \(Y\) are
independent random elements and either \(f\ge0\) or \(f(X,Y)\) is
integrable, then
\[
 \E f(X,Y)=\E g(X),\qquad g(x)\coloneqq \E f(x,Y),
\]
because the joint law is the product of the marginal laws. This is the
formal content of every argument that is carried out for a fixed value of
an independent component such as \(V_0\) and then averaged over it. With
\(\mu_2\) the counting measure, the theorem says that
\(\E\sum_k\xi_k=\sum_k\E\xi_k\) for nonnegative terms, and for terms with
\(\sum_k\E|\xi_k|<\infty\). For an orthogonal decomposition
\(\R^n=L\oplus L^\perp\), Lebesgue measure on \(\R^n\) is the product of the
Lebesgue measures on \(L\) and on \(L^\perp\), so that
\(\int_{\R^n}f\,dx=\int_{L^\perp}\int_Lf(l+w)\,dl\,dw\).

The final averaging in \PR{Lemma}{appx-an:lem:surface}{C.5} is an
instance of the first case. Let \(V_0\in\R^r\) and
\(G\sim N(0,I_r)\) be independent, and let \(\sigma>0\).
For a fixed threshold set \(Q\), it gives
\[
 \E\!\left[R(G)^3\ind\{V_0+\sigma G\in
                   Q^{\sigma/R(G)}\setminus Q\}\right]
 =\E_G\!\left[R(G)^3\Pp_{V_0}\{V_0+\sigma G\in
                   Q^{\sigma/R(G)}\setminus Q\}\right].
\]
Here \(R\) is the projection radius in
Fact~\ref{tb:projection-moments}, and \(Q^s\) shifts every
unit-normal threshold outward by \(s\). Conditional on \(G\), both
the width \(\sigma/R(G)\) and the translation \(\sigma G\) are fixed.
Concentration of \(V_0\) can therefore be applied directly. Tonelli
requires no integrability check at this stage because the integrand
is nonnegative.

\usedin{conditioning on independent components in Appendices~\PN{appx-sec:smoothing}{B}
and~\PN{appx-sec:support}{C}; absolute convergence of the translated-weight
series in \PR{Lemma}{appx-an:lem:restoration}{C.2}; final averaging in
\PR{Lemma}{appx-an:lem:surface}{C.5}; conditioning on independent pairs
in the retained proof in \PR{Appendix}{appx-app:concentration}{D}.}
\emph{Source.} \TBc{Fol99}, Section~2.5.

To differentiate the interpolation, we first justify moving a derivative
inside an expectation and then expand along a line segment. An integrable
bound is needed for the former operation; the latter is ordinary
one-dimensional Taylor expansion applied to a multivariate function.

\begin{fact}[Differentiation under the integral sign; Taylor's formula]
\label{tb:taylor}
(a) Let \(I\) be an open interval and let \(f(t,\omega)\) be
measurable and integrable in \(\omega\) for each \(t\in I\).
Suppose that, outside one null set, it is differentiable in \(t\).
On each compact subinterval \(J\subset I\), assume
\(|\partial_t f(t,\omega)|\le g_J(\omega)\) for all \(t\in J\),
where \(g_J\) is integrable. Then \(\frac d{dt}\E f(t,\cdot)
=\E\,\partial_t f(t,\cdot)\) on \(I\).

(b) For an integer \(q\ge1\), \(f\in C^q(\R^n)\), and \(x,v\in\R^n\),
\[
 f(x+v)=\sum_{k=0}^{q-1}\frac1{k!}D^kf(x)[v^{\otimes k}]
 +\int_0^1\frac{(1-s)^{q-1}}{(q-1)!}D^qf(x+sv)[v^{\otimes q}]\,ds.
\]
For \(q=1\) this is \(f(x+v)-f(x)=\int_0^1Df(x+sv)[v]\,ds\).
\end{fact}

\begin{proof}
For (a), fix an interior time and a compact neighborhood \(J\).
The mean value theorem bounds nearby difference quotients by \(g_J\),
and they converge pointwise, so dominated convergence applies. For (b),
the function \(s\mapsto f(x+sv)\) is of class \(C^q\) with \(k\)th derivative
\(D^kf(x+sv)[v^{\otimes k}]\); apply the one-dimensional Taylor formula with
integral remainder at \(s=0\) and evaluate it at \(s=1\).
\end{proof}

\usedin{the signed Hessian comparison and its Taylor expansion in
Subsection~\PN{appx-an:sec:path}{C.2}; the restoration of the cubic term
to the full path in \PR{Lemma}{appx-an:lem:identity}{C.1}; the segment
integral in \PR{Lemma}{appx-lem:endpoint-gap}{B.3}.}
\emph{Source.} \TBc{Fol99}, Section~2.3.

Differentiating the thinning probability changes one mask at a time.
Here a mask is a variable that retains a summand when it equals one
and deletes it when it equals zero. The following finite-product identity
isolates the change in retention probability; changes in amplitude and
Gaussian variance are separate contributions.

\begin{fact}[Differentiating independent Bernoulli masks]\label{tb:bernoulli}
Let \(B_1,\ldots,B_N\) be independent Bernoulli variables with parameter
\(q\in(0,1)\). For a function \(H:\{0,1\}^N\to\R\), write
\(B^{i,b}\) for the vector obtained by setting its \(i\)th entry to
\(b\in\{0,1\}\). Then
\[
 \frac{d}{dq}\E_q H(B)
 =\sum_{i=1}^N\E_q\{H(B^{i,1})-H(B^{i,0})\}.
\]
If \(q=q(t)\) and \(H=H_t\) are differentiable, then
\[
 \frac{d}{dt}\E_{q(t)}H_t(B)
 =\E_{q(t)}\partial_t H_t(B)
 +q'(t)\sum_i\E_{q(t)}
       \{H_t(B^{i,1})-H_t(B^{i,0})\}.
\]
\end{fact}

\begin{proof}
Temporarily give the masks separate parameters \(q_1,\ldots,q_N\).
Conditioning on all masks except the \(i\)th makes their expectation
affine in \(q_i\), with derivative equal to the insertion--deletion
difference in the first display. Differentiating along
\(q_1=\cdots=q_N=q\) gives the sum. The second formula is the chain
rule. Independent random arguments may be included by conditioning
and then applying Fact~\ref{tb:taylor}(a).
\end{proof}

\usedin{Subsection~\PN{appx-an:sec:path}{C.2}, before
\PR{Lemma}{appx-an:lem:identity}{C.1}, to differentiate the masks in
the Bernoulli--Gauss path.}

In the paper, the amplitude and retention probability are
\(\gamma=1+t\) and \(\varpi=(1-t)/(1+t)^2\), for \(0<t<1\).
The support-coordinate summands \(\zeta_i^0\) are centered, with
\(\sum_i\E\zeta_i^0(\zeta_i^0)^\top=I_r\). Write
\(Y_i\coloneqq\gamma\mathbb B_{i,t}\zeta_i^0\) and
\(U_t\coloneqq\sum_iY_i+\sqrt t\,G\), where all masks and \(G\)
are mutually independent and independent of the summands. To expand the
contribution of summand \(i\), remove it and put \(H_i\coloneqq U_t-Y_i\).
Let \(\zeta_i\) be an independent copy of \(\zeta_i^0\), with all copies
independent of each other and of the path variables. The signed Hessian
comparison in~\PE{appx-an:signed-hessian}{C.6} has scalar weight
\[
 \nu_t(s)\coloneqq \frac{4-2(3-t)s}{1+t},\qquad 0\le s\le1.
\]
Direct integration gives
\[
 \int_0^1\nu_t(s)\,ds=1,\qquad
 \int_0^1s\nu_t(s)\,ds=\frac{2t}{3(1+t)}.
\]
For a twice continuously differentiable scalar function \(g\),
Taylor's formula and the first identity give
\[
 g(0)-\int_0^1\nu_t(s)g(s)\,ds
 =-\frac{2t}{3(1+t)}g'(0)
   -\int_0^1\nu_t(s)\int_0^s(s-u)g''(u)\,du\,ds.
\]
For the logistic approximation \(F=F_{a,h}\) defined below, take
\(g(s)\coloneqq D^2F(H_i+s\gamma\zeta_i)[\zeta_i,\zeta_i]\).
Then \(g'(0)=\gamma D^3F(H_i)[\zeta_i^{\otimes3}]\), so
the factor \(1/2\) in the signed Hessian comparison turns the first
term into
\(-\tfrac12\{2t/(3\gamma)\}\gamma
 D^3F(H_i)[\zeta_i^{\otimes3}]
 =-tD^3F(H_i)[\zeta_i^{\otimes3}]/3\). This explains the
cubic coefficient before the derivatives are restored to the full path.
The weight \(\nu_t\) is signed; these identities use its
integrals, not a probability interpretation.

Repeated differentiation of the logistic approximation produces products
of face probabilities. The partition form of the higher-order chain
rule records which directions differentiate the same inner factor.

\begin{fact}[Fa\`a di Bruno's formula]\label{tb:faa}
Let \(\phi\) be a smooth function on an open interval of \(\R\) and let
\(g:\R^n\to\R\) be smooth with values in that interval. For \(q\ge1\), let
\(\Pi_q\) be the set of partitions \(\mathcal P\) of \(\{1,\ldots,q\}\) into
nonempty blocks, let \(|\mathcal P|\) be the number of blocks, and for a
block \(J\) let \(z_J\) be the list of directions \(z_a\) with
\(a\in J\). Then
\[
 D^q(\phi\circ g)(x)[z_1,\ldots,z_q]
 =\sum_{\mathcal P\in\Pi_q}\phi^{(|\mathcal P|)}(g(x))
  \prod_{J\in\mathcal P}D^{|J|}g(x)[z_J].
\]
\end{fact}

\begin{proof}
For \(q=1\) this is the chain rule. Differentiate the right side for \(q\)
in the direction \(z_{q+1}\). By the product rule, the derivative either
falls on \(\phi^{(|\mathcal P|)}(g(x))\), which gives
\(\phi^{(|\mathcal P|+1)}(g(x))\,Dg(x)[z_{q+1}]\) and adds the block
\(\{q+1\}\) to \(\mathcal P\), or falls on one factor
\(D^{|J|}g(x)[z_J]\), which adds \(q+1\) to the block \(J\). Every partition
of \(\{1,\ldots,q+1\}\) arises exactly once in this way, by removing
\(q+1\) from its block.
\end{proof}

\emph{Source.} \TBc{Com74}; the induction above proves the form used here.

For a finite list of face normals \(u_1,\ldots,u_m\in\R^n\), thresholds
\(a_1,\ldots,a_m\in\R\), and a width \(h>0\), put
\[
 q_\alpha(x)\coloneqq(u_\alpha\cdot x-a_\alpha)/h,\qquad
 w_\alpha(x)\coloneqq e^{q_\alpha(x)},\qquad T(x)\coloneqq\sum_\alpha w_\alpha(x),
\]
\[
 F(x)\coloneqq\frac1{1+T(x)},\qquad
 \pi_\alpha(x)\coloneqq F(x)w_\alpha(x),\qquad
 \theta(x)\coloneqq1-F(x)=\sum_\alpha\pi_\alpha(x).
\]
Here \(F\) is the logistic approximation to the indicator of the
threshold set. The baseline outcome has weight one and probability
\(F\); face \(\alpha\) has probability \(\pi_\alpha\). These
probabilities describe the smoothing formula, not the original event.
One and two differentiations give
\begin{align*}
 DF(x)[z]&=-h^{-1}\sum_\alpha F\pi_\alpha(u_\alpha\cdot z),\\
 D^2F(x)[z_1,z_2]
 &=h^{-2}\left\{2\sum_{\alpha,\beta}F\pi_\alpha\pi_\beta
       (u_\alpha\cdot z_1)(u_\beta\cdot z_2)
       -\sum_\alpha F\pi_\alpha
       (u_\alpha\cdot z_1)(u_\alpha\cdot z_2)\right\}.
\end{align*}
All weights are evaluated at \(x\). In the first formula their sum is
\(F\theta\), which is small away from the smoothed boundary.
The next result retains this factor at every derivative order.

\begin{fact}[Derivatives of the logistic approximation]\label{tb:logistic}
With the notation above, let \(q\ge1\).
For directions \(z_1,\ldots,z_q\), with the blocks \(J_1,\ldots,J_{|\mathcal P|}\)
of \(\mathcal P\) listed in any fixed order,
\[
 D^qF(x)[z_1,\ldots,z_q]
 =h^{-q}\sum_{\mathcal P\in\Pi_q}(-1)^{|\mathcal P|}|\mathcal P|!
 \sum_{\alpha_1,\ldots,\alpha_{|\mathcal P|}}
 F(x)\prod_{b=1}^{|\mathcal P|}\pi_{\alpha_b}(x)
 \prod_{b=1}^{|\mathcal P|}\prod_{a\in J_b}(u_{\alpha_b}\cdot z_a),
\]
Moreover, \(\sum_{\mathcal P\in\Pi_q}|\mathcal P|!\le q!\,2^{q-1}\), and hence
\[
 |D^qF(x)[z_1,\ldots,z_q]|
 \le q!\,2^{q-1}h^{-q}F(x)\theta(x)
       \prod_{a=1}^q\max_\alpha|u_\alpha\cdot z_a|.
\]
\end{fact}

\begin{proof}
Apply Fact~\ref{tb:faa} with \(\phi(\tau)\coloneqq (1+\tau)^{-1}\) and \(g=T\). Then
\(\phi^{(k)}(\tau)=(-1)^kk!\,(1+\tau)^{-k-1}\), and since each \(w_\alpha\) is
an exponential of an affine function,
\(D^lT(x)[z_1,\ldots,z_l]=h^{-l}\sum_\alpha w_\alpha(x)\prod_{a\le l}
(u_\alpha\cdot z_a)\). Substituting, and using
\((1+T)^{-|\mathcal P|-1}\prod_bw_{\alpha_b}=F\prod_b\pi_{\alpha_b}\) and
\(\sum_b|J_b|=q\), gives the formula. The number
\(\sum_{\mathcal P}|\mathcal P|!\) counts ordered partitions
\((J_1,\ldots,J_k)\) of \(\{1,\ldots,q\}\). Writing the elements of
\(J_1\), then of \(J_2\), and so on, in increasing order inside each block,
gives a permutation together with the set of positions at which a new block
starts, which is a subset of the \(q-1\) gaps between consecutive entries.
Every ordered partition arises from at least one of these
\(q!\,2^{q-1}\) pairs. For a partition with \(k\ge1\) blocks, summing the probability
weights gives \(F\theta^k\le F\theta\). Each product of projections
is bounded by \(\prod_a\max_\alpha|u_\alpha\cdot z_a|\).
These two observations prove the stated envelope without a factor
depending on the dimension.
\end{proof}

The early estimate uses this formula for orders three and four. After
the two Gaussian integrations by parts in the late estimate, the explicit
first and second derivatives above suffice. In both cases the factor
\(F\theta\) connects the derivative estimate to concentration near
the boundary.

\usedin{the general derivative expansion and local envelope in
Subsection~\PN{appx-sec:mixture}{B.2}, namely
\PE{appx-eq:partition-derivative}{B.1}--\PE{appx-eq:local-envelope}{B.3};
the explicit first and second derivatives in
\PR{Lemma}{appx-an:lem:restoration}{C.2}.}
\emph{Source.} The formula follows from the partition chain rule in
Fact~\ref{tb:faa}; the proof above is self-contained.

\subsection{Gaussian identities}\label{tb:sec:gauss}

The mean estimate separates a normal coordinate from the tangential
coordinates of a face. Orthogonality makes these Gaussian components
independent; a second elementary identity controls the density change
when we move outward from the face.

\begin{fact}[Gaussian density in orthogonal coordinates]\label{tb:gauss-coord}
(a) Let \(\R^n=L\oplus L^\perp\) be an orthogonal decomposition. Then
\(\varphi_n(l+w)=\varphi_L(l)\varphi_{L^\perp}(w)\) for \(l\in L\) and
\(w\in L^\perp\), so the projections of a standard Gaussian vector onto
\(L\) and \(L^\perp\) are independent standard Gaussian vectors. In
particular, for a unit vector \(u\), \(a\in\R\), and \(y\perp u\),
\(\varphi_n(au+y)=\varphi_1(a)\varphi_{n-1}(y)\).

(b) For \(x,v\in\R^n\),
\(\varphi_n(x+v)=\varphi_n(x)\exp\{-x\cdot v-\tfrac12\norm v^2\}\).
\end{fact}

\begin{proof}
Both follow from \(\norm{l+w}^2=\norm l^2+\norm w^2\) and from expanding
\(\norm{x+v}^2\).
\end{proof}

\usedin{the normal--tangential decomposition in
\PR{Lemma}{appx-an:lem:mean}{C.3}; the facet density factorization
and density comparison along normal tubes in
\PR{Lemma}{appx-an:lem:surface}{C.5}.}

Gaussian integration by parts replaces two directional derivatives by
a centered quadratic function of the Gaussian vector. The factor
\(\sigma^{-2}\) below records the Gaussian noise scale exactly.

\begin{fact}[Gaussian integration by parts]\label{tb:stein}
Let \(G\sim N(0,I_n)\). In (a), assume \(f\in C^1(\R^n)\) and
that \(f\) and its first derivatives have polynomial growth.
In (b), assume in addition that \(f\in C^2\) and its second
derivatives have polynomial growth. Then
\begin{enumerate}
\item[(a)] \(\E[Gf(G)]=\E\nabla f(G)\);
\item[(b)] \(\E[(GG^\top-I_n)f(G)]=\E D^2f(G)\); equivalently, for \(z\in\R^n\),
\(\E D^2f(G)[z,z]=\E\bigl[\tr\{(GG^\top-I_n)zz^\top\}f(G)\bigr]\), where
\(\tr\{(GG^\top-I_n)zz^\top\}=(z\cdot G)^2-\norm z^2\);
\item[(c)] if \(s\ge0\) is an integer, \(v\in\R^n\), \(\sigma>0\),
\(U\coloneqq v+\sigma G\), and
\(F\in C^{s+2}(\R^n)\) has derivatives of polynomial growth, then for
\(z,z_1,\ldots,z_s\in\R^n\),
\[
 \E D^{s+2}F(U)[z,z,z_1,\ldots,z_s]
 =\sigma^{-2}\E\bigl[\tr\{(GG^\top-I_n)zz^\top\}\,
 D^sF(U)[z_1,\ldots,z_s]\bigr].
\]
\end{enumerate}
\end{fact}

\begin{proof}
For (a) in one coordinate, \(x\varphi_1(x)=-\varphi_1'(x)\), so integration by
parts gives
\[\int xf(x)\varphi_1(x)\,dx=\int f'(x)\varphi_1(x)\,dx,\]
and the boundary terms vanish because of the growth condition. The other coordinates
are handled by Fact~\ref{tb:fubini}. For (b), apply (a) to \(\partial_if\) and
to \(G_kf\): \(\E[G_iG_kf(G)]=\E\partial_i(G_kf)(G)=\delta_{ik}\E f(G)
+\E[G_k\partial_if(G)]=\delta_{ik}\E f(G)+\E\partial_{ik}f(G)\). For (c), the
function \(f(g)\coloneqq D^sF(v+\sigma g)[z_1,\ldots,z_s]\) has
\(D^2f(g)[z,z]=\sigma^2D^{s+2}F(v+\sigma g)[z,z,z_1,\ldots,z_s]\); apply (b).
\end{proof}

\usedin{Subsection~\PN{sec:outline-late}{2.4} and
\PR{Lemma}{appx-an:lem:restoration}{C.2}, conditional on \(\mathcal H\)
with \(v=V_0+v_i\) and \(z=Z_i\); the identity
\(\E[Z\cdot\nabla\Psi(Z)]=\E\tr(\Sigma D^2\Psi(Z))\)
in \PR{Lemma}{appx-lem:endpoint-gap}{B.3}.}
The identity also holds conditionally when the base and directions are
measurable with respect to a sigma-field independent of \(G\), provided
the conditional growth assumptions hold. Averaging that identity further
requires integrability of both sides. In the paper the fixed width
\(h>0\), bounded face components, and conditional covariance bound
supply this integrability.

\emph{Source.} Stein~\TBc{Ste81}; the proof above is complete.

Changing the Gaussian variance contributes a Hessian trace to the path
derivative. We give the formula for a general covariance and then use
its standard-Gaussian form in the support coordinates.

\begin{fact}[Heat equation for the Gaussian component]\label{tb:heat}
Let \(Z\sim N(0,\Sigma)\) in \(\R^p\), where \(\Sigma\) may be singular, and
let \(f\in C^2(\R^p)\) be bounded with bounded first and second derivatives.
Then, for \(x\in\R^p\) and \(t>0\),
\[
 \frac d{dt}\E f(x+\sqrt t Z)
 =\tfrac12\E\tr\{\Sigma D^2f(x+\sqrt t Z)\}.
\]
\end{fact}

\begin{proof}
Write \(Z=\Gamma G\) with \(\Gamma\Gamma^\top=\Sigma\) and \(G\) standard
Gaussian. By Fact~\ref{tb:taylor}(a), which applies because the difference
quotients near \(t_0>0\) are bounded by
\(\norm\nabla f_\infty\norm{\Gamma G}/(2\sqrt{t_0/2})\)
when the times lie in \([t_0/2,3t_0/2]\), the derivative equals
\((2\sqrt t)^{-1}\E[G\cdot\Gamma^\top\nabla f(x+\sqrt t\Gamma G)]\).
Fact~\ref{tb:stein}(a) applied to the function \(g\mapsto
\Gamma^\top\nabla f(x+\sqrt t\Gamma g)\) turns this into
\((2\sqrt t)^{-1}\sqrt t\,\E\tr\{\Gamma^\top D^2f\,\Gamma\}
=\frac12\E\tr\{\Sigma D^2f\}\).
\end{proof}

In the support \(\R^r\), the covariance is \(I_r\), so this gives
\(\frac d{dt}\E_G F(v+\sqrt tG)=\frac12\E_G\tr D^2F(v+\sqrt tG)\).
The base \(v\) is held fixed for this differentiation. If its law also
depends on time, its change contributes separately.

\usedin{Subsection~\PN{sec:outline-path}{2.2} and the derivation of
\PE{appx-an:signed-hessian}{C.6} in
Subsection~\PN{appx-an:sec:path}{C.2}. The Gaussian component supplies
the covariance term that cancels the constant term of the Hessian
expansion.}

The surface estimate needs the size of a few normal projections,
not the Euclidean length of the full Gaussian vector. The following
bound explains why only a logarithm of the number of faces enters.

\begin{fact}[Maxima of one- and two-dimensional Gaussian projections]
\label{tb:projection-moments}
Let \(u_1,\ldots,u_m\) be unit vectors in \(\R^r\), with \(m\ge1\).
For linearly independent \(u_j,u_k\), let \(\Pi_{jk}\) be orthogonal
projection onto their span, and put
\[
 \mathcal M(g)\coloneqq \max\left\{\max_j|u_j\cdot g|,
   \max_{j,k:\,u_j,u_k\text{ independent}}\norm{\Pi_{jk}g}\right\},
 \qquad R(g)\coloneqq 1+\mathcal M(g).
\]
An empty second maximum is omitted. Then \(R\) is \(1\)-Lipschitz.
For \(G\sim N(0,I_r)\) and \(q\in\{2,3\}\),
\[
 \E R(G)^q\le C_q\{\log(em)\}^{q/2}.
\]
\end{fact}

\begin{proof}
Every projection is a contraction, so each function in the maximum
is \(1\)-Lipschitz, as is their maximum. A scalar standard normal
satisfies \(\Pp(|Z|>s)\le2e^{-s^2/2}\); the length of a
two-dimensional standard normal satisfies
\(\Pp(\norm{G_2}>s)=e^{-s^2/2}\). A union bound gives
\[
 \Pp\{\mathcal M(G)>s\}\le\min\{1,2(m+m^2)e^{-s^2/2}\}.
\]
The projections need not be independent. Integrate
\(q\int_0^\infty s^{q-1}\Pp\{\mathcal M(G)>s\}\,ds\), splitting
at \(\sqrt{8\log(em)}\), and use
\((1+x)^q\le2^{q-1}(1+x^q)\).
\end{proof}

\usedin{\PE{appx-an:projection-moments}{C.28} in
\PR{Lemma}{appx-an:lem:surface}{C.5}. For the full face list \(m=2p\), and also for its subsystems, these moments
are of orders \(L\) and \(L^{3/2}\), where \(L=\log(ep)\).}

\subsection{Matrix and entropy inequalities}\label{tb:sec:matrix}

Both time ranges use the same observation: the summed second moments
control a pair of projections even when other random coefficients
remain inside the expectation.

\begin{fact}[Conditional covariance contraction]\label{tb:covariance}
Let \(Z_1,\ldots,Z_N\in\R^r\) have finite second moments, and let
\(\mathcal H\) be a sigma-field such that
\(\sum_i\E[Z_iZ_i^\top\mid\mathcal H]\preceq I_r\) almost surely.
For \(\mathcal H\)-measurable unit vectors \(u,v\),
\[
 \sum_i\E[|(u\cdot Z_i)(v\cdot Z_i)|\mid\mathcal H]\le1.
\]
If real random coefficients satisfy \(|c_i|\le a\) almost surely,
where \(a\ge0\) is constant, then
\[
 \left\|\sum_i\E[c_iZ_iZ_i^\top\mid\mathcal H]\right\|_\op\le a.
\]
No independence of the coefficients and vectors is required.
\end{fact}

\begin{proof}
Conditional Cauchy--Schwarz, followed by Cauchy--Schwarz over the
index \(i\), bounds the first sum by
\[
 \left(\sum_i\E[(u\cdot Z_i)^2\mid\mathcal H]\right)^{1/2}
 \left(\sum_i\E[(v\cdot Z_i)^2\mid\mathcal H]\right)^{1/2}\le1.
\]
For the second claim, each unit vector \(u\) gives
\(\bigl|\sum_i\E[c_i(u\cdot Z_i)^2\mid\mathcal H]\bigr|\le a\).
The matrix is symmetric, so taking the supremum of its absolute
quadratic forms proves the operator-norm bound.
\end{proof}

\usedin{\PE{appx-an:pair-contraction}{C.2} and the early estimates in
Subsection~\PN{appx-an:sec:early}{C.3}; the coefficient matrices in
\PR{Lemma}{appx-an:lem:restoration}{C.2}. Bounded face projections
control the remaining directions by powers of \(\delta\).}

After conditioning, the coefficient matrix has bounded operator norm.
The following duality transfers that bound to the Gaussian matrix, whose
nuclear norm we estimate using its positive and negative parts.

\begin{fact}[Duality between the operator and nuclear norms]\label{tb:nuclear}
(a) For every square matrix \(\Xi\) and every symmetric matrix \(M\),
\(|\tr(\Xi M)|\le\norm\Xi_\op\norm M_*\). Moreover
\(\norm M_*=\max\{\tr(\Xi M):\norm\Xi_\op\le1\}\), so that \(\norm\cdot_*\) is
a norm on symmetric matrices.

(b) If \(M=P-N\) with \(P\succeq0\) and \(N\succeq0\), then
\(\norm M_*\le\tr P+\tr N=2\tr P-\tr M\).
\end{fact}

\begin{proof}
For (a), let \(M=\sum_i\mu_ie_ie_i^\top\) be a spectral decomposition. Then
\(\tr(\Xi M)=\sum_i\mu_i\,e_i^\top\Xi e_i\) and \(|e_i^\top\Xi e_i|\le
\norm\Xi_\op\), which gives the inequality. Equality is attained at
\(\Xi\coloneqq \sum_i\operatorname{sgn}(\mu_i)e_ie_i^\top\). A supremum of linear
functionals satisfies the triangle inequality. For (b), the norm of a
positive semidefinite matrix is its trace, so the triangle inequality gives
\(\norm M_*\le\tr P+\tr N\), and \(\tr N=\tr P-\tr M\).
\end{proof}

\usedin{the final contraction in \PR{Lemma}{appx-an:lem:restoration}{C.2};
the reduction \PE{appx-an:matrix-trace-reduction}{C.22} in
\PR{Lemma}{appx-an:lem:matrix}{C.4}, with
\(P=p_\alpha\mu_\alpha\mu_\alpha^\top\) and
\(N=p_\alpha(I_r-K_\alpha)\).}

An individual logistic weight can have arbitrarily small mass, so its
logarithmic bound should be summed before the mass is estimated. The
log-sum inequality performs that summation; its integral version then
allows us to average over the independent base.

\begin{fact}[Log-sum inequality and a concave function]\label{tb:logsum}
(a) For positive numbers \(y_1,\ldots,y_l\) and \(z_1,\ldots,z_l\) with sums
\(y\) and \(z\),
\(\sum_jy_j\log(z_j/y_j)\le y\log(z/y)\).

(b) The function \((x,y)\mapsto x\log(y/x)\) is concave on \((0,\infty)^2\).
Hence, for positive integrable random variables \(X,Y\),
\(\E[X\log(Y/X)]\le\E X\,\log(\E Y/\E X)\).
The expectation on the left may equal \(-\infty\); its positive
part is integrable because \(x\log(y/x)\le y/e\).
\end{fact}

\begin{proof}
For (a), Jensen's inequality for the concave function \(\log\), with the
weights \(y_j/y\), gives \(\sum_j(y_j/y)\log(z_j/y_j)\le\log\sum_j(y_j/y)(z_j/y_j)
=\log(z/y)\). For (b), the function \(f(x,y)\coloneqq x\log y-x\log x\) has Hessian
\(\begin{pmatrix}-1/x&1/y\\1/y&-x/y^2\end{pmatrix}\), whose determinant is
zero and whose trace is negative, so the Hessian is negative semidefinite.
For the expectation bound, put \(\mu\coloneqq\E X\) and
\(\eta\coloneqq\E Y\). The tangent-plane inequality for this concave
function gives
\[
 f(X,Y)\le f(\mu,\eta)
   +f_x(\mu,\eta)(X-\mu)+f_y(\mu,\eta)(Y-\eta).
\]
The right side is integrable. Taking expectations proves the claim,
including when the negative part of the left side has infinite mean.
\end{proof}

\usedin{the reference-mass estimate in
\PR{Lemma}{appx-an:lem:matrix}{C.4}: first sum over tuples sharing
their first face, then average over \(V_0\). Both the reference mass
and the logistic mass contain the factor \(h\sqrt L\), which cancels
inside the logarithm.}
\emph{Source.} \TBc{CT06}, Theorem~2.7.1, for (a).

\subsection{Log-concavity}\label{tb:sec:logconcave}

A function \(f:\R^n\to[0,\infty)\) is \emph{log-concave} if
\(f((1-\alpha)x+\alpha y)\ge f(x)^{1-\alpha}f(y)^\alpha\) for all \(x,y\) and \(\alpha\in[0,1]\);
when \(f>0\) this says that \(\log f\) is concave. At the endpoints
\(\alpha=0,1\) the inequality means equality, without a convention
for \(0^0\). We need the fact that integrating out coordinates
preserves log-concavity: it will turn a Gaussian convolution into a
concave log normalizing constant.

\begin{fact}[Pr\'ekopa--Leindler inequality and Pr\'ekopa's theorem]
\label{tb:prekopa}
(a) Let \(0<\alpha<1\) and let \(f,g,h:\R^n\to[0,\infty)\) be bounded,
continuous, and integrable, with
\(h((1-\alpha)x+\alpha y)\ge f(x)^{1-\alpha}g(y)^\alpha\) for all \(x,y\). Then
\(\int h\ge(\int f)^{1-\alpha}(\int g)^\alpha\).

(b) Let \(f:\R^a\times\R^b\to[0,\infty)\) be bounded, continuous, and
log-concave, with \(f(z,\cdot)\) integrable for every \(z\). Then
\(z\mapsto\int f(z,y)\,dy\) is log-concave on \(\R^a\).
\end{fact}

\begin{proof}
Part (a) is the Pr\'ekopa--Leindler inequality; we use the standard
integrable-function theorem in \TBc{Gar02}, Theorem~7.1.
To deduce (b), fix \(z_0,z_1\) and \(0<\alpha<1\), and put
\(z_\alpha\coloneqq(1-\alpha)z_0+\alpha z_1\).
Log-concavity gives
\[
 f(z_\alpha,(1-\alpha)y_0+\alpha y_1)
 \ge f(z_0,y_0)^{1-\alpha}f(z_1,y_1)^\alpha.
\]
The three sections are bounded, continuous, and integrable by hypothesis,
so (a) applies to them. Its conclusion is exactly the log-concavity of
\(z\mapsto\int f(z,y)\,dy\). The endpoint cases are immediate.
\end{proof}

\emph{Source.} Pr\'ekopa~\TBc{Pre71,Pre73} and Leindler~\TBc{Lei72}.
The marginalization result is Theorem~6 of \TBc{Pre73};
\TBc{Gar02}, Section~7, treats the integral inequality.

We now apply marginal log-concavity to a Gaussian convolution. Its
logarithmic Hessian is the covariance of the tilted law minus the
identity, so concavity gives the matrix inequality needed in the proof.

\begin{fact}[Gaussian laws tilted by a log-concave weight]\label{tb:tilt}
Let \(w:\R^r\to[0,\infty)\) be bounded, continuous, and log-concave, with
\(p\coloneqq \E w(G)>0\) for \(G\sim N(0,I_r)\). Let \(\mu\) and \(K\) be the mean and
covariance matrix of the probability law with density \(w(g)\varphi_r(g)/p\).
Then \(K\preceq I_r\), and
\(\E[(GG^\top-I_r)w(G)]=p(K-I_r+\mu\mu^\top)\).
\end{fact}

\begin{proof}
The function \((z,g)\mapsto e^{-\norm{g-z}^2/2}w(g)\) is log-concave, since
\(\norm{g-z}^2\) is jointly convex in \((z,g)\), and it is bounded and
continuous, with an integrable \(g\)-section for every \(z\).
The integral is positive because \(p>0\) and the Gaussian kernel is
strictly positive. By Fact~\ref{tb:prekopa}(b), \(\ell(z)\coloneqq \log\int e^{-\norm{g-z}^2/2}
w(g)\,dg\) is concave. Since \(\norm{g-z}^2=\norm g^2-2z\cdot g+\norm z^2\),
\[
 \ell(z)=\text{const}-\tfrac12\norm z^2+\log M(z),\qquad
 M(z)\coloneqq \int e^{z\cdot g}w(g)\varphi_r(g)\,dg=p\,\E_\nu e^{z\cdot g},
\]
where \(\nu\) is the tilted law. The function \(M\) is finite and smooth: on every bounded set of
\(z\)'s, each derivative of the integrand is bounded by a constant
times \((1+\norm g)^k e^{c\norm g}\varphi_r(g)\), an integrable function.
Thus differentiation under the integral is valid, and the Hessian of \(\log M\) at \(0\) is the
covariance \(K\) of \(\nu\). A concave function of class \(C^2\) has a
negative semidefinite Hessian, so \(K-I_r\preceq0\). The identity is the
decomposition of the second moment \(\E_\nu gg^\top=K+\mu\mu^\top\).
\end{proof}

\usedin{the Gaussian-tilt calculation in
\PR{Lemma}{appx-an:lem:matrix}{C.4}, where \(w=w_\alpha\) and
\(K=K_\alpha\), and its explanation in
Subsection~\PN{sec:outline-late}{2.4}.}
\emph{Source.} The covariance contraction is a classical consequence of
Pr\'ekopa's marginalization theorem~\TBc{Pre73}, Theorem~6, as the
preceding proof shows. For smooth positive weights, it also follows
from the Brascamp--Lieb variance inequality~\TBc{BL76}, applied to
linear functions under the potential \(\norm g^2/2-\log w(g)\), whose
Hessian is at least \(I_r\). The same convolution argument covers
indicators of convex sets by the measurable form of Pr\'ekopa's
theorem. This convex-cell specialization is also used in \TBc{10}.

It remains to check that every weight produced by the derivative
expansion meets the preceding hypothesis. Its logarithm has a particularly
simple form: an affine function minus a multiple of a log-sum-exponential.

\begin{fact}[Log-concavity of the logistic weights]\label{tb:lc-weights}
With the notation of Fact~\ref{tb:logistic}, for every \(k\ge0\) and every
tuple \((\alpha_1,\ldots,\alpha_k)\), the function
\(F\pi_{\alpha_1}\cdots\pi_{\alpha_k}=\exp\{\sum_{a\le k}q_{\alpha_a}\}
(1+T)^{-k-1}\), where \(q_\alpha\coloneqq (u_\alpha\cdot x-a_\alpha)/h\), is
log-concave on \(\R^n\), and so is its composition
\(g\mapsto F\pi_{\alpha_1}\cdots\pi_{\alpha_k}(v+\sigma g)\) with any affine map.
\end{fact}

\begin{proof}
The logarithm is a sum of the affine functions \(q_{\alpha_a}\) and
\(-(k+1)\log(1+\sum_\alpha e^{q_\alpha})\). The function
\(\log(1+\sum_\alpha e^{q_\alpha})\) is convex. To include the baseline
term explicitly, put \(q_0\coloneqq0\). H\"older's inequality gives
\[
 \sum_{\alpha=0}^m e^{(1-s)q_\alpha(x)+s q_\alpha(y)}
 \le\left(\sum_{\alpha=0}^m e^{q_\alpha(x)}\right)^{1-s}
     \left(\sum_{\alpha=0}^m e^{q_\alpha(y)}\right)^s,
 \qquad 0<s<1.
\]
Taking logarithms proves the assertion. Composition
with an affine map preserves concavity.
\end{proof}

\usedin{the application of covariance contraction to every
\(w_\alpha=F\prod_a\pi_{\alpha_a}\) in
\PR{Lemma}{appx-an:lem:matrix}{C.4}.}
\emph{Source.} Convexity of log-sum-exp is standard; see
Boyd and Vandenberghe~\TBc{BV04}, Section~3.1.5, pp.~72--74.
The baseline term \(1=e^{q_0}\), with \(q_0=0\), is included in that sum.

\subsection{Polyhedra: faces, normal cones, and nearest points}\label{tb:sec:poly}

The surface estimate treats a fixed polyhedron geometrically. Unlike
the smoothing formula, its boundary integral depends only on the set,
so repeated and redundant inequalities may now be removed.
A \emph{polyhedron} is a set
\[
 Q\coloneqq \bigcap_{j=1}^m\{x\in\R^n:u_j\cdot x\le a_j\},
\]
with unit vectors \(u_j\), called \emph{face normals}, and numbers
\(a_j\), called \emph{offsets}. Each face normal points outward.
The polyhedron is closed and convex. It is
\emph{full-dimensional} if it has an interior point, and \emph{proper} if it
is neither empty nor \(\R^n\). The description is \emph{irredundant} if
deleting any one inequality changes the set; deleting redundant or repeated
inequalities one at a time turns every description into an irredundant one
of the same set. For \(j\le m\), put \(H_j\coloneqq \{x:u_j\cdot x=a_j\}\) and
\(F_j\coloneqq Q\cap H_j\). A \emph{face} of \(Q\) is \(Q\) itself or a set
\(Q\cap\{u\cdot x=c\}\) with \(u\ne0\) and \(Q\subseteq\{u\cdot x\le c\}\).
A \emph{facet} is a face of dimension \(n-1\), and a \emph{ridge}, when \(n\ge2\), is a face
of dimension \(n-2\). In the plane the facets are the sides of a polygon
and the ridges its corners, in \(\R^3\) the facets are the faces of a solid
and the ridges its edges, and on the line the facets are the endpoints of an
interval, which has no ridges. We use these standard properties of faces:
the relative interiors of the faces of \(Q\) are pairwise disjoint and cover
\(Q\); a face of a face of \(Q\) is a face of \(Q\); and a face \(G'\)
of a face \(G\) with \(G'\ne G\) lies in the relative boundary
\(G\setminus\relint G\) of \(G\) (\TBc{Roc70}, Section~18; \TBc{Zie95}, Chapter~2).

The next fact identifies the pieces that contribute to the boundary
calculation. In particular, every ridge belongs to exactly two facets;
this allows us to combine their coefficients before taking absolute values.

\begin{fact}[Facets and ridges of a polyhedron]\label{tb:facets}
Let \(Q\) be a proper, full-dimensional polyhedron in \(\R^n\) with an
irredundant description. Statements about ridges are omitted when \(n=1\).
\begin{enumerate}
\item[(a)] The face normals \(u_1,\ldots,u_m\) are pairwise distinct, each
\(F_j\) is a facet, \(Q\) has no other facets, and
\(\partial Q=F_1\cup\cdots\cup F_m\).
\item[(b)] For \(j\ne k\), \(F_j\cap F_k\subseteq H_j\cap H_k\), and
\(H_j\cap H_k\) is empty if \(u_k=-u_j\), and an affine subspace of
dimension \(n-2\) otherwise. In particular \(F_j\cap F_k\) lies in a
finite union of affine subspaces of dimension at most \(n-2\).
\item[(c)] Each \(F_j\) is a full-dimensional polyhedron in the hyperplane
\(H_j\). Write the points of \(H_j\) as \(a_ju_j+y\) with \(y\perp u_j\), and,
for \(k\ne j\) with \(u_k\ne\pm u_j\), put
\(\alpha_{jk}\coloneqq \arccos(u_j\cdot u_k)\in(0,\pi)\),
\[
 \nu_{k\mid j}\coloneqq \frac{u_k-(u_j\cdot u_k)u_j}{\sin\alpha_{jk}},\qquad
 d_{k\mid j}\coloneqq \frac{a_k-(u_j\cdot u_k)a_j}{\sin\alpha_{jk}} .
\]
Then \(F_j\) is described in the coordinate \(y\) by the inequalities
\(\nu_{k\mid j}\cdot y\le d_{k\mid j}\), and its facets are exactly the
ridges \(E_{jk}\coloneqq F_j\cap F_k\) of \(Q\) of dimension \(n-2\); the outward
unit normal of \(F_j\) at \(E_{jk}\) is \(\nu_{k\mid j}\).
\item[(d)] Every ridge of \(Q\) lies in exactly two facets, so the ridges
correspond to the unordered pairs \(\{j,k\}\) with
\(\dim(F_j\cap F_k)=n-2\); for such a pair \(u_k\ne\pm u_j\), and the
only inequalities active at a point of \(\relint E_{jk}\) are \(j\) and \(k\).
\end{enumerate}
\end{fact}

\begin{proof}
(a) Fix \(j\). Irredundancy gives a point \(x_j\) with \(u_j\cdot x_j>a_j\)
and \(u_k\cdot x_j\le a_k\) for \(k\ne j\). Let \(x_0\) be an interior point of
\(Q\), and let \(x_*\) be the point where the segment from \(x_0\) to \(x_j\)
crosses \(H_j\). For \(k\ne j\), \(u_k\cdot x_*<a_k\), because \(x_*\) is a
convex combination of \(x_0\), where the inequality is strict, with positive
weight, and of \(x_j\), where it holds weakly. Hence a neighborhood of
\(x_*\) in \(H_j\) lies in \(F_j\), and \(\dim F_j=n-1\). Two inequalities
with the same normal cannot both be irredundant, since the one with the
larger offset is implied by the other. A point of \(Q\) at
which no inequality is active is interior, and at a point where
inequality \(j\) is active the points \(x+\varepsilon u_j\) lie outside \(Q\),
so \(\partial Q=\bigcup_jF_j\). If \(G=Q\cap\{u\cdot x=c\}\) is a facet and
\(x\in\relint G\), some inequality \(j\) is active at \(x\); on \(G\) the
affine function \(u_j\cdot x-a_j\le0\) attains its maximum at a relative
interior point and is therefore zero, so \(G\subseteq H_j\), and since \(G\) has
dimension \(n-1\), \(G=Q\cap H_j=F_j\).

(b) If \(u_k=-u_j\) and \(H_j\cap H_k\ne\emptyset\), then \(Q\) would lie in
the hyperplane \(H_j\), contradicting full dimension. If \(u_k\ne\pm u_j\),
the two equations are independent.

(c) Suppose \(n\ge2\). The set \(F_j\) is the polyhedron in \(H_j\) cut out by the remaining
inequalities, which take the stated form in the coordinate \(y\) because
\(u_k\cdot(a_ju_j+y)=a_ju_j\cdot u_k+(u_k-(u_j\cdot u_k)u_j)\cdot y\); an
inequality with \(u_k=-u_j\) holds automatically on \(H_j\) since
\(F_j\ne\emptyset\). After removing inequalities redundant within \(H_j\), part (a)
applied there shows that the facets of \(F_j\) are
the faces of \(F_j\) of dimension \(n-2\), hence faces of \(Q\) of that
dimension, that is, ridges of \(Q\) in \(F_j\); each equals
\(F_j\cap H_k=F_j\cap F_k\) for the inequality \(k\) that defines it.

(d) Let \(E\) be a ridge and \(y\in\relint E\), and let \(J\) be the set of
inequalities active at \(y\). On \(E\) each affine function
\(u_i\cdot x-a_i\), \(i\in J\), is at most zero and vanishes at the relative
interior point \(y\), so it vanishes on \(E\); thus \(E\subseteq F_i\) for
\(i\in J\) and \(\dim\operatorname{span}\{u_i:i\in J\}\le2\). A single active
inequality is impossible, since then \(y\) would lie in the relative interior
of \(F_j\) and the only face of \(F_j\) containing it is \(F_j\) itself, of
dimension \(n-1\). Opposite face normals \(u_i=-u_j\) are excluded by (b). Hence
the active face normals span a plane \(\Pi\). Near \(y\), \(Q\) coincides with
\(y\) plus the product of \(\Pi^\perp\) and a planar cone \(K\subseteq\Pi\)
with nonempty interior, defined by the active inequalities and contained in
a wedge of opening \(\pi-\alpha_{jk}<\pi\); this cone has exactly two
boundary rays. Each active inequality \(i\) has \(\dim F_i=n-1\), so
\(F_i\) meets every neighborhood of \(y\) in a set of dimension \(n-1\), and
the line \(\{u_i\cdot v=0\}\) must therefore carry a boundary ray of \(K\).
With distinct face normals this leaves exactly two active inequalities.
\end{proof}

\usedin{the facet--ridge decomposition leading to
\PE{appx-an:surface-identity}{C.30} and the pairing of the two terms
attached to each ridge in \PR{Lemma}{appx-an:lem:surface}{C.5}.}
\emph{Source.} These are standard facts about faces of polyhedra; see
\TBc{Roc70}, Sections~18 and~19, and \TBc{Zie95}, Chapter~2.
The planar-cone argument above explains why precisely two facets
meet at a ridge, including for unbounded polyhedra.

To attach a tube to a face, we must know which outward directions
keep that face as the nearest part of the polyhedron. These directions
form its normal cone, determined by the active inequalities.

\begin{fact}[Normal cones]\label{tb:normalcone}
Let \(Q\) be a polyhedron as in Fact~\ref{tb:facets}. For \(x\in Q\) let
\(J(x)\coloneqq \{j:u_j\cdot x=a_j\}\), and let
\[
 N_Q(x)\coloneqq \{v\in\R^n:v\cdot(q-x)\le0\text{ for all }q\in Q\}
\]
be the \emph{normal cone} of \(Q\) at \(x\). Then
\(N_Q(x)=\{\sum_{j\in J(x)}c_ju_j:c_j\ge0\}\). In particular
\(N_Q(x)=\{0\}\) at interior points, \(N_Q(x)=\{cu_j:c\ge0\}\) for
\(x\in\relint F_j\), and \(N_Q(x)=\{c_ju_j+c_ku_k:c_j,c_k\ge0\}\), a sector of
angle \(\alpha_{jk}\) in the plane spanned by \(u_j\) and \(u_k\), for
\(x\in\relint E_{jk}\).
\end{fact}

\begin{proof}
Put \(T\coloneqq \{d:u_j\cdot d\le0\text{ for }j\in J(x)\}\). A vector \(v\) lies in
\(N_Q(x)\) if and only if \(v\cdot d\le0\) for all \(d\in T\): if \(d\in T\)
then \(x+\varepsilon d\in Q\) for small \(\varepsilon>0\), because the
inactive inequalities are strict at \(x\), and conversely \(q-x\in T\) for
\(q\in Q\). By Farkas's lemma, which says that \(v\cdot d\le0\) for all
\(d\) with \(u_j\cdot d\le0\) (\(j\in J\)) if and only if \(v\) is a
nonnegative combination of the \(u_j\), \(j\in J\), this is the stated
cone. For \(x\in\relint F_j\), \(J(x)=\{j\}\), since an index \(k\ne j\) in
\(J(x)\) would put \(x\) in \(F_j\cap F_k\), a face of \(F_j\) of dimension
at most \(n-2\), and so in the relative boundary of \(F_j\). For
\(x\in\relint E_{jk}\), \(J(x)=\{j,k\}\) by Fact~\ref{tb:facets}(d).
\end{proof}

\usedin{the choice of the normal interval over a facet and the normal
sector over a ridge in \PR{Lemma}{appx-an:lem:surface}{C.5}.}
\emph{Source.} \TBc{Roc70}, Sections~22 and~23, for Farkas's lemma and normal
cones; \TBc{RW98}, Theorem~6.46.

The nearest-point characterization makes the outward tubes disjoint.
It applies to any nonempty closed convex set; for such a set we use
\(N_Q(x)=\{v:v\cdot(q-x)\le0\text{ for all }q\in Q\}\).

\begin{fact}[Nearest point]\label{tb:projection}
Let \(Q\subseteq\R^n\) be nonempty, closed, and convex. For each
\(y\in\R^n\) there is a unique point \(\prj_Q(y)\in Q\) with
\(\norm{y-\prj_Q(y)}=\dist(y,Q)\), and \(x=\prj_Q(y)\) if and only if
\(x\in Q\) and \(y-x\in N_Q(x)\). Consequently every \(y\in\R^n\) can be
written in exactly one way as \(y=x+v\) with \(x\in Q\) and \(v\in N_Q(x)\),
and then \(\norm v=\dist(y,Q)\). A direction \(v\) lies in \(N_Q(x)\) if and
only if \(x\) is the nearest point of \(x+s v\) in \(Q\) for every \(s>0\).
\end{fact}

\begin{proof}
For \(x,q\in Q\), \(\norm{y-q}^2=\norm{y-x}^2+\norm{x-q}^2-2(y-x)\cdot(q-x)\).
If \(y-x\in N_Q(x)\), the last term is nonnegative, so
\(\norm{y-q}^2\ge\norm{y-x}^2+\norm{x-q}^2>\norm{y-x}^2\) for \(q\ne x\); thus
\(x\) is the unique nearest point. Conversely, if \((y-x)\cdot(q-x)>0\) for
some \(q\in Q\), then moving from \(x\) toward \(q\) decreases the distance
to \(y\) for small steps, by the displayed identity applied to
\(x+\varepsilon(q-x)\in Q\), so \(x\) is not nearest. For existence, choose \(q_0\in Q\). A minimizing sequence may be
restricted to the closed ball centered at \(y\) with radius
\(\norm{y-q_0}+1\); it therefore has a convergent subsequence in
finite dimension. Closedness puts its limit in \(Q\), and continuity
of the norm makes that limit a minimizer. The last statement follows because \(N_Q(x)\) is a
cone.
\end{proof}

\usedin{disjointness of the variable-width tubes in
\PR{Lemma}{appx-an:lem:surface}{C.5}; each tube point has its nearest
point in the relative interior of exactly one face.}
\emph{Source.} \TBc{RW98}, Example~2.25, for existence and uniqueness,
and Proposition~6.17 for the characterization by the normal cone.

\subsection{Surface measure, Gaussian surface integrals, and the
divergence theorem}\label{tb:sec:surface}

We integrate each facet with its ordinary Euclidean area measure.
Hausdorff measure gives a common notation for facets and ridges in
all dimensions and shows why their lower-dimensional overlaps do not
affect a surface integral.

\begin{fact}[Hausdorff measure on flat pieces]\label{tb:hausdorff}
For \(0\le k\le n\), the \(k\)-dimensional Hausdorff measure \(\Haus^k\) on
\(\R^n\) is obtained from the covering formula
\[
 \Haus^k(S)\coloneqq \lim_{\epsilon\downarrow0}\inf\Bigl\{\sum_i\omega_k
 \Bigl(\tfrac12\operatorname{diam}C_i\Bigr)^k:S\subseteq\bigcup_iC_i,\
 \operatorname{diam}C_i\le\epsilon\Bigr\},
\]
where the infimum runs over countable covers and \(\omega_k\) is
the volume of the unit ball of \(\R^k\).
For \(k=0\), every nonempty covering set contributes one, including
a singleton. The formula defines an outer measure, whose restriction
to Borel sets is the measure used here. (a) On an
affine subspace of dimension \(k\), \(\Haus^k\) is \(k\)-dimensional Lebesgue
measure; in particular \(\Haus^0\) is counting measure and \(\Haus^n\) is
Lebesgue measure on \(\R^n\). (b) \(\Haus^k\) is invariant under isometries.
(c) A set contained in a finite union of affine subspaces of dimension less
than \(k\) has \(\Haus^k\)-measure zero.
\end{fact}

Part (c) follows from (a), since each such subspace lies in a
\(k\)-dimensional affine subspace in which it is a Lebesgue null set. For a
proper, full-dimensional polyhedron \(Q\) with an irredundant description,
and for a function \(f\ge0\) or \(f\) integrable on \(\partial Q\), the
\emph{surface integral} is
\[
 \int_{\partial Q}f\,d\Haus^{n-1}=\sum_{j=1}^m\int_{F_j}f\,d\Haus^{n-1},
\]
where on each facet \(\Haus^{n-1}\) is Lebesgue measure on the hyperplane
\(H_j\) by part~(a). By Fact~\ref{tb:facets}(b) and part~(c), distinct
facets overlap in sets of \(\Haus^{n-1}\)-measure zero, so the sum is the
integral of \(f\) over the set \(\partial Q\) with respect to
\(\Haus^{n-1}\). It depends on the set \(Q\) and not on its description,
which justifies deleting redundant inequalities before computing it.
The Gaussian surface integral \(\surf(Q)=\int_{\partial Q}\Hh\,\varphi_r\,
d\Haus^{r-1}\) of Subsection~\PN{sec:outline-late}{2.4} is the case
\(n=r\) and \(f=\Hh\varphi_r\), with \(\Hh(g)\coloneqq \norm g^2-r\).
The integral is signed; we take its absolute value only after summing
the facet integrals. We set \(\surf(\varnothing)=\surf(\R^r)\coloneqq0\).
Lower-dimensional polyhedra are excluded from this definition and are
handled as exceptional levels in Fact~\ref{tb:variation}.

\usedin{the definition of \(\surf\) in
Subsection~\PN{sec:outline-late}{2.4} and
\PR{Lemma}{appx-an:lem:surface}{C.5}.}
\emph{Source.} \TBc{EG15}, Chapter~2.

Integrating a tangential divergence over a facet converts its interior
integral into contributions from its ridges. We prove the needed divergence
formula by slicing, then use a cutoff to allow an unbounded polyhedron
under precisely the stated decay assumptions.

\begin{samepage}
\begin{fact}[Gauss--Green formula (divergence theorem) for a polyhedron]\label{tb:divergence}
Let \(Q\subseteq\R^n\) (\(n\ge1\)) be a proper, full-dimensional polyhedron
with an irredundant description, facets \(F_1,\ldots,F_m\), and unit
face normals \(u_1,\ldots,u_m\). Let \(X:\R^n\to\R^n\) be of class \(C^1\), with
\(\norm{X(x)}+|\operatorname{div}X(x)|\le C(1+\norm x)^k\varphi_n(x)\) for
some \(C\) and \(k\). Then
\[
 \int_Q\operatorname{div}X\,dx=\sum_{j=1}^m\int_{F_j}X\cdot u_j\,d\Haus^{n-1}.
\]
For \(n=1\) the facets are the endpoints of an interval, \(\Haus^0\) is
counting measure, and the identity is the fundamental theorem of calculus.
\end{fact}
\end{samepage}

\begin{proof}
We first prove the formula for a compactly supported \(C^1\) vector
field. For a unit vector \(e\) and a compactly supported
\(C^1\) scalar function \(f\), we claim that
\begin{equation}
 \int_Q\partial_ef\,dx=\sum_j(e\cdot u_j)\int_{F_j}f\,d\Haus^{n-1},
 \tag{\(*\)}
\end{equation}
since the theorem follows by taking \(e=e_i\) and \(f=X_i\) and summing over
\(i\). Write \(x=w+s e\) with \(w\in e^\perp\). By Fubini's theorem,
\(\int_Q\partial_ef\,dx=\int_{e^\perp}\int_{I(w)}\partial_sf(w+s e)\,ds\,dw\),
where \(I(w)\coloneqq \{s:w+s e\in Q\}\) is a closed interval, possibly empty, a point,
or unbounded, because \(Q\) is convex. Since \(f\) tends to zero along every
line, the inner integral equals \(f(w+s_+(w)e)-f(w+s_-(w)e)\), where
\(s_\pm(w)\) are the endpoints of a nonempty \(I(w)\).
An empty interval contributes zero, and a term with an infinite
endpoint is read as zero.

Consider the upper endpoints \(x=w+s_+(w)e\). If \(x\in\relint F_j\), then
only the inequality \(j\) is active at \(x\) (Fact~\ref{tb:normalcone}), so
\(x+s e\in Q\) for small \(s>0\) unless \(e\cdot u_j>0\), and \(x\) is an upper
endpoint if and only if \(e\cdot u_j>0\). The upper endpoints that lie in no
\(\relint F_j\) lie in \(\bigcup_{j\ne k}F_j\cap F_k\), which is a subset of a
finite union of affine subspaces of dimension at most \(n-2\)
(Fact~\ref{tb:facets}(b)); its projection along \(e\) onto \(e^\perp\) is a
Lebesgue null set of \(e^\perp\). For fixed \(j\) with \(e\cdot u_j>0\), the
map \(x\mapsto x-(x\cdot e)e\) is injective on the set of upper endpoints in
\(\relint F_j\), since a line meets \(Q\) in one interval, and it is the
restriction to \(H_j\) of a linear map whose Jacobian with respect to
\(\Haus^{n-1}\) on \(H_j\) is \(e\cdot u_j\): if \(b_1,\ldots,b_{n-1}\) is an
orthonormal basis of \(u_j^\perp\), the volume of the projection along \(e\)
of the cube they span equals the volume of the parallelepiped spanned by
\(b_1,\ldots,b_{n-1},e\), which is \(|e\cdot u_j|\). Therefore
\(\int f(w+s_+(w)e)\,dw=\sum_{j:\,e\cdot u_j>0}(e\cdot u_j)\int_{F_j}f\,d\Haus^{n-1}\),
the sets \(F_j\setminus\relint F_j\) being \(\Haus^{n-1}\)-null. The same
argument for the lower endpoints gives
\(\int f(w+s_-(w)e)\,dw
=\sum_{j:\,e\cdot u_j<0}|e\cdot u_j|\int_{F_j}f\,d\Haus^{n-1}\).
Subtracting this lower-endpoint integral from the upper-endpoint
integral proves \((*)\). Facets with \(e\cdot u_j=0\) contribute
nothing. Summing over coordinate directions proves the compactly
supported vector-field formula.

For the stated field \(X\), choose smooth cutoffs \(\chi_R\) equal to
one on the ball of radius \(R\), zero outside the ball of radius
\(2R\), with \(0\le\chi_R\le1\) and
\(\norm{\nabla\chi_R}\le C/R\). Apply the formula to \(\chi_RX\).
The product rule gives
\[
 \operatorname{div}(\chi_RX)
 =\chi_R\operatorname{div}X+\nabla\chi_R\cdot X.
\]
The first term converges in integral by the assumed integrable bound;
the second has integral in absolute value at most
\((C/R)\int_{\{\norm x\ge R\}}\norm{X(x)}\,dx\to0\).
On each of the finitely many facets, polynomial times Gaussian decay
is integrable, so dominated convergence also applies to the boundary
integrals. This proves the formula without requiring separate bounds
on the component derivatives of \(X\).
\end{proof}

\usedin{the divergence calculation inside a facet in
\PR{Lemma}{appx-an:lem:surface}{C.5}. The vector field is
\(X(y)\coloneqq y\varphi_{r-1}(y)\), with divergence
\(\{r-1-\norm y^2\}\varphi_{r-1}(y)\); the boundary terms are the
ridge contributions in \PE{appx-an:surface-identity}{C.30}.}
\emph{Source.} This is the classical Gauss--Green formula specialized
to a polyhedron; see \TBc{EG15}, Chapter~5. The cutoff argument above
extends the compact-support formula to the stated Gaussian-decay
hypothesis on an unbounded polyhedron.

For the surface integral used in the paper, this divergence
calculation produces coefficients whose separate expressions contain
\(1/\sin\alpha_{jk}\). We combine the two contributions of a ridge
before estimating them. Let
\(s_j\coloneqq \int_{F_j}\varphi_r\,d\Haus^{r-1}\),
\(\rho_{jk}\coloneqq \int_{E_{jk}}\varphi_r\,d\Haus^{r-2}\), and let
\(z_{jk}\) be the projection of a ridge point onto
\(\operatorname{span}(u_j,u_k)\). Applying the divergence formula in
each facet gives
\[
 \surf(Q)=\sum_j(a_j^2-1)s_j
       -\sum_{\{j,k\}}(d_{k\mid j}+d_{j\mid k})\rho_{jk}.
\]
The offsets \(d_{k\mid j}\) from Fact~\ref{tb:facets}(c) satisfy
\[
 d_{k\mid j}+d_{j\mid k}
 =2\sin(\alpha_{jk}/2)\,
       \frac{u_j+u_k}{\norm{u_j+u_k}}\cdot z_{jk},\qquad
 |d_{k\mid j}+d_{j\mid k}|\le\alpha_{jk}\norm{z_{jk}}.
\]
Indeed, their sum is
\((1-\cos\alpha_{jk})(a_j+a_k)/\sin\alpha_{jk}\), and
\(a_j+a_k=(u_j+u_k)\cdot z_{jk}\).
The inverse sine cancels. The remaining factor \(\alpha_{jk}\)
is precisely the angle of the normal sector whose mass we use next.

\subsection{Tubes around a polyhedron}\label{tb:sec:tubes}

We now compare the paired boundary terms with Gaussian mass outside
the polyhedron. The Gaussian density changes rapidly when a boundary point has a large
normal projection. We therefore shorten the outward tube at that point.
Its width may vary along a face: the required integration formula still
follows from orthogonal coordinates and Fubini's theorem.

Let \(Q\) be a proper, full-dimensional polyhedron in \(\R^n\) with
facets \(F_j\), ridges \(E_{jk}\), and an irredundant description by
unit face normals \(u_j\). For a ridge, put
\(L_{jk}\coloneqq \operatorname{span}(u_j,u_k)\) and
\(z_{jk}\coloneqq \Pi_{jk}x\), for any \(x\in E_{jk}\). This vector is constant
on the ridge because
\(u_j\cdot z_{jk}=a_j\) and \(u_k\cdot z_{jk}=a_k\).
Let \(\eta(x)>0\) be a measurable function on the facets and ridges.
We use the same value of \(\eta\) at a point whenever it belongs to
more than one closed face; only relative interiors enter the tubes.
The outward tubes are
\begin{align*}
 \mathcal T_j&\coloneqq \{x+s u_j:x\in\relint F_j,\ 0<s<\eta(x)\},\\
 \mathcal T_{jk}&\coloneqq \{x+v:x\in\relint E_{jk},\ v\in N_Q(x),\
                         0<\norm v<\eta(x)\}.
\end{align*}
For \(x\in\relint E_{jk}\), the normal cone is a sector of angle
\(\alpha_{jk}\) in \(L_{jk}\). Write
\(S_{jk}(a)\coloneqq N_Q(x)\cap\{v:0<\norm v<a\}\); its area is
\(\alpha_{jk}a^2/2\). For \(Q=\bigcap_j\{u_j\cdot x\le a_j\}\),
the notation \(Q^s\) increases every threshold by \(s\); it does not
denote a Euclidean neighborhood.

\begin{fact}[Variable-width normal tubes]\label{tb:tubes}
(a) The tubes are measurable, and tubes over distinct facets and
ridges are pairwise disjoint.
Each tube point \(g=x+v\) has nearest point \(x\) in \(Q\), and
\(\dist(g,Q)=\norm v\).

(b) For a nonnegative measurable function \(f\),
\begin{align*}
 \int_{\mathcal T_j}f(g)\,dg
 &=\int_{F_j}\int_0^{\eta(x)} f(x+s u_j)\,ds\,d\Haus^{n-1}(x),\\
 \int_{\mathcal T_{jk}}f(g)\,dg
 &=\int_{E_{jk}}\int_{S_{jk}(\eta(x))}f(x+v)\,dv\,d\Haus^{n-2}(x).
\end{align*}
The same formulas hold whenever \(f\) is integrable on the relevant tube.

(c) Use the function \(R\) in Fact~\ref{tb:projection-moments}, built
from the face normals of \(Q\), and choose \(\eta(x)\coloneqq 1/(4R(x))\).
Then all these tubes lie in \(\{g:g\in Q^{1/R(g)}\setminus Q\}\).
For an absolute constant \(c>0\), their normal integrals obey
\begin{align*}
 \int_0^{\eta(x)}R(x+s u_j)^3\varphi_n(x+s u_j)\,ds
 &\ge c\{1+a_j^2\}\varphi_n(x),\qquad x\in F_j,\\
 \int_{S_{jk}(\eta(x))}R(x+v)^3\varphi_n(x+v)\,dv
 &\ge c\alpha_{jk}\norm{z_{jk}}\varphi_n(x),\qquad x\in E_{jk}.
\end{align*}
For \(n=1\), there are no ridge tubes.
\end{fact}

\begin{proof}
(a) Facts~\ref{tb:normalcone} and~\ref{tb:projection} identify \(x\)
as the unique nearest point of \(x+v\). The relative interiors of
distinct faces are disjoint, so the tubes are disjoint.

(b) The map \((x,s)\mapsto x+s u_j\) is an affine isometry from
\(H_j\times\R\) onto \(\R^n\). Likewise, \((x,v)\mapsto x+v\)
is an affine isometry from \((z_{jk}+L_{jk}^\perp)\times L_{jk}\)
onto \(\R^n\). The domains are measurable in these coordinates because \(\eta\)
is measurable; their images under the affine isometries are therefore
measurable as well. The Jacobians are one. Apply Fubini to the indicator
of the domain \(0<s<\eta(x)\), or \(0<\norm v<\eta(x)\) in the
normal sector. No derivative of \(\eta\) is involved, and no
smoothness of \(\eta\) is needed.

(c) For a facet displacement \(v=s u_j\), and for a ridge displacement
\(v\in L_{jk}\), respectively,
\[
 |x\cdot v|\le\mathcal M(x)\norm v\le\tfrac14.
\]
The ridge case uses \(x\cdot v=z_{jk}\cdot v\). Since
\(\norm v\le1/4\), Fact~\ref{tb:gauss-coord}(b) gives
\(\varphi_n(x+v)\ge e^{-1/4-1/32}\varphi_n(x)\).
The Lipschitz property and \(R\ge1\) also give
\[
 \tfrac34R(x)\le R(x+v)\le\tfrac54R(x).
\]
The facet integral is therefore at least
\(cR(x)^3\eta(x)\varphi_n(x)=c'R(x)^2\varphi_n(x)\), which
dominates \(c''(1+a_j^2)\varphi_n(x)\) because
\(R(x)\ge1+|u_j\cdot x|=1+|a_j|\). The ridge integral is at least
\(c\alpha_{jk}R(x)^3\eta(x)^2\varphi_n(x)
=c'\alpha_{jk}R(x)\varphi_n(x)\); use
\(R(x)\ge\norm{z_{jk}}\).

Finally, a tube point \(g=x+v\) lies outside \(Q\), and each
unit-normal inequality defining \(Q\) satisfies
\[
 u_j\cdot g-a_j\le\norm v<\frac1{4R(x)}\le\frac1{R(g)}.
\]
This proves the inclusion in (c), even if the threshold system used
to define \(Q^s\) also contains redundant inequalities.
\end{proof}

Together with the paired ridge coefficients derived above, part (c)
gives
\[
 |\surf(Q)|\le C\E\left[R(G)^3
              \ind\{G\in Q^{1/R(G)}\setminus Q\}\right].
\]
The power three serves both dimensions: integration along a facet
uses one power of the tube width, leaving \(R^2\), while integration
over a ridge sector uses two, leaving \(R\).

\usedin{\PE{appx-an:adaptive-tubes}{C.29} in
\PR{Lemma}{appx-an:lem:surface}{C.5}. After setting
\(Q'\coloneqq (Q-V_0)/\sigma\), the event is
\(V_0+\sigma G\in Q^{\sigma/R(G)}\setminus Q\).
Conditional on \(G\), its width is fixed; concentration of \(V_0\)
and Fact~\ref{tb:projection-moments} give the averaged surface bound.}
\emph{Source.} The facet construction is Nazarov's outward-normal
argument~\TBc{Naz03}; see the detailed proof in \TBc{CCK17N},
Lemmas~1 and~2. Disjointness over arbitrary faces follows from the
classical projection characterization in \TBc{RW98}, Proposition~6.17.
The variable widths, ridge-sector integrals, and weighted estimates
needed here are derived above. For the related use of facet and ridge
tubes in Gaussian approximation, see \TBc{10}.

\subsection{The coarea formula and the first variation of a Gaussian
integral}\label{tb:sec:coarea}

The logistic mixture moves all thresholds together. To differentiate
its Gaussian integral, we integrate first over the level sets of the
maximum constraint. The coarea formula provides this change of variables.

\begin{fact}[Coarea formula]\label{tb:coarea}
Let \(\psi:\R^n\to\R\) be Lipschitz, so that \(\nabla\psi\) exists almost
everywhere. For every measurable \(f\ge0\),
\[
 \int_{\R^n}f\,\norm{\nabla\psi}\,dx=\int_{-\infty}^\infty
 \Bigl(\int_{\{\psi=s\}}f\,d\Haus^{n-1}\Bigr)ds ,
\]
For signed \(f\), the same identity holds if
\(\int |f|\norm{\nabla\psi}\,dx<\infty\), in particular if
\(f\in L^1(\R^n)\), since \(\psi\) is Lipschitz.
\end{fact}

\emph{Source.} \TBc{EG15}, Section~3.4. The proofs use the formula only for
a maximum of finitely many affine functions with gradients of equal length, and
for these Fact~\ref{tb:variation} gives a direct proof from Fubini's theorem.

For the finite maximum used here, the coarea formula can be proved
directly by partitioning the space according to the maximizing face.
This proof also identifies the exceptional levels and explains why the
independent Gamma shift permits differentiation at every fixed threshold.

\begin{fact}[First variation of a polyhedral sublevel integral]\label{tb:variation}
Let \(u_1,\ldots,u_m\) be unit vectors in \(\R^n\) (\(m\ge1\)), \(\kappa>0\),
\(b_1,\ldots,b_m\in\R\), and
\[
 \psi(g)\coloneqq \max_{j\le m}\{\kappa\,u_j\cdot g-b_j\},\qquad Q_s\coloneqq \{g:\psi(g)\le s\}.
\]
(a) For \(f\ge0\) measurable, or \(f\) integrable, and \(s\in\R\),
\[
 \int_{Q_s}f\,dg=\kappa^{-1}\int_{-\infty}^s\Bigl(\int_{\{\psi=s'\}}f\,d\Haus^{n-1}
 \Bigr)ds'.
\]
(b) If \(f\) is integrable, then \(\Lambda(s)\coloneqq \int_{Q_s}f\,dg\) is absolutely
continuous, and for almost every \(s\)
\[
 \Lambda'(s)=\kappa^{-1}\int_{\{\psi=s\}}f\,d\Haus^{n-1},\qquad
 |\Lambda'(s)|\le\kappa^{-1}\sum_{j=1}^m\int_{\{\kappa u_j\cdot g-b_j=s\}}|f|\,d\Haus^{n-1}.
\]
(c) If \(s>\inf\psi\), then \(Q_s\) is a proper, full-dimensional polyhedron
with face normals \(u_j\) and offsets \((s+b_j)/\kappa\), \(\{\psi=s\}=\partial Q_s\),
and \(\int_{\{\psi=s\}}f\,d\Haus^{n-1}\) is the surface integral of \(f\) over
\(\partial Q_s\) of Subsection~\ref{tb:sec:surface}. For \(s<\inf\psi\) the set
\(Q_s\) is empty, and \(s=\inf\psi\) is the only other level at which \(Q_s\)
may fail to be full-dimensional.

(d) For \(\Lambda\) from (b), assume in addition that it is
Lipschitz. If \(S\) is a real random variable with a density, then \(s\mapsto\E\Lambda(s+S)\) is differentiable with derivative
\(\E\Lambda'(s+S)\), where \(\Lambda'\) is defined almost everywhere.
\end{fact}

\begin{proof}
(a) Let \(C_j\) be the set where \(j\) is the smallest index attaining the
maximum defining \(\psi\); the Borel sets \(C_1,\ldots,C_m\) partition
\(\R^n\). On \(C_j\), \(\psi(g)=\kappa v-b_j\) in the coordinates
\(g=v u_j+y\), \(y\perp u_j\). By Fact~\ref{tb:fubini}, in these coordinates,
\[
 \int_{C_j\cap Q_s}f\,dg=\int_{u_j^\perp}\int_{-\infty}^{(s+b_j)/\kappa}
 (f\ind_{C_j})(v u_j+y)\,dv\,dy .
\]
The substitution \(s'=\kappa v-b_j\) turns the inner integral into
\(\kappa^{-1}\int_{-\infty}^s(f\ind_{C_j})(v u_j+y)\,ds'\) with
\(v=(s'+b_j)/\kappa\), and a further application of Fubini's theorem
exchanges the integrals. For fixed \(s'\), the set of \(v u_j+y\) in \(C_j\)
with this \(v\) is the slice \(C_j\cap\{\psi=s'\}\), and on the hyperplane
containing it \(\Haus^{n-1}\) is Lebesgue measure \(dy\)
(Fact~\ref{tb:hausdorff}(a)). Hence
\(\int_{C_j\cap Q_s}f\,dg=\kappa^{-1}\int_{-\infty}^s\int_{C_j\cap\{\psi=s'\}}
f\,d\Haus^{n-1}\,ds'\). Summing over \(j\) gives (a), since the slices
\(C_j\cap\{\psi=s'\}\) partition \(\{\psi=s'\}\).

(b) Applying (a) to \(|f|\) and letting \(s\to\infty\) shows that
\(I(s')\coloneqq \kappa^{-1}\int_{\{\psi=s'\}}f\,d\Haus^{n-1}\) is integrable on
\(\R\), and (a) says \(\Lambda(s)=\int_{-\infty}^sI(s')\,ds'\). Hence \(\Lambda\) is
absolutely continuous, and \(\Lambda'=I\) almost everywhere by Lebesgue's
differentiation theorem. The bound follows because
\(\{\psi=s\}\subseteq\bigcup_j\{\kappa u_j\cdot g-b_j=s\}\), a union of \(m\)
hyperplanes.

(c) The set \(Q_s\) is the polyhedron \(\bigcap_j\{u_j\cdot g\le(s+b_j)/\kappa\}\),
and \(\psi\) is convex and continuous. If \(\psi(g)=s>\inf\psi\), choose \(g'\)
with \(\psi(g')<s\); for \(\tau>0\), \(g=(g_\tau+\tau g')/(1+\tau)\) with
\(g_\tau\coloneqq g+\tau(g-g')\), so convexity gives
\(\psi(g_\tau)\ge s+\tau(s-\psi(g'))>s\). Thus \(g\) is a limit of points
outside \(Q_s\) and \(g\in\partial Q_s\); the open set \(\{\psi<s\}\) is
nonempty and lies in \(Q_s\), which is therefore full-dimensional, and it
is a proper subset because \(\psi(g)\ge\kappa u_1\cdot g-b_1\to\infty\) along
\(u_1\). A point with \(\psi<s\) is interior, so \(\partial Q_s=\{\psi=s\}\). After
the redundant and repeated inequalities are deleted, \(\partial Q_s=\{\psi=s\}\)
is the union of the sets \(Q_s\cap\{\kappa u_j\cdot g-b_j=s\}\) of the
remaining inequalities, which are faces of \(Q_s\) or empty. A face of
dimension \(n-1\) is a facet (Fact~\ref{tb:facets}(a)) and any other face is
\(\Haus^{n-1}\)-null, so the integral over \(\{\psi=s\}\) is the sum over the
facets of the irredundant description.

(d) The function \(\Lambda\) is bounded by \(\int|f|\), so the expectation
exists. A Lipschitz function is absolutely continuous and differentiable outside a
Lebesgue null set \(\mathcal N\). The random variable \(s+S\) avoids
\(\mathcal N\) almost surely, because \(S\) has a density, so the difference
quotients of \(\Lambda(s+S)\) converge almost surely and are bounded by the
Lipschitz constant. Dominated convergence gives the claim.
\end{proof}

\usedin{the first-variation paragraph in the proof of
\PR{Lemma}{appx-an:lem:matrix}{C.4}, immediately after
\PE{appx-an:gamma-mixture}{C.25}. Here
\(\psi(g)=\max_j\{u_j\cdot(V_0+\sigma g)-a_j\}\),
\(\kappa=\sigma\), \(b_j=a_j-u_j\cdot V_0\), and
\(f(g)=(\norm g^2-r)\varphi_r(g)\). Parts (a)--(c) give
\(\sigma^{-1}\surf((Q(a+s\mathbf1)-V_0)/\sigma)\) for almost every
\(s\). The independent shift \(-h\log T_\ell\), with
\(T_\ell\sim\operatorname{Gamma}(\ell,1)\), has a density, so part (d)
permits averaging these derivatives. A uniform bound
\(Cm(r+1)/\sigma\) follows from (b) and the Gaussian integral over
each supporting hyperplane; it justifies the interchange and does
not enter the final dimension-free bound.}

\clearpage
\subsection{Where each fact is used}\label{tb:sec:index}

The table follows the order of the proof. Each row identifies the
operation for which the corresponding facts are needed.

\medskip
\begin{center}
\small
\renewcommand{\arraystretch}{1.18}
\begin{tabular}{@{}p{0.27\textwidth}p{0.68\textwidth}@{}}
Facts~\ref{tb:fubini}, \ref{tb:taylor} &
Conditioning, Taylor remainders, and exchange of expectations with
convergent series in Appendices B--C.\\
Facts~\ref{tb:bernoulli}, \ref{tb:heat} &
Mask differentiation, the heat equation, and the signed Hessian
comparison in Subsection~\PN{appx-an:sec:path}{C.2}.\\
Facts~\ref{tb:faa}, \ref{tb:logistic}, \ref{tb:covariance} &
Derivative envelopes in Appendix~\PN{appx-sec:smoothing}{B}; early
bounds in \PN{appx-an:sec:early}{C.3} and coefficient matrices in
\PR{Lemma}{appx-an:lem:restoration}{C.2}.\\
Facts~\ref{tb:stein}, \ref{tb:nuclear} &
Two Gaussian integrations by parts and matrix contraction in
\PR{Lemma}{appx-an:lem:restoration}{C.2}.\\
Facts~\ref{tb:gauss-coord}, \ref{tb:logsum} &
Reference-face mean bounds and summation of reference masses in
\PR{Lemmas}{appx-an:lem:mean}{C.3} and~\PN{appx-an:lem:matrix}{C.4}.\\
Facts~\ref{tb:prekopa}--\ref{tb:lc-weights} &
Logistic-tilt covariance contraction in
\PR{Lemma}{appx-an:lem:matrix}{C.4}.\\
Facts~\ref{tb:facets}--\ref{tb:divergence} &
Surface measure and the facet--ridge identity in
\PR{Lemma}{appx-an:lem:surface}{C.5}.\\
Facts~\ref{tb:tubes}, \ref{tb:projection-moments} &
Adaptive normal tubes and Gaussian projection moments in
\PR{Lemma}{appx-an:lem:surface}{C.5}.\\
Facts~\ref{tb:coarea}, \ref{tb:variation} &
First variation of the Gamma mixture in
\PR{Lemma}{appx-an:lem:matrix}{C.4}.\\
\end{tabular}
\end{center}

\subsection{Sources}\label{tb:sec:sources}

\begin{description}\footnotesize\setlength{\itemsep}{0pt}\setlength{\parskip}{0pt}
\item[{[BL76]}] H.~J. Brascamp and E.~H. Lieb (1976). On extensions of the
Brunn--Minkowski and Pr\'ekopa--Leindler theorems, including inequalities
for log concave functions, and with an application to the diffusion
equation. \emph{Journal of Functional Analysis} \textbf{22}, 366--389.
\item[{[BV04]}] S. Boyd and L. Vandenberghe (2004).
\emph{Convex Optimization}. Cambridge University Press, Cambridge.
\item[{[CCK17N]}] V. Chernozhukov, D. Chetverikov, and K. Kato (2017).
Detailed proof of Nazarov's inequality. arXiv:1711.10696.
\item[{[Com74]}] L. Comtet (1974). \emph{Advanced Combinatorics}. Reidel,
Dordrecht.
\item[{[CT06]}] T.~M. Cover and J.~A. Thomas (2006). \emph{Elements of
Information Theory}, 2nd edition. Wiley.
\item[{[EG15]}] L.~C. Evans and R.~F. Gariepy (2015). \emph{Measure Theory
and Fine Properties of Functions}, revised edition. CRC Press.
\item[{[Fol99]}] G.~B. Folland (1999). \emph{Real Analysis: Modern
Techniques and Their Applications}, 2nd edition. Wiley.
\item[{[Gar02]}] R.~J. Gardner (2002). The Brunn--Minkowski inequality.
\emph{Bulletin of the American Mathematical Society} \textbf{39}, 355--405.
\item[{[10]}] Z. Gao and W. Zhang (2026). Cubic-root Gaussian approximation
under unrestricted covariance. arXiv:2608.30221v1.
\item[{[Lei72]}] L. Leindler (1972). On a certain converse of H\"older's
inequality II. \emph{Acta Scientiarum Mathematicarum (Szeged)} \textbf{33},
217--223.
\item[{[Naz03]}] F. Nazarov (2003). On the maximal perimeter of a convex
set in \(\R^n\) with respect to a Gaussian measure. In V.~D. Milman
and G. Schechtman (eds.), \emph{Geometric Aspects of Functional Analysis},
Lecture Notes in Mathematics \textbf{1807}, 169--187. Springer.
\item[{[Pre71]}] A. Pr\'ekopa (1971). Logarithmic concave measures with
application to stochastic programming. \emph{Acta Scientiarum Mathematicarum
(Szeged)} \textbf{32}, 301--316.
\item[{[Pre73]}] A. Pr\'ekopa (1973). On logarithmic concave measures and
functions. \emph{Acta Scientiarum Mathematicarum (Szeged)} \textbf{34},
335--343.
\item[{[Roc70]}] R.~T. Rockafellar (1970). \emph{Convex Analysis}. Princeton
University Press.
\item[{[RW98]}] R.~T. Rockafellar and R.~J.-B. Wets (1998). \emph{Variational
Analysis}. Springer, Berlin.
\item[{[Ste81]}] C.~M. Stein (1981). Estimation of the mean of a multivariate
normal distribution. \emph{The Annals of Statistics} \textbf{9}, 1135--1151.
\item[{[Zie95]}] G.~M. Ziegler (1995). \emph{Lectures on Polytopes}.
Springer, New York.
\end{description}



% Fallback paper numbers below are checked against the v3 auxiliary file.

\makeatletter
\providecommand{\E}{\mathbb E}
\providecommand{\Pp}{\mathbb P}
\providecommand{\R}{\mathbb R}
\providecommand{\tr}{\operatorname{tr}}
\providecommand{\op}{\mathrm{op}}
\providecommand{\norm}[1]{\left\lVert#1\right\rVert}
\providecommand{\ind}{\mathbf 1}
\providecommand{\dd}{\,d}
\providecommand{\Hh}{H_2}
\providecommand{\surf}{\mathcal J}
\providecommand{\relint}{\operatorname{relint}}
\providecommand{\aff}{\operatorname{aff}}
\providecommand{\cone}{\operatorname{cone}}
\providecommand{\dist}{\operatorname{dist}}
\providecommand{\prj}{\operatorname{pr}}
\providecommand{\Haus}{\mathcal H}
\providecommand{\TBc}[1]{[#1]}
% \PN{label}{C.5} prints a number of the paper (a \ref in the paper);
% \PR{Lemma}{label}{C.5} prints "Lemma C.5"; \PE{label}{C.21} an equation.
\providecommand{\PN}[2]{\ifdefined\toolboxinpaper\ref{#1}\else#2\fi}
\providecommand{\PR}[3]{#1~\PN{#2}{#3}}
\providecommand{\PE}[2]{(\PN{#1}{#2})}
\providecommand{\usedin}[1]{\par\smallskip\noindent{\small\emph{Used in:} #1}\par\smallskip}
\@ifundefined{c@fact}{\newtheorem{fact}{Fact}[section]}{}
\makeatother

\section{Toolbox: facts from analysis and geometry}\label{app:toolbox}

The late interpolation estimate reduces high derivatives to two objects:
a Gaussian matrix weighted by the logistic approximation, and a signed
Gaussian integral over a polyhedral boundary. This appendix supplies the
analysis and geometry behind those reductions. The calculus facts also
justify the early estimate and the smoothing argument in
Appendices~\PN{appx-sec:smoothing}{B} and~\PN{appx-sec:support}{C}.
Each fact is stated with the assumptions needed for its use, followed by
a proof or a standard reference and a precise pointer to the manuscript.

We first treat differentiation, Gaussian integration by parts, and
covariance contraction. Log-concavity then controls the Gaussian law
tilted by a logistic weight. The geometric facts explain how to pair the
two contributions at each ridge and compare the resulting boundary
integral with the mass of short outward tubes. Because the faces are
flat, the required area and tube formulas follow from orthogonal
coordinates and Fubini's theorem. Finally, the first-variation formula
connects the logistic mixture to the surface estimate, including the
exceptional levels at which the polyhedron loses dimension.
Subsection~\ref{tb:sec:index} records these dependencies in one table.

\subsection{Notation}\label{tb:sec:notation}

The smoothing and interpolation arguments take place in the covariance
support \(\R^r\). They use a fixed list of \(m=2p\) unit face normals
\(u_1,\ldots,u_m\), which may repeat, and the threshold sets
\[
 A_a\coloneqq\{x\in\R^r:u_\alpha\cdot x\le a_\alpha,
                          \ \alpha=1,\ldots,m\},\qquad a\in\R^m.
\]
All thresholds are finite; inconsistent inequalities give the empty set.
We keep the full list when defining the logistic approximation, even if
some inequalities are redundant. We delete redundant inequalities only
when describing the facets of a fixed full-dimensional polyhedron.
For a specified system of unit-normal inequalities, \(Q^s\) increases
each threshold by \(s\); this notation depends on the system and does
not denote a Euclidean neighborhood. We write \(L\coloneqq\log(ep)\).

In the general facts below, \(\R^n\) carries the inner product
\(x\cdot y\) and Euclidean norm \(\norm x\). This dimension is \(r\)
in the support, or \(r-1\) when we work within a facet. We write \(dx\) for Lebesgue measure, \(\varphi_n\) and
\(\Phi_n\) for the standard Gaussian density and law on \(\R^n\), and \(G\)
for a standard Gaussian vector. For a linear subspace \(L\), \(\varphi_L\)
is the standard Gaussian density on \(L\) with respect to Lebesgue measure
on \(L\). In dimension zero, Lebesgue and Gaussian measure both assign
mass one to the origin, and \(\varphi_0(0)\coloneqq1\). This convention
includes the tangential calculation at an endpoint when \(r=1\).
For a nonempty set \(S\), \(\ind_S\) is its indicator function, \(\aff S\)
is the smallest affine subspace containing it, the dimension of \(S\) is
that of \(\aff S\), and \(\relint S\) is the interior of \(S\) relative to
\(\aff S\). The distance from a point \(y\) to a set \(S\) is
\(\dist(y,S)\coloneqq \inf_{x\in S}\norm{y-x}\).

For an integer \(q\ge1\) and \(f\in C^q(\R^n)\), the derivative of order \(q\) in the directions
\(z_1,\ldots,z_q\in\R^n\) is
\[
 D^qf(x)[z_1,\ldots,z_q]
 \coloneqq \partial_{s_1}\cdots\partial_{s_q}f(x+s_1z_1+\cdots+s_qz_q)\Big|_{s=0}
 =\sum_{i_1,\ldots,i_q}\partial_{i_1\cdots i_q}f(x)\,
   z_{1,i_1}\cdots z_{q,i_q},
\]
which is symmetric and multilinear in the directions. We write
\(D^qf(x)[z^{\otimes q}]\) when all directions equal \(z\), with
\(D^0f\coloneqq f\). The first derivative satisfies
\(Df(x)[z]=\nabla f(x)\cdot z\), and we identify the second derivative
with its Hessian matrix. The Loewner order \(M\preceq N\) means that
\(N-M\) is positive semidefinite, \(\norm\Xi_\op\) is the operator norm
of a matrix \(\Xi\), and \(\norm M_*\) is the \emph{nuclear norm}, which for
a symmetric matrix \(M\) with eigenvalues \(\mu_i\) equals \(\sum_i|\mu_i|\).

\subsection{Integration and calculus}\label{tb:sec:calc}

We repeatedly condition on an independent component and then average the
resulting bound. The following theorem specifies when these operations,
and the corresponding exchanges of integrals or sums, are valid.

\begin{fact}[Fubini and Tonelli]\label{tb:fubini}
Let \((S_1,\mu_1)\) and \((S_2,\mu_2)\) be \(\sigma\)-finite measure
spaces and let \(f\) be a measurable function on \(S_1\times S_2\). If
\(f\ge0\), then \(s_1\mapsto\int f(s_1,s_2)\,\mu_2(ds_2)\) is measurable
and
\[
 \int\Bigl(\int f\,d\mu_2\Bigr)d\mu_1
 =\int\Bigl(\int f\,d\mu_1\Bigr)d\mu_2
 =\int f\,d(\mu_1\otimes\mu_2)\in[0,\infty].
\]
If \(\int|f|\,d(\mu_1\otimes\mu_2)<\infty\), the same equalities hold, and
the inner integrals are finite almost everywhere.
\end{fact}

Three special cases occur in the proofs. If \(X\) and \(Y\) are
independent random elements and either \(f\ge0\) or \(f(X,Y)\) is
integrable, then
\[
 \E f(X,Y)=\E g(X),\qquad g(x)\coloneqq \E f(x,Y),
\]
because the joint law is the product of the marginal laws. This is the
formal content of every argument that is carried out for a fixed value of
an independent component such as \(V_0\) and then averaged over it. With
\(\mu_2\) the counting measure, the theorem says that
\(\E\sum_k\xi_k=\sum_k\E\xi_k\) for nonnegative terms, and for terms with
\(\sum_k\E|\xi_k|<\infty\). For an orthogonal decomposition
\(\R^n=L\oplus L^\perp\), Lebesgue measure on \(\R^n\) is the product of the
Lebesgue measures on \(L\) and on \(L^\perp\), so that
\(\int_{\R^n}f\,dx=\int_{L^\perp}\int_Lf(l+w)\,dl\,dw\).

The final averaging in \PR{Lemma}{appx-an:lem:surface}{C.5} is an
instance of the first case. Let \(V_0\in\R^r\) and
\(G\sim N(0,I_r)\) be independent, and let \(\sigma>0\).
For a fixed threshold set \(Q\), it gives
\[
 \E\!\left[R(G)^3\ind\{V_0+\sigma G\in
                   Q^{\sigma/R(G)}\setminus Q\}\right]
 =\E_G\!\left[R(G)^3\Pp_{V_0}\{V_0+\sigma G\in
                   Q^{\sigma/R(G)}\setminus Q\}\right].
\]
Here \(R\) is the projection radius in
Fact~\ref{tb:projection-moments}, and \(Q^s\) shifts every
unit-normal threshold outward by \(s\). Conditional on \(G\), both
the width \(\sigma/R(G)\) and the translation \(\sigma G\) are fixed.
Concentration of \(V_0\) can therefore be applied directly. Tonelli
requires no integrability check at this stage because the integrand
is nonnegative.

\usedin{conditioning on independent components in Appendices~\PN{appx-sec:smoothing}{B}
and~\PN{appx-sec:support}{C}; absolute convergence of the translated-weight
series in \PR{Lemma}{appx-an:lem:restoration}{C.2}; final averaging in
\PR{Lemma}{appx-an:lem:surface}{C.5}; conditioning on independent pairs
in the retained proof in \PR{Appendix}{appx-app:concentration}{D}.}
\emph{Source.} \TBc{Fol99}, Section~2.5.

To differentiate the interpolation, we first justify moving a derivative
inside an expectation and then expand along a line segment. An integrable
bound is needed for the former operation; the latter is ordinary
one-dimensional Taylor expansion applied to a multivariate function.

\begin{fact}[Differentiation under the integral sign; Taylor's formula]
\label{tb:taylor}
(a) Let \(I\) be an open interval and let \(f(t,\omega)\) be
measurable and integrable in \(\omega\) for each \(t\in I\).
Suppose that, outside one null set, it is differentiable in \(t\).
On each compact subinterval \(J\subset I\), assume
\(|\partial_t f(t,\omega)|\le g_J(\omega)\) for all \(t\in J\),
where \(g_J\) is integrable. Then \(\frac d{dt}\E f(t,\cdot)
=\E\,\partial_t f(t,\cdot)\) on \(I\).

(b) For an integer \(q\ge1\), \(f\in C^q(\R^n)\), and \(x,v\in\R^n\),
\[
 f(x+v)=\sum_{k=0}^{q-1}\frac1{k!}D^kf(x)[v^{\otimes k}]
 +\int_0^1\frac{(1-s)^{q-1}}{(q-1)!}D^qf(x+sv)[v^{\otimes q}]\,ds.
\]
For \(q=1\) this is \(f(x+v)-f(x)=\int_0^1Df(x+sv)[v]\,ds\).
\end{fact}

\begin{proof}
For (a), fix an interior time and a compact neighborhood \(J\).
The mean value theorem bounds nearby difference quotients by \(g_J\),
and they converge pointwise, so dominated convergence applies. For (b),
the function \(s\mapsto f(x+sv)\) is of class \(C^q\) with \(k\)th derivative
\(D^kf(x+sv)[v^{\otimes k}]\); apply the one-dimensional Taylor formula with
integral remainder at \(s=0\) and evaluate it at \(s=1\).
\end{proof}

\usedin{the signed Hessian comparison and its Taylor expansion in
Subsection~\PN{appx-an:sec:path}{C.2}; the restoration of the cubic term
to the full path in \PR{Lemma}{appx-an:lem:identity}{C.1}; the segment
integral in \PR{Lemma}{appx-lem:endpoint-gap}{B.3}.}
\emph{Source.} \TBc{Fol99}, Section~2.3.

Differentiating the thinning probability changes one mask at a time.
Here a mask is a variable that retains a summand when it equals one
and deletes it when it equals zero. The following finite-product identity
isolates the change in retention probability; changes in amplitude and
Gaussian variance are separate contributions.

\begin{fact}[Differentiating independent Bernoulli masks]\label{tb:bernoulli}
Let \(B_1,\ldots,B_N\) be independent Bernoulli variables with parameter
\(q\in(0,1)\). For a function \(H:\{0,1\}^N\to\R\), write
\(B^{i,b}\) for the vector obtained by setting its \(i\)th entry to
\(b\in\{0,1\}\). Then
\[
 \frac{d}{dq}\E_q H(B)
 =\sum_{i=1}^N\E_q\{H(B^{i,1})-H(B^{i,0})\}.
\]
If \(q=q(t)\) and \(H=H_t\) are differentiable, then
\[
 \frac{d}{dt}\E_{q(t)}H_t(B)
 =\E_{q(t)}\partial_t H_t(B)
 +q'(t)\sum_i\E_{q(t)}
       \{H_t(B^{i,1})-H_t(B^{i,0})\}.
\]
\end{fact}

\begin{proof}
Temporarily give the masks separate parameters \(q_1,\ldots,q_N\).
Conditioning on all masks except the \(i\)th makes their expectation
affine in \(q_i\), with derivative equal to the insertion--deletion
difference in the first display. Differentiating along
\(q_1=\cdots=q_N=q\) gives the sum. The second formula is the chain
rule. Independent random arguments may be included by conditioning
and then applying Fact~\ref{tb:taylor}(a).
\end{proof}

\usedin{Subsection~\PN{appx-an:sec:path}{C.2}, before
\PR{Lemma}{appx-an:lem:identity}{C.1}, to differentiate the masks in
the Bernoulli--Gauss path.}

In the paper, the amplitude and retention probability are
\(\gamma=1+t\) and \(\varpi=(1-t)/(1+t)^2\), for \(0<t<1\).
The support-coordinate summands \(\zeta_i^0\) are centered, with
\(\sum_i\E\zeta_i^0(\zeta_i^0)^\top=I_r\). Write
\(Y_i\coloneqq\gamma\mathbb B_{i,t}\zeta_i^0\) and
\(U_t\coloneqq\sum_iY_i+\sqrt t\,G\), where all masks and \(G\)
are mutually independent and independent of the summands. To expand the
contribution of summand \(i\), remove it and put \(H_i\coloneqq U_t-Y_i\).
Let \(\zeta_i\) be an independent copy of \(\zeta_i^0\), with all copies
independent of each other and of the path variables. The signed Hessian
comparison in~\PE{appx-an:signed-hessian}{C.6} has scalar weight
\[
 \nu_t(s)\coloneqq \frac{4-2(3-t)s}{1+t},\qquad 0\le s\le1.
\]
Direct integration gives
\[
 \int_0^1\nu_t(s)\,ds=1,\qquad
 \int_0^1s\nu_t(s)\,ds=\frac{2t}{3(1+t)}.
\]
For a twice continuously differentiable scalar function \(g\),
Taylor's formula and the first identity give
\[
 g(0)-\int_0^1\nu_t(s)g(s)\,ds
 =-\frac{2t}{3(1+t)}g'(0)
   -\int_0^1\nu_t(s)\int_0^s(s-u)g''(u)\,du\,ds.
\]
For the logistic approximation \(F=F_{a,h}\) defined below, take
\(g(s)\coloneqq D^2F(H_i+s\gamma\zeta_i)[\zeta_i,\zeta_i]\).
Then \(g'(0)=\gamma D^3F(H_i)[\zeta_i^{\otimes3}]\), so
the factor \(1/2\) in the signed Hessian comparison turns the first
term into
\(-\tfrac12\{2t/(3\gamma)\}\gamma
 D^3F(H_i)[\zeta_i^{\otimes3}]
 =-tD^3F(H_i)[\zeta_i^{\otimes3}]/3\). This explains the
cubic coefficient before the derivatives are restored to the full path.
The weight \(\nu_t\) is signed; these identities use its
integrals, not a probability interpretation.

Repeated differentiation of the logistic approximation produces products
of face probabilities. The partition form of the higher-order chain
rule records which directions differentiate the same inner factor.

\begin{fact}[Fa\`a di Bruno's formula]\label{tb:faa}
Let \(\phi\) be a smooth function on an open interval of \(\R\) and let
\(g:\R^n\to\R\) be smooth with values in that interval. For \(q\ge1\), let
\(\Pi_q\) be the set of partitions \(\mathcal P\) of \(\{1,\ldots,q\}\) into
nonempty blocks, let \(|\mathcal P|\) be the number of blocks, and for a
block \(J\) let \(z_J\) be the list of directions \(z_a\) with
\(a\in J\). Then
\[
 D^q(\phi\circ g)(x)[z_1,\ldots,z_q]
 =\sum_{\mathcal P\in\Pi_q}\phi^{(|\mathcal P|)}(g(x))
  \prod_{J\in\mathcal P}D^{|J|}g(x)[z_J].
\]
\end{fact}

\begin{proof}
For \(q=1\) this is the chain rule. Differentiate the right side for \(q\)
in the direction \(z_{q+1}\). By the product rule, the derivative either
falls on \(\phi^{(|\mathcal P|)}(g(x))\), which gives
\(\phi^{(|\mathcal P|+1)}(g(x))\,Dg(x)[z_{q+1}]\) and adds the block
\(\{q+1\}\) to \(\mathcal P\), or falls on one factor
\(D^{|J|}g(x)[z_J]\), which adds \(q+1\) to the block \(J\). Every partition
of \(\{1,\ldots,q+1\}\) arises exactly once in this way, by removing
\(q+1\) from its block.
\end{proof}

\emph{Source.} \TBc{Com74}; the induction above proves the form used here.

For a finite list of face normals \(u_1,\ldots,u_m\in\R^n\), thresholds
\(a_1,\ldots,a_m\in\R\), and a width \(h>0\), put
\[
 q_\alpha(x)\coloneqq(u_\alpha\cdot x-a_\alpha)/h,\qquad
 w_\alpha(x)\coloneqq e^{q_\alpha(x)},\qquad T(x)\coloneqq\sum_\alpha w_\alpha(x),
\]
\[
 F(x)\coloneqq\frac1{1+T(x)},\qquad
 \pi_\alpha(x)\coloneqq F(x)w_\alpha(x),\qquad
 \theta(x)\coloneqq1-F(x)=\sum_\alpha\pi_\alpha(x).
\]
Here \(F\) is the logistic approximation to the indicator of the
threshold set. The baseline outcome has weight one and probability
\(F\); face \(\alpha\) has probability \(\pi_\alpha\). These
probabilities describe the smoothing formula, not the original event.
One and two differentiations give
\begin{align*}
 DF(x)[z]&=-h^{-1}\sum_\alpha F\pi_\alpha(u_\alpha\cdot z),\\
 D^2F(x)[z_1,z_2]
 &=h^{-2}\left\{2\sum_{\alpha,\beta}F\pi_\alpha\pi_\beta
       (u_\alpha\cdot z_1)(u_\beta\cdot z_2)
       -\sum_\alpha F\pi_\alpha
       (u_\alpha\cdot z_1)(u_\alpha\cdot z_2)\right\}.
\end{align*}
All weights are evaluated at \(x\). In the first formula their sum is
\(F\theta\), which is small away from the smoothed boundary.
The next result retains this factor at every derivative order.

\begin{fact}[Derivatives of the logistic approximation]\label{tb:logistic}
With the notation above, let \(q\ge1\).
For directions \(z_1,\ldots,z_q\), with the blocks \(J_1,\ldots,J_{|\mathcal P|}\)
of \(\mathcal P\) listed in any fixed order,
\[
 D^qF(x)[z_1,\ldots,z_q]
 =h^{-q}\sum_{\mathcal P\in\Pi_q}(-1)^{|\mathcal P|}|\mathcal P|!
 \sum_{\alpha_1,\ldots,\alpha_{|\mathcal P|}}
 F(x)\prod_{b=1}^{|\mathcal P|}\pi_{\alpha_b}(x)
 \prod_{b=1}^{|\mathcal P|}\prod_{a\in J_b}(u_{\alpha_b}\cdot z_a),
\]
Moreover, \(\sum_{\mathcal P\in\Pi_q}|\mathcal P|!\le q!\,2^{q-1}\), and hence
\[
 |D^qF(x)[z_1,\ldots,z_q]|
 \le q!\,2^{q-1}h^{-q}F(x)\theta(x)
       \prod_{a=1}^q\max_\alpha|u_\alpha\cdot z_a|.
\]
\end{fact}

\begin{proof}
Apply Fact~\ref{tb:faa} with \(\phi(\tau)\coloneqq (1+\tau)^{-1}\) and \(g=T\). Then
\(\phi^{(k)}(\tau)=(-1)^kk!\,(1+\tau)^{-k-1}\), and since each \(w_\alpha\) is
an exponential of an affine function,
\(D^lT(x)[z_1,\ldots,z_l]=h^{-l}\sum_\alpha w_\alpha(x)\prod_{a\le l}
(u_\alpha\cdot z_a)\). Substituting, and using
\((1+T)^{-|\mathcal P|-1}\prod_bw_{\alpha_b}=F\prod_b\pi_{\alpha_b}\) and
\(\sum_b|J_b|=q\), gives the formula. The number
\(\sum_{\mathcal P}|\mathcal P|!\) counts ordered partitions
\((J_1,\ldots,J_k)\) of \(\{1,\ldots,q\}\). Writing the elements of
\(J_1\), then of \(J_2\), and so on, in increasing order inside each block,
gives a permutation together with the set of positions at which a new block
starts, which is a subset of the \(q-1\) gaps between consecutive entries.
Every ordered partition arises from at least one of these
\(q!\,2^{q-1}\) pairs. For a partition with \(k\ge1\) blocks, summing the probability
weights gives \(F\theta^k\le F\theta\). Each product of projections
is bounded by \(\prod_a\max_\alpha|u_\alpha\cdot z_a|\).
These two observations prove the stated envelope without a factor
depending on the dimension.
\end{proof}

The early estimate uses this formula for orders three and four. After
the two Gaussian integrations by parts in the late estimate, the explicit
first and second derivatives above suffice. In both cases the factor
\(F\theta\) connects the derivative estimate to concentration near
the boundary.

\usedin{the general derivative expansion and local envelope in
Subsection~\PN{appx-sec:mixture}{B.2}, namely
\PE{appx-eq:partition-derivative}{B.1}--\PE{appx-eq:local-envelope}{B.3};
the explicit first and second derivatives in
\PR{Lemma}{appx-an:lem:restoration}{C.2}.}
\emph{Source.} The formula follows from the partition chain rule in
Fact~\ref{tb:faa}; the proof above is self-contained.

\subsection{Gaussian identities}\label{tb:sec:gauss}

The mean estimate separates a normal coordinate from the tangential
coordinates of a face. Orthogonality makes these Gaussian components
independent; a second elementary identity controls the density change
when we move outward from the face.

\begin{fact}[Gaussian density in orthogonal coordinates]\label{tb:gauss-coord}
(a) Let \(\R^n=L\oplus L^\perp\) be an orthogonal decomposition. Then
\(\varphi_n(l+w)=\varphi_L(l)\varphi_{L^\perp}(w)\) for \(l\in L\) and
\(w\in L^\perp\), so the projections of a standard Gaussian vector onto
\(L\) and \(L^\perp\) are independent standard Gaussian vectors. In
particular, for a unit vector \(u\), \(a\in\R\), and \(y\perp u\),
\(\varphi_n(au+y)=\varphi_1(a)\varphi_{n-1}(y)\).

(b) For \(x,v\in\R^n\),
\(\varphi_n(x+v)=\varphi_n(x)\exp\{-x\cdot v-\tfrac12\norm v^2\}\).
\end{fact}

\begin{proof}
Both follow from \(\norm{l+w}^2=\norm l^2+\norm w^2\) and from expanding
\(\norm{x+v}^2\).
\end{proof}

\usedin{the normal--tangential decomposition in
\PR{Lemma}{appx-an:lem:mean}{C.3}; the facet density factorization
and density comparison along normal tubes in
\PR{Lemma}{appx-an:lem:surface}{C.5}.}

Gaussian integration by parts replaces two directional derivatives by
a centered quadratic function of the Gaussian vector. The factor
\(\sigma^{-2}\) below records the Gaussian noise scale exactly.

\begin{fact}[Gaussian integration by parts]\label{tb:stein}
Let \(G\sim N(0,I_n)\). In (a), assume \(f\in C^1(\R^n)\) and
that \(f\) and its first derivatives have polynomial growth.
In (b), assume in addition that \(f\in C^2\) and its second
derivatives have polynomial growth. Then
\begin{enumerate}
\item[(a)] \(\E[Gf(G)]=\E\nabla f(G)\);
\item[(b)] \(\E[(GG^\top-I_n)f(G)]=\E D^2f(G)\); equivalently, for \(z\in\R^n\),
\(\E D^2f(G)[z,z]=\E\bigl[\tr\{(GG^\top-I_n)zz^\top\}f(G)\bigr]\), where
\(\tr\{(GG^\top-I_n)zz^\top\}=(z\cdot G)^2-\norm z^2\);
\item[(c)] if \(s\ge0\) is an integer, \(v\in\R^n\), \(\sigma>0\),
\(U\coloneqq v+\sigma G\), and
\(F\in C^{s+2}(\R^n)\) has derivatives of polynomial growth, then for
\(z,z_1,\ldots,z_s\in\R^n\),
\[
 \E D^{s+2}F(U)[z,z,z_1,\ldots,z_s]
 =\sigma^{-2}\E\bigl[\tr\{(GG^\top-I_n)zz^\top\}\,
 D^sF(U)[z_1,\ldots,z_s]\bigr].
\]
\end{enumerate}
\end{fact}

\begin{proof}
For (a) in one coordinate, \(x\varphi_1(x)=-\varphi_1'(x)\), so integration by
parts gives
\[\int xf(x)\varphi_1(x)\,dx=\int f'(x)\varphi_1(x)\,dx,\]
and the boundary terms vanish because of the growth condition. The other coordinates
are handled by Fact~\ref{tb:fubini}. For (b), apply (a) to \(\partial_if\) and
to \(G_kf\): \(\E[G_iG_kf(G)]=\E\partial_i(G_kf)(G)=\delta_{ik}\E f(G)
+\E[G_k\partial_if(G)]=\delta_{ik}\E f(G)+\E\partial_{ik}f(G)\). For (c), the
function \(f(g)\coloneqq D^sF(v+\sigma g)[z_1,\ldots,z_s]\) has
\(D^2f(g)[z,z]=\sigma^2D^{s+2}F(v+\sigma g)[z,z,z_1,\ldots,z_s]\); apply (b).
\end{proof}

\usedin{Subsection~\PN{sec:outline-late}{2.4} and
\PR{Lemma}{appx-an:lem:restoration}{C.2}, conditional on \(\mathcal H\)
with \(v=V_0+v_i\) and \(z=Z_i\); the identity
\(\E[Z\cdot\nabla\Psi(Z)]=\E\tr(\Sigma D^2\Psi(Z))\)
in \PR{Lemma}{appx-lem:endpoint-gap}{B.3}.}
The identity also holds conditionally when the base and directions are
measurable with respect to a sigma-field independent of \(G\), provided
the conditional growth assumptions hold. Averaging that identity further
requires integrability of both sides. In the paper the fixed width
\(h>0\), bounded face components, and conditional covariance bound
supply this integrability.

\emph{Source.} Stein~\TBc{Ste81}; the proof above is complete.

Changing the Gaussian variance contributes a Hessian trace to the path
derivative. We give the formula for a general covariance and then use
its standard-Gaussian form in the support coordinates.

\begin{fact}[Heat equation for the Gaussian component]\label{tb:heat}
Let \(Z\sim N(0,\Sigma)\) in \(\R^p\), where \(\Sigma\) may be singular, and
let \(f\in C^2(\R^p)\) be bounded with bounded first and second derivatives.
Then, for \(x\in\R^p\) and \(t>0\),
\[
 \frac d{dt}\E f(x+\sqrt t Z)
 =\tfrac12\E\tr\{\Sigma D^2f(x+\sqrt t Z)\}.
\]
\end{fact}

\begin{proof}
Write \(Z=\Gamma G\) with \(\Gamma\Gamma^\top=\Sigma\) and \(G\) standard
Gaussian. By Fact~\ref{tb:taylor}(a), which applies because the difference
quotients near \(t_0>0\) are bounded by
\(\norm\nabla f_\infty\norm{\Gamma G}/(2\sqrt{t_0/2})\)
when the times lie in \([t_0/2,3t_0/2]\), the derivative equals
\((2\sqrt t)^{-1}\E[G\cdot\Gamma^\top\nabla f(x+\sqrt t\Gamma G)]\).
Fact~\ref{tb:stein}(a) applied to the function \(g\mapsto
\Gamma^\top\nabla f(x+\sqrt t\Gamma g)\) turns this into
\((2\sqrt t)^{-1}\sqrt t\,\E\tr\{\Gamma^\top D^2f\,\Gamma\}
=\frac12\E\tr\{\Sigma D^2f\}\).
\end{proof}

In the support \(\R^r\), the covariance is \(I_r\), so this gives
\(\frac d{dt}\E_G F(v+\sqrt tG)=\frac12\E_G\tr D^2F(v+\sqrt tG)\).
The base \(v\) is held fixed for this differentiation. If its law also
depends on time, its change contributes separately.

\usedin{Subsection~\PN{sec:outline-path}{2.2} and the derivation of
\PE{appx-an:signed-hessian}{C.6} in
Subsection~\PN{appx-an:sec:path}{C.2}. The Gaussian component supplies
the covariance term that cancels the constant term of the Hessian
expansion.}

The surface estimate needs the size of a few normal projections,
not the Euclidean length of the full Gaussian vector. The following
bound explains why only a logarithm of the number of faces enters.

\begin{fact}[Maxima of one- and two-dimensional Gaussian projections]
\label{tb:projection-moments}
Let \(u_1,\ldots,u_m\) be unit vectors in \(\R^r\), with \(m\ge1\).
For linearly independent \(u_j,u_k\), let \(\Pi_{jk}\) be orthogonal
projection onto their span, and put
\[
 \mathcal M(g)\coloneqq \max\left\{\max_j|u_j\cdot g|,
   \max_{j,k:\,u_j,u_k\text{ independent}}\norm{\Pi_{jk}g}\right\},
 \qquad R(g)\coloneqq 1+\mathcal M(g).
\]
An empty second maximum is omitted. Then \(R\) is \(1\)-Lipschitz.
For \(G\sim N(0,I_r)\) and \(q\in\{2,3\}\),
\[
 \E R(G)^q\le C_q\{\log(em)\}^{q/2}.
\]
\end{fact}

\begin{proof}
Every projection is a contraction, so each function in the maximum
is \(1\)-Lipschitz, as is their maximum. A scalar standard normal
satisfies \(\Pp(|Z|>s)\le2e^{-s^2/2}\); the length of a
two-dimensional standard normal satisfies
\(\Pp(\norm{G_2}>s)=e^{-s^2/2}\). A union bound gives
\[
 \Pp\{\mathcal M(G)>s\}\le\min\{1,2(m+m^2)e^{-s^2/2}\}.
\]
The projections need not be independent. Integrate
\(q\int_0^\infty s^{q-1}\Pp\{\mathcal M(G)>s\}\,ds\), splitting
at \(\sqrt{8\log(em)}\), and use
\((1+x)^q\le2^{q-1}(1+x^q)\).
\end{proof}

\usedin{\PE{appx-an:projection-moments}{C.28} in
\PR{Lemma}{appx-an:lem:surface}{C.5}. For the full face list \(m=2p\), and also for its subsystems, these moments
are of orders \(L\) and \(L^{3/2}\), where \(L=\log(ep)\).}

\subsection{Matrix and entropy inequalities}\label{tb:sec:matrix}

Both time ranges use the same observation: the summed second moments
control a pair of projections even when other random coefficients
remain inside the expectation.

\begin{fact}[Conditional covariance contraction]\label{tb:covariance}
Let \(Z_1,\ldots,Z_N\in\R^r\) have finite second moments, and let
\(\mathcal H\) be a sigma-field such that
\(\sum_i\E[Z_iZ_i^\top\mid\mathcal H]\preceq I_r\) almost surely.
For \(\mathcal H\)-measurable unit vectors \(u,v\),
\[
 \sum_i\E[|(u\cdot Z_i)(v\cdot Z_i)|\mid\mathcal H]\le1.
\]
If real random coefficients satisfy \(|c_i|\le a\) almost surely,
where \(a\ge0\) is constant, then
\[
 \left\|\sum_i\E[c_iZ_iZ_i^\top\mid\mathcal H]\right\|_\op\le a.
\]
No independence of the coefficients and vectors is required.
\end{fact}

\begin{proof}
Conditional Cauchy--Schwarz, followed by Cauchy--Schwarz over the
index \(i\), bounds the first sum by
\[
 \left(\sum_i\E[(u\cdot Z_i)^2\mid\mathcal H]\right)^{1/2}
 \left(\sum_i\E[(v\cdot Z_i)^2\mid\mathcal H]\right)^{1/2}\le1.
\]
For the second claim, each unit vector \(u\) gives
\(\bigl|\sum_i\E[c_i(u\cdot Z_i)^2\mid\mathcal H]\bigr|\le a\).
The matrix is symmetric, so taking the supremum of its absolute
quadratic forms proves the operator-norm bound.
\end{proof}

\usedin{\PE{appx-an:pair-contraction}{C.2} and the early estimates in
Subsection~\PN{appx-an:sec:early}{C.3}; the coefficient matrices in
\PR{Lemma}{appx-an:lem:restoration}{C.2}. Bounded face projections
control the remaining directions by powers of \(\delta\).}

After conditioning, the coefficient matrix has bounded operator norm.
The following duality transfers that bound to the Gaussian matrix, whose
nuclear norm we estimate using its positive and negative parts.

\begin{fact}[Duality between the operator and nuclear norms]\label{tb:nuclear}
(a) For every square matrix \(\Xi\) and every symmetric matrix \(M\),
\(|\tr(\Xi M)|\le\norm\Xi_\op\norm M_*\). Moreover
\(\norm M_*=\max\{\tr(\Xi M):\norm\Xi_\op\le1\}\), so that \(\norm\cdot_*\) is
a norm on symmetric matrices.

(b) If \(M=P-N\) with \(P\succeq0\) and \(N\succeq0\), then
\(\norm M_*\le\tr P+\tr N=2\tr P-\tr M\).
\end{fact}

\begin{proof}
For (a), let \(M=\sum_i\mu_ie_ie_i^\top\) be a spectral decomposition. Then
\(\tr(\Xi M)=\sum_i\mu_i\,e_i^\top\Xi e_i\) and \(|e_i^\top\Xi e_i|\le
\norm\Xi_\op\), which gives the inequality. Equality is attained at
\(\Xi\coloneqq \sum_i\operatorname{sgn}(\mu_i)e_ie_i^\top\). A supremum of linear
functionals satisfies the triangle inequality. For (b), the norm of a
positive semidefinite matrix is its trace, so the triangle inequality gives
\(\norm M_*\le\tr P+\tr N\), and \(\tr N=\tr P-\tr M\).
\end{proof}

\usedin{the final contraction in \PR{Lemma}{appx-an:lem:restoration}{C.2};
the reduction \PE{appx-an:matrix-trace-reduction}{C.22} in
\PR{Lemma}{appx-an:lem:matrix}{C.4}, with
\(P=p_\alpha\mu_\alpha\mu_\alpha^\top\) and
\(N=p_\alpha(I_r-K_\alpha)\).}

An individual logistic weight can have arbitrarily small mass, so its
logarithmic bound should be summed before the mass is estimated. The
log-sum inequality performs that summation; its integral version then
allows us to average over the independent base.

\begin{fact}[Log-sum inequality and a concave function]\label{tb:logsum}
(a) For positive numbers \(y_1,\ldots,y_l\) and \(z_1,\ldots,z_l\) with sums
\(y\) and \(z\),
\(\sum_jy_j\log(z_j/y_j)\le y\log(z/y)\).

(b) The function \((x,y)\mapsto x\log(y/x)\) is concave on \((0,\infty)^2\).
Hence, for positive integrable random variables \(X,Y\),
\(\E[X\log(Y/X)]\le\E X\,\log(\E Y/\E X)\).
The expectation on the left may equal \(-\infty\); its positive
part is integrable because \(x\log(y/x)\le y/e\).
\end{fact}

\begin{proof}
For (a), Jensen's inequality for the concave function \(\log\), with the
weights \(y_j/y\), gives \(\sum_j(y_j/y)\log(z_j/y_j)\le\log\sum_j(y_j/y)(z_j/y_j)
=\log(z/y)\). For (b), the function \(f(x,y)\coloneqq x\log y-x\log x\) has Hessian
\(\begin{pmatrix}-1/x&1/y\\1/y&-x/y^2\end{pmatrix}\), whose determinant is
zero and whose trace is negative, so the Hessian is negative semidefinite.
For the expectation bound, put \(\mu\coloneqq\E X\) and
\(\eta\coloneqq\E Y\). The tangent-plane inequality for this concave
function gives
\[
 f(X,Y)\le f(\mu,\eta)
   +f_x(\mu,\eta)(X-\mu)+f_y(\mu,\eta)(Y-\eta).
\]
The right side is integrable. Taking expectations proves the claim,
including when the negative part of the left side has infinite mean.
\end{proof}

\usedin{the reference-mass estimate in
\PR{Lemma}{appx-an:lem:matrix}{C.4}: first sum over tuples sharing
their first face, then average over \(V_0\). Both the reference mass
and the logistic mass contain the factor \(h\sqrt L\), which cancels
inside the logarithm.}
\emph{Source.} \TBc{CT06}, Theorem~2.7.1, for (a).

\subsection{Log-concavity}\label{tb:sec:logconcave}

A function \(f:\R^n\to[0,\infty)\) is \emph{log-concave} if
\(f((1-\alpha)x+\alpha y)\ge f(x)^{1-\alpha}f(y)^\alpha\) for all \(x,y\) and \(\alpha\in[0,1]\);
when \(f>0\) this says that \(\log f\) is concave. At the endpoints
\(\alpha=0,1\) the inequality means equality, without a convention
for \(0^0\). We need the fact that integrating out coordinates
preserves log-concavity: it will turn a Gaussian convolution into a
concave log normalizing constant.

\begin{fact}[Pr\'ekopa--Leindler inequality and Pr\'ekopa's theorem]
\label{tb:prekopa}
(a) Let \(0<\alpha<1\) and let \(f,g,h:\R^n\to[0,\infty)\) be bounded,
continuous, and integrable, with
\(h((1-\alpha)x+\alpha y)\ge f(x)^{1-\alpha}g(y)^\alpha\) for all \(x,y\). Then
\(\int h\ge(\int f)^{1-\alpha}(\int g)^\alpha\).

(b) Let \(f:\R^a\times\R^b\to[0,\infty)\) be bounded, continuous, and
log-concave, with \(f(z,\cdot)\) integrable for every \(z\). Then
\(z\mapsto\int f(z,y)\,dy\) is log-concave on \(\R^a\).
\end{fact}

\begin{proof}
Part (a) is the Pr\'ekopa--Leindler inequality; we use the standard
integrable-function theorem in \TBc{Gar02}, Theorem~7.1.
To deduce (b), fix \(z_0,z_1\) and \(0<\alpha<1\), and put
\(z_\alpha\coloneqq(1-\alpha)z_0+\alpha z_1\).
Log-concavity gives
\[
 f(z_\alpha,(1-\alpha)y_0+\alpha y_1)
 \ge f(z_0,y_0)^{1-\alpha}f(z_1,y_1)^\alpha.
\]
The three sections are bounded, continuous, and integrable by hypothesis,
so (a) applies to them. Its conclusion is exactly the log-concavity of
\(z\mapsto\int f(z,y)\,dy\). The endpoint cases are immediate.
\end{proof}

\emph{Source.} Pr\'ekopa~\TBc{Pre71,Pre73} and Leindler~\TBc{Lei72}.
The marginalization result is Theorem~6 of \TBc{Pre73};
\TBc{Gar02}, Section~7, treats the integral inequality.

We now apply marginal log-concavity to a Gaussian convolution. Its
logarithmic Hessian is the covariance of the tilted law minus the
identity, so concavity gives the matrix inequality needed in the proof.

\begin{fact}[Gaussian laws tilted by a log-concave weight]\label{tb:tilt}
Let \(w:\R^r\to[0,\infty)\) be bounded, continuous, and log-concave, with
\(p\coloneqq \E w(G)>0\) for \(G\sim N(0,I_r)\). Let \(\mu\) and \(K\) be the mean and
covariance matrix of the probability law with density \(w(g)\varphi_r(g)/p\).
Then \(K\preceq I_r\), and
\(\E[(GG^\top-I_r)w(G)]=p(K-I_r+\mu\mu^\top)\).
\end{fact}

\begin{proof}
The function \((z,g)\mapsto e^{-\norm{g-z}^2/2}w(g)\) is log-concave, since
\(\norm{g-z}^2\) is jointly convex in \((z,g)\), and it is bounded and
continuous, with an integrable \(g\)-section for every \(z\).
The integral is positive because \(p>0\) and the Gaussian kernel is
strictly positive. By Fact~\ref{tb:prekopa}(b), \(\ell(z)\coloneqq \log\int e^{-\norm{g-z}^2/2}
w(g)\,dg\) is concave. Since \(\norm{g-z}^2=\norm g^2-2z\cdot g+\norm z^2\),
\[
 \ell(z)=\text{const}-\tfrac12\norm z^2+\log M(z),\qquad
 M(z)\coloneqq \int e^{z\cdot g}w(g)\varphi_r(g)\,dg=p\,\E_\nu e^{z\cdot g},
\]
where \(\nu\) is the tilted law. The function \(M\) is finite and smooth: on every bounded set of
\(z\)'s, each derivative of the integrand is bounded by a constant
times \((1+\norm g)^k e^{c\norm g}\varphi_r(g)\), an integrable function.
Thus differentiation under the integral is valid, and the Hessian of \(\log M\) at \(0\) is the
covariance \(K\) of \(\nu\). A concave function of class \(C^2\) has a
negative semidefinite Hessian, so \(K-I_r\preceq0\). The identity is the
decomposition of the second moment \(\E_\nu gg^\top=K+\mu\mu^\top\).
\end{proof}

\usedin{the Gaussian-tilt calculation in
\PR{Lemma}{appx-an:lem:matrix}{C.4}, where \(w=w_\alpha\) and
\(K=K_\alpha\), and its explanation in
Subsection~\PN{sec:outline-late}{2.4}.}
\emph{Source.} The covariance contraction is a classical consequence of
Pr\'ekopa's marginalization theorem~\TBc{Pre73}, Theorem~6, as the
preceding proof shows. For smooth positive weights, it also follows
from the Brascamp--Lieb variance inequality~\TBc{BL76}, applied to
linear functions under the potential \(\norm g^2/2-\log w(g)\), whose
Hessian is at least \(I_r\). The same convolution argument covers
indicators of convex sets by the measurable form of Pr\'ekopa's
theorem. This convex-cell specialization is also used in \TBc{10}.

It remains to check that every weight produced by the derivative
expansion meets the preceding hypothesis. Its logarithm has a particularly
simple form: an affine function minus a multiple of a log-sum-exponential.

\begin{fact}[Log-concavity of the logistic weights]\label{tb:lc-weights}
With the notation of Fact~\ref{tb:logistic}, for every \(k\ge0\) and every
tuple \((\alpha_1,\ldots,\alpha_k)\), the function
\(F\pi_{\alpha_1}\cdots\pi_{\alpha_k}=\exp\{\sum_{a\le k}q_{\alpha_a}\}
(1+T)^{-k-1}\), where \(q_\alpha\coloneqq (u_\alpha\cdot x-a_\alpha)/h\), is
log-concave on \(\R^n\), and so is its composition
\(g\mapsto F\pi_{\alpha_1}\cdots\pi_{\alpha_k}(v+\sigma g)\) with any affine map.
\end{fact}

\begin{proof}
The logarithm is a sum of the affine functions \(q_{\alpha_a}\) and
\(-(k+1)\log(1+\sum_\alpha e^{q_\alpha})\). The function
\(\log(1+\sum_\alpha e^{q_\alpha})\) is convex. To include the baseline
term explicitly, put \(q_0\coloneqq0\). H\"older's inequality gives
\[
 \sum_{\alpha=0}^m e^{(1-s)q_\alpha(x)+s q_\alpha(y)}
 \le\left(\sum_{\alpha=0}^m e^{q_\alpha(x)}\right)^{1-s}
     \left(\sum_{\alpha=0}^m e^{q_\alpha(y)}\right)^s,
 \qquad 0<s<1.
\]
Taking logarithms proves the assertion. Composition
with an affine map preserves concavity.
\end{proof}

\usedin{the application of covariance contraction to every
\(w_\alpha=F\prod_a\pi_{\alpha_a}\) in
\PR{Lemma}{appx-an:lem:matrix}{C.4}.}
\emph{Source.} Convexity of log-sum-exp is standard; see
Boyd and Vandenberghe~\TBc{BV04}, Section~3.1.5, pp.~72--74.
The baseline term \(1=e^{q_0}\), with \(q_0=0\), is included in that sum.

\subsection{Polyhedra: faces, normal cones, and nearest points}\label{tb:sec:poly}

The surface estimate treats a fixed polyhedron geometrically. Unlike
the smoothing formula, its boundary integral depends only on the set,
so repeated and redundant inequalities may now be removed.
A \emph{polyhedron} is a set
\[
 Q\coloneqq \bigcap_{j=1}^m\{x\in\R^n:u_j\cdot x\le a_j\},
\]
with unit vectors \(u_j\), called \emph{face normals}, and numbers
\(a_j\), called \emph{offsets}. Each face normal points outward.
The polyhedron is closed and convex. It is
\emph{full-dimensional} if it has an interior point, and \emph{proper} if it
is neither empty nor \(\R^n\). The description is \emph{irredundant} if
deleting any one inequality changes the set; deleting redundant or repeated
inequalities one at a time turns every description into an irredundant one
of the same set. For \(j\le m\), put \(H_j\coloneqq \{x:u_j\cdot x=a_j\}\) and
\(F_j\coloneqq Q\cap H_j\). A \emph{face} of \(Q\) is \(Q\) itself or a set
\(Q\cap\{u\cdot x=c\}\) with \(u\ne0\) and \(Q\subseteq\{u\cdot x\le c\}\).
A \emph{facet} is a face of dimension \(n-1\), and a \emph{ridge}, when \(n\ge2\), is a face
of dimension \(n-2\). In the plane the facets are the sides of a polygon
and the ridges its corners, in \(\R^3\) the facets are the faces of a solid
and the ridges its edges, and on the line the facets are the endpoints of an
interval, which has no ridges. We use these standard properties of faces:
the relative interiors of the faces of \(Q\) are pairwise disjoint and cover
\(Q\); a face of a face of \(Q\) is a face of \(Q\); and a face \(G'\)
of a face \(G\) with \(G'\ne G\) lies in the relative boundary
\(G\setminus\relint G\) of \(G\) (\TBc{Roc70}, Section~18; \TBc{Zie95}, Chapter~2).

The next fact identifies the pieces that contribute to the boundary
calculation. In particular, every ridge belongs to exactly two facets;
this allows us to combine their coefficients before taking absolute values.

\begin{fact}[Facets and ridges of a polyhedron]\label{tb:facets}
Let \(Q\) be a proper, full-dimensional polyhedron in \(\R^n\) with an
irredundant description. Statements about ridges are omitted when \(n=1\).
\begin{enumerate}
\item[(a)] The face normals \(u_1,\ldots,u_m\) are pairwise distinct, each
\(F_j\) is a facet, \(Q\) has no other facets, and
\(\partial Q=F_1\cup\cdots\cup F_m\).
\item[(b)] For \(j\ne k\), \(F_j\cap F_k\subseteq H_j\cap H_k\), and
\(H_j\cap H_k\) is empty if \(u_k=-u_j\), and an affine subspace of
dimension \(n-2\) otherwise. In particular \(F_j\cap F_k\) lies in a
finite union of affine subspaces of dimension at most \(n-2\).
\item[(c)] Each \(F_j\) is a full-dimensional polyhedron in the hyperplane
\(H_j\). Write the points of \(H_j\) as \(a_ju_j+y\) with \(y\perp u_j\), and,
for \(k\ne j\) with \(u_k\ne\pm u_j\), put
\(\alpha_{jk}\coloneqq \arccos(u_j\cdot u_k)\in(0,\pi)\),
\[
 \nu_{k\mid j}\coloneqq \frac{u_k-(u_j\cdot u_k)u_j}{\sin\alpha_{jk}},\qquad
 d_{k\mid j}\coloneqq \frac{a_k-(u_j\cdot u_k)a_j}{\sin\alpha_{jk}} .
\]
Then \(F_j\) is described in the coordinate \(y\) by the inequalities
\(\nu_{k\mid j}\cdot y\le d_{k\mid j}\), and its facets are exactly the
ridges \(E_{jk}\coloneqq F_j\cap F_k\) of \(Q\) of dimension \(n-2\); the outward
unit normal of \(F_j\) at \(E_{jk}\) is \(\nu_{k\mid j}\).
\item[(d)] Every ridge of \(Q\) lies in exactly two facets, so the ridges
correspond to the unordered pairs \(\{j,k\}\) with
\(\dim(F_j\cap F_k)=n-2\); for such a pair \(u_k\ne\pm u_j\), and the
only inequalities active at a point of \(\relint E_{jk}\) are \(j\) and \(k\).
\end{enumerate}
\end{fact}

\begin{proof}
(a) Fix \(j\). Irredundancy gives a point \(x_j\) with \(u_j\cdot x_j>a_j\)
and \(u_k\cdot x_j\le a_k\) for \(k\ne j\). Let \(x_0\) be an interior point of
\(Q\), and let \(x_*\) be the point where the segment from \(x_0\) to \(x_j\)
crosses \(H_j\). For \(k\ne j\), \(u_k\cdot x_*<a_k\), because \(x_*\) is a
convex combination of \(x_0\), where the inequality is strict, with positive
weight, and of \(x_j\), where it holds weakly. Hence a neighborhood of
\(x_*\) in \(H_j\) lies in \(F_j\), and \(\dim F_j=n-1\). Two inequalities
with the same normal cannot both be irredundant, since the one with the
larger offset is implied by the other. A point of \(Q\) at
which no inequality is active is interior, and at a point where
inequality \(j\) is active the points \(x+\varepsilon u_j\) lie outside \(Q\),
so \(\partial Q=\bigcup_jF_j\). If \(G=Q\cap\{u\cdot x=c\}\) is a facet and
\(x\in\relint G\), some inequality \(j\) is active at \(x\); on \(G\) the
affine function \(u_j\cdot x-a_j\le0\) attains its maximum at a relative
interior point and is therefore zero, so \(G\subseteq H_j\), and since \(G\) has
dimension \(n-1\), \(G=Q\cap H_j=F_j\).

(b) If \(u_k=-u_j\) and \(H_j\cap H_k\ne\emptyset\), then \(Q\) would lie in
the hyperplane \(H_j\), contradicting full dimension. If \(u_k\ne\pm u_j\),
the two equations are independent.

(c) Suppose \(n\ge2\). The set \(F_j\) is the polyhedron in \(H_j\) cut out by the remaining
inequalities, which take the stated form in the coordinate \(y\) because
\(u_k\cdot(a_ju_j+y)=a_ju_j\cdot u_k+(u_k-(u_j\cdot u_k)u_j)\cdot y\); an
inequality with \(u_k=-u_j\) holds automatically on \(H_j\) since
\(F_j\ne\emptyset\). After removing inequalities redundant within \(H_j\), part (a)
applied there shows that the facets of \(F_j\) are
the faces of \(F_j\) of dimension \(n-2\), hence faces of \(Q\) of that
dimension, that is, ridges of \(Q\) in \(F_j\); each equals
\(F_j\cap H_k=F_j\cap F_k\) for the inequality \(k\) that defines it.

(d) Let \(E\) be a ridge and \(y\in\relint E\), and let \(J\) be the set of
inequalities active at \(y\). On \(E\) each affine function
\(u_i\cdot x-a_i\), \(i\in J\), is at most zero and vanishes at the relative
interior point \(y\), so it vanishes on \(E\); thus \(E\subseteq F_i\) for
\(i\in J\) and \(\dim\operatorname{span}\{u_i:i\in J\}\le2\). A single active
inequality is impossible, since then \(y\) would lie in the relative interior
of \(F_j\) and the only face of \(F_j\) containing it is \(F_j\) itself, of
dimension \(n-1\). Opposite face normals \(u_i=-u_j\) are excluded by (b). Hence
the active face normals span a plane \(\Pi\). Near \(y\), \(Q\) coincides with
\(y\) plus the product of \(\Pi^\perp\) and a planar cone \(K\subseteq\Pi\)
with nonempty interior, defined by the active inequalities and contained in
a wedge of opening \(\pi-\alpha_{jk}<\pi\); this cone has exactly two
boundary rays. Each active inequality \(i\) has \(\dim F_i=n-1\), so
\(F_i\) meets every neighborhood of \(y\) in a set of dimension \(n-1\), and
the line \(\{u_i\cdot v=0\}\) must therefore carry a boundary ray of \(K\).
With distinct face normals this leaves exactly two active inequalities.
\end{proof}

\usedin{the facet--ridge decomposition leading to
\PE{appx-an:surface-identity}{C.30} and the pairing of the two terms
attached to each ridge in \PR{Lemma}{appx-an:lem:surface}{C.5}.}
\emph{Source.} These are standard facts about faces of polyhedra; see
\TBc{Roc70}, Sections~18 and~19, and \TBc{Zie95}, Chapter~2.
The planar-cone argument above explains why precisely two facets
meet at a ridge, including for unbounded polyhedra.

To attach a tube to a face, we must know which outward directions
keep that face as the nearest part of the polyhedron. These directions
form its normal cone, determined by the active inequalities.

\begin{fact}[Normal cones]\label{tb:normalcone}
Let \(Q\) be a polyhedron as in Fact~\ref{tb:facets}. For \(x\in Q\) let
\(J(x)\coloneqq \{j:u_j\cdot x=a_j\}\), and let
\[
 N_Q(x)\coloneqq \{v\in\R^n:v\cdot(q-x)\le0\text{ for all }q\in Q\}
\]
be the \emph{normal cone} of \(Q\) at \(x\). Then
\(N_Q(x)=\{\sum_{j\in J(x)}c_ju_j:c_j\ge0\}\). In particular
\(N_Q(x)=\{0\}\) at interior points, \(N_Q(x)=\{cu_j:c\ge0\}\) for
\(x\in\relint F_j\), and \(N_Q(x)=\{c_ju_j+c_ku_k:c_j,c_k\ge0\}\), a sector of
angle \(\alpha_{jk}\) in the plane spanned by \(u_j\) and \(u_k\), for
\(x\in\relint E_{jk}\).
\end{fact}

\begin{proof}
Put \(T\coloneqq \{d:u_j\cdot d\le0\text{ for }j\in J(x)\}\). A vector \(v\) lies in
\(N_Q(x)\) if and only if \(v\cdot d\le0\) for all \(d\in T\): if \(d\in T\)
then \(x+\varepsilon d\in Q\) for small \(\varepsilon>0\), because the
inactive inequalities are strict at \(x\), and conversely \(q-x\in T\) for
\(q\in Q\). By Farkas's lemma, which says that \(v\cdot d\le0\) for all
\(d\) with \(u_j\cdot d\le0\) (\(j\in J\)) if and only if \(v\) is a
nonnegative combination of the \(u_j\), \(j\in J\), this is the stated
cone. For \(x\in\relint F_j\), \(J(x)=\{j\}\), since an index \(k\ne j\) in
\(J(x)\) would put \(x\) in \(F_j\cap F_k\), a face of \(F_j\) of dimension
at most \(n-2\), and so in the relative boundary of \(F_j\). For
\(x\in\relint E_{jk}\), \(J(x)=\{j,k\}\) by Fact~\ref{tb:facets}(d).
\end{proof}

\usedin{the choice of the normal interval over a facet and the normal
sector over a ridge in \PR{Lemma}{appx-an:lem:surface}{C.5}.}
\emph{Source.} \TBc{Roc70}, Sections~22 and~23, for Farkas's lemma and normal
cones; \TBc{RW98}, Theorem~6.46.

The nearest-point characterization makes the outward tubes disjoint.
It applies to any nonempty closed convex set; for such a set we use
\(N_Q(x)=\{v:v\cdot(q-x)\le0\text{ for all }q\in Q\}\).

\begin{fact}[Nearest point]\label{tb:projection}
Let \(Q\subseteq\R^n\) be nonempty, closed, and convex. For each
\(y\in\R^n\) there is a unique point \(\prj_Q(y)\in Q\) with
\(\norm{y-\prj_Q(y)}=\dist(y,Q)\), and \(x=\prj_Q(y)\) if and only if
\(x\in Q\) and \(y-x\in N_Q(x)\). Consequently every \(y\in\R^n\) can be
written in exactly one way as \(y=x+v\) with \(x\in Q\) and \(v\in N_Q(x)\),
and then \(\norm v=\dist(y,Q)\). A direction \(v\) lies in \(N_Q(x)\) if and
only if \(x\) is the nearest point of \(x+s v\) in \(Q\) for every \(s>0\).
\end{fact}

\begin{proof}
For \(x,q\in Q\), \(\norm{y-q}^2=\norm{y-x}^2+\norm{x-q}^2-2(y-x)\cdot(q-x)\).
If \(y-x\in N_Q(x)\), the last term is nonnegative, so
\(\norm{y-q}^2\ge\norm{y-x}^2+\norm{x-q}^2>\norm{y-x}^2\) for \(q\ne x\); thus
\(x\) is the unique nearest point. Conversely, if \((y-x)\cdot(q-x)>0\) for
some \(q\in Q\), then moving from \(x\) toward \(q\) decreases the distance
to \(y\) for small steps, by the displayed identity applied to
\(x+\varepsilon(q-x)\in Q\), so \(x\) is not nearest. For existence, choose \(q_0\in Q\). A minimizing sequence may be
restricted to the closed ball centered at \(y\) with radius
\(\norm{y-q_0}+1\); it therefore has a convergent subsequence in
finite dimension. Closedness puts its limit in \(Q\), and continuity
of the norm makes that limit a minimizer. The last statement follows because \(N_Q(x)\) is a
cone.
\end{proof}

\usedin{disjointness of the variable-width tubes in
\PR{Lemma}{appx-an:lem:surface}{C.5}; each tube point has its nearest
point in the relative interior of exactly one face.}
\emph{Source.} \TBc{RW98}, Example~2.25, for existence and uniqueness,
and Proposition~6.17 for the characterization by the normal cone.

\subsection{Surface measure, Gaussian surface integrals, and the
divergence theorem}\label{tb:sec:surface}

We integrate each facet with its ordinary Euclidean area measure.
Hausdorff measure gives a common notation for facets and ridges in
all dimensions and shows why their lower-dimensional overlaps do not
affect a surface integral.

\begin{fact}[Hausdorff measure on flat pieces]\label{tb:hausdorff}
For \(0\le k\le n\), the \(k\)-dimensional Hausdorff measure \(\Haus^k\) on
\(\R^n\) is obtained from the covering formula
\[
 \Haus^k(S)\coloneqq \lim_{\epsilon\downarrow0}\inf\Bigl\{\sum_i\omega_k
 \Bigl(\tfrac12\operatorname{diam}C_i\Bigr)^k:S\subseteq\bigcup_iC_i,\
 \operatorname{diam}C_i\le\epsilon\Bigr\},
\]
where the infimum runs over countable covers and \(\omega_k\) is
the volume of the unit ball of \(\R^k\).
For \(k=0\), every nonempty covering set contributes one, including
a singleton. The formula defines an outer measure, whose restriction
to Borel sets is the measure used here. (a) On an
affine subspace of dimension \(k\), \(\Haus^k\) is \(k\)-dimensional Lebesgue
measure; in particular \(\Haus^0\) is counting measure and \(\Haus^n\) is
Lebesgue measure on \(\R^n\). (b) \(\Haus^k\) is invariant under isometries.
(c) A set contained in a finite union of affine subspaces of dimension less
than \(k\) has \(\Haus^k\)-measure zero.
\end{fact}

Part (c) follows from (a), since each such subspace lies in a
\(k\)-dimensional affine subspace in which it is a Lebesgue null set. For a
proper, full-dimensional polyhedron \(Q\) with an irredundant description,
and for a function \(f\ge0\) or \(f\) integrable on \(\partial Q\), the
\emph{surface integral} is
\[
 \int_{\partial Q}f\,d\Haus^{n-1}=\sum_{j=1}^m\int_{F_j}f\,d\Haus^{n-1},
\]
where on each facet \(\Haus^{n-1}\) is Lebesgue measure on the hyperplane
\(H_j\) by part~(a). By Fact~\ref{tb:facets}(b) and part~(c), distinct
facets overlap in sets of \(\Haus^{n-1}\)-measure zero, so the sum is the
integral of \(f\) over the set \(\partial Q\) with respect to
\(\Haus^{n-1}\). It depends on the set \(Q\) and not on its description,
which justifies deleting redundant inequalities before computing it.
The Gaussian surface integral \(\surf(Q)=\int_{\partial Q}\Hh\,\varphi_r\,
d\Haus^{r-1}\) of Subsection~\PN{sec:outline-late}{2.4} is the case
\(n=r\) and \(f=\Hh\varphi_r\), with \(\Hh(g)\coloneqq \norm g^2-r\).
The integral is signed; we take its absolute value only after summing
the facet integrals. We set \(\surf(\varnothing)=\surf(\R^r)\coloneqq0\).
Lower-dimensional polyhedra are excluded from this definition and are
handled as exceptional levels in Fact~\ref{tb:variation}.

\usedin{the definition of \(\surf\) in
Subsection~\PN{sec:outline-late}{2.4} and
\PR{Lemma}{appx-an:lem:surface}{C.5}.}
\emph{Source.} \TBc{EG15}, Chapter~2.

Integrating a tangential divergence over a facet converts its interior
integral into contributions from its ridges. We prove the needed divergence
formula by slicing, then use a cutoff to allow an unbounded polyhedron
under precisely the stated decay assumptions.

\begin{samepage}
\begin{fact}[Gauss--Green formula (divergence theorem) for a polyhedron]\label{tb:divergence}
Let \(Q\subseteq\R^n\) (\(n\ge1\)) be a proper, full-dimensional polyhedron
with an irredundant description, facets \(F_1,\ldots,F_m\), and unit
face normals \(u_1,\ldots,u_m\). Let \(X:\R^n\to\R^n\) be of class \(C^1\), with
\(\norm{X(x)}+|\operatorname{div}X(x)|\le C(1+\norm x)^k\varphi_n(x)\) for
some \(C\) and \(k\). Then
\[
 \int_Q\operatorname{div}X\,dx=\sum_{j=1}^m\int_{F_j}X\cdot u_j\,d\Haus^{n-1}.
\]
For \(n=1\) the facets are the endpoints of an interval, \(\Haus^0\) is
counting measure, and the identity is the fundamental theorem of calculus.
\end{fact}
\end{samepage}

\begin{proof}
We first prove the formula for a compactly supported \(C^1\) vector
field. For a unit vector \(e\) and a compactly supported
\(C^1\) scalar function \(f\), we claim that
\begin{equation}
 \int_Q\partial_ef\,dx=\sum_j(e\cdot u_j)\int_{F_j}f\,d\Haus^{n-1},
 \tag{\(*\)}
\end{equation}
since the theorem follows by taking \(e=e_i\) and \(f=X_i\) and summing over
\(i\). Write \(x=w+s e\) with \(w\in e^\perp\). By Fubini's theorem,
\(\int_Q\partial_ef\,dx=\int_{e^\perp}\int_{I(w)}\partial_sf(w+s e)\,ds\,dw\),
where \(I(w)\coloneqq \{s:w+s e\in Q\}\) is a closed interval, possibly empty, a point,
or unbounded, because \(Q\) is convex. Since \(f\) tends to zero along every
line, the inner integral equals \(f(w+s_+(w)e)-f(w+s_-(w)e)\), where
\(s_\pm(w)\) are the endpoints of a nonempty \(I(w)\).
An empty interval contributes zero, and a term with an infinite
endpoint is read as zero.

Consider the upper endpoints \(x=w+s_+(w)e\). If \(x\in\relint F_j\), then
only the inequality \(j\) is active at \(x\) (Fact~\ref{tb:normalcone}), so
\(x+s e\in Q\) for small \(s>0\) unless \(e\cdot u_j>0\), and \(x\) is an upper
endpoint if and only if \(e\cdot u_j>0\). The upper endpoints that lie in no
\(\relint F_j\) lie in \(\bigcup_{j\ne k}F_j\cap F_k\), which is a subset of a
finite union of affine subspaces of dimension at most \(n-2\)
(Fact~\ref{tb:facets}(b)); its projection along \(e\) onto \(e^\perp\) is a
Lebesgue null set of \(e^\perp\). For fixed \(j\) with \(e\cdot u_j>0\), the
map \(x\mapsto x-(x\cdot e)e\) is injective on the set of upper endpoints in
\(\relint F_j\), since a line meets \(Q\) in one interval, and it is the
restriction to \(H_j\) of a linear map whose Jacobian with respect to
\(\Haus^{n-1}\) on \(H_j\) is \(e\cdot u_j\): if \(b_1,\ldots,b_{n-1}\) is an
orthonormal basis of \(u_j^\perp\), the volume of the projection along \(e\)
of the cube they span equals the volume of the parallelepiped spanned by
\(b_1,\ldots,b_{n-1},e\), which is \(|e\cdot u_j|\). Therefore
\(\int f(w+s_+(w)e)\,dw=\sum_{j:\,e\cdot u_j>0}(e\cdot u_j)\int_{F_j}f\,d\Haus^{n-1}\),
the sets \(F_j\setminus\relint F_j\) being \(\Haus^{n-1}\)-null. The same
argument for the lower endpoints gives
\(\int f(w+s_-(w)e)\,dw
=\sum_{j:\,e\cdot u_j<0}|e\cdot u_j|\int_{F_j}f\,d\Haus^{n-1}\).
Subtracting this lower-endpoint integral from the upper-endpoint
integral proves \((*)\). Facets with \(e\cdot u_j=0\) contribute
nothing. Summing over coordinate directions proves the compactly
supported vector-field formula.

For the stated field \(X\), choose smooth cutoffs \(\chi_R\) equal to
one on the ball of radius \(R\), zero outside the ball of radius
\(2R\), with \(0\le\chi_R\le1\) and
\(\norm{\nabla\chi_R}\le C/R\). Apply the formula to \(\chi_RX\).
The product rule gives
\[
 \operatorname{div}(\chi_RX)
 =\chi_R\operatorname{div}X+\nabla\chi_R\cdot X.
\]
The first term converges in integral by the assumed integrable bound;
the second has integral in absolute value at most
\((C/R)\int_{\{\norm x\ge R\}}\norm{X(x)}\,dx\to0\).
On each of the finitely many facets, polynomial times Gaussian decay
is integrable, so dominated convergence also applies to the boundary
integrals. This proves the formula without requiring separate bounds
on the component derivatives of \(X\).
\end{proof}

\usedin{the divergence calculation inside a facet in
\PR{Lemma}{appx-an:lem:surface}{C.5}. The vector field is
\(X(y)\coloneqq y\varphi_{r-1}(y)\), with divergence
\(\{r-1-\norm y^2\}\varphi_{r-1}(y)\); the boundary terms are the
ridge contributions in \PE{appx-an:surface-identity}{C.30}.}
\emph{Source.} This is the classical Gauss--Green formula specialized
to a polyhedron; see \TBc{EG15}, Chapter~5. The cutoff argument above
extends the compact-support formula to the stated Gaussian-decay
hypothesis on an unbounded polyhedron.

For the surface integral used in the paper, this divergence
calculation produces coefficients whose separate expressions contain
\(1/\sin\alpha_{jk}\). We combine the two contributions of a ridge
before estimating them. Let
\(s_j\coloneqq \int_{F_j}\varphi_r\,d\Haus^{r-1}\),
\(\rho_{jk}\coloneqq \int_{E_{jk}}\varphi_r\,d\Haus^{r-2}\), and let
\(z_{jk}\) be the projection of a ridge point onto
\(\operatorname{span}(u_j,u_k)\). Applying the divergence formula in
each facet gives
\[
 \surf(Q)=\sum_j(a_j^2-1)s_j
       -\sum_{\{j,k\}}(d_{k\mid j}+d_{j\mid k})\rho_{jk}.
\]
The offsets \(d_{k\mid j}\) from Fact~\ref{tb:facets}(c) satisfy
\[
 d_{k\mid j}+d_{j\mid k}
 =2\sin(\alpha_{jk}/2)\,
       \frac{u_j+u_k}{\norm{u_j+u_k}}\cdot z_{jk},\qquad
 |d_{k\mid j}+d_{j\mid k}|\le\alpha_{jk}\norm{z_{jk}}.
\]
Indeed, their sum is
\((1-\cos\alpha_{jk})(a_j+a_k)/\sin\alpha_{jk}\), and
\(a_j+a_k=(u_j+u_k)\cdot z_{jk}\).
The inverse sine cancels. The remaining factor \(\alpha_{jk}\)
is precisely the angle of the normal sector whose mass we use next.

\subsection{Tubes around a polyhedron}\label{tb:sec:tubes}

We now compare the paired boundary terms with Gaussian mass outside
the polyhedron. The Gaussian density changes rapidly when a boundary point has a large
normal projection. We therefore shorten the outward tube at that point.
Its width may vary along a face: the required integration formula still
follows from orthogonal coordinates and Fubini's theorem.

Let \(Q\) be a proper, full-dimensional polyhedron in \(\R^n\) with
facets \(F_j\), ridges \(E_{jk}\), and an irredundant description by
unit face normals \(u_j\). For a ridge, put
\(L_{jk}\coloneqq \operatorname{span}(u_j,u_k)\) and
\(z_{jk}\coloneqq \Pi_{jk}x\), for any \(x\in E_{jk}\). This vector is constant
on the ridge because
\(u_j\cdot z_{jk}=a_j\) and \(u_k\cdot z_{jk}=a_k\).
Let \(\eta(x)>0\) be a measurable function on the facets and ridges.
We use the same value of \(\eta\) at a point whenever it belongs to
more than one closed face; only relative interiors enter the tubes.
The outward tubes are
\begin{align*}
 \mathcal T_j&\coloneqq \{x+s u_j:x\in\relint F_j,\ 0<s<\eta(x)\},\\
 \mathcal T_{jk}&\coloneqq \{x+v:x\in\relint E_{jk},\ v\in N_Q(x),\
                         0<\norm v<\eta(x)\}.
\end{align*}
For \(x\in\relint E_{jk}\), the normal cone is a sector of angle
\(\alpha_{jk}\) in \(L_{jk}\). Write
\(S_{jk}(a)\coloneqq N_Q(x)\cap\{v:0<\norm v<a\}\); its area is
\(\alpha_{jk}a^2/2\). For \(Q=\bigcap_j\{u_j\cdot x\le a_j\}\),
the notation \(Q^s\) increases every threshold by \(s\); it does not
denote a Euclidean neighborhood.

\begin{fact}[Variable-width normal tubes]\label{tb:tubes}
(a) The tubes are measurable, and tubes over distinct facets and
ridges are pairwise disjoint.
Each tube point \(g=x+v\) has nearest point \(x\) in \(Q\), and
\(\dist(g,Q)=\norm v\).

(b) For a nonnegative measurable function \(f\),
\begin{align*}
 \int_{\mathcal T_j}f(g)\,dg
 &=\int_{F_j}\int_0^{\eta(x)} f(x+s u_j)\,ds\,d\Haus^{n-1}(x),\\
 \int_{\mathcal T_{jk}}f(g)\,dg
 &=\int_{E_{jk}}\int_{S_{jk}(\eta(x))}f(x+v)\,dv\,d\Haus^{n-2}(x).
\end{align*}
The same formulas hold whenever \(f\) is integrable on the relevant tube.

(c) Use the function \(R\) in Fact~\ref{tb:projection-moments}, built
from the face normals of \(Q\), and choose \(\eta(x)\coloneqq 1/(4R(x))\).
Then all these tubes lie in \(\{g:g\in Q^{1/R(g)}\setminus Q\}\).
For an absolute constant \(c>0\), their normal integrals obey
\begin{align*}
 \int_0^{\eta(x)}R(x+s u_j)^3\varphi_n(x+s u_j)\,ds
 &\ge c\{1+a_j^2\}\varphi_n(x),\qquad x\in F_j,\\
 \int_{S_{jk}(\eta(x))}R(x+v)^3\varphi_n(x+v)\,dv
 &\ge c\alpha_{jk}\norm{z_{jk}}\varphi_n(x),\qquad x\in E_{jk}.
\end{align*}
For \(n=1\), there are no ridge tubes.
\end{fact}

\begin{proof}
(a) Facts~\ref{tb:normalcone} and~\ref{tb:projection} identify \(x\)
as the unique nearest point of \(x+v\). The relative interiors of
distinct faces are disjoint, so the tubes are disjoint.

(b) The map \((x,s)\mapsto x+s u_j\) is an affine isometry from
\(H_j\times\R\) onto \(\R^n\). Likewise, \((x,v)\mapsto x+v\)
is an affine isometry from \((z_{jk}+L_{jk}^\perp)\times L_{jk}\)
onto \(\R^n\). The domains are measurable in these coordinates because \(\eta\)
is measurable; their images under the affine isometries are therefore
measurable as well. The Jacobians are one. Apply Fubini to the indicator
of the domain \(0<s<\eta(x)\), or \(0<\norm v<\eta(x)\) in the
normal sector. No derivative of \(\eta\) is involved, and no
smoothness of \(\eta\) is needed.

(c) For a facet displacement \(v=s u_j\), and for a ridge displacement
\(v\in L_{jk}\), respectively,
\[
 |x\cdot v|\le\mathcal M(x)\norm v\le\tfrac14.
\]
The ridge case uses \(x\cdot v=z_{jk}\cdot v\). Since
\(\norm v\le1/4\), Fact~\ref{tb:gauss-coord}(b) gives
\(\varphi_n(x+v)\ge e^{-1/4-1/32}\varphi_n(x)\).
The Lipschitz property and \(R\ge1\) also give
\[
 \tfrac34R(x)\le R(x+v)\le\tfrac54R(x).
\]
The facet integral is therefore at least
\(cR(x)^3\eta(x)\varphi_n(x)=c'R(x)^2\varphi_n(x)\), which
dominates \(c''(1+a_j^2)\varphi_n(x)\) because
\(R(x)\ge1+|u_j\cdot x|=1+|a_j|\). The ridge integral is at least
\(c\alpha_{jk}R(x)^3\eta(x)^2\varphi_n(x)
=c'\alpha_{jk}R(x)\varphi_n(x)\); use
\(R(x)\ge\norm{z_{jk}}\).

Finally, a tube point \(g=x+v\) lies outside \(Q\), and each
unit-normal inequality defining \(Q\) satisfies
\[
 u_j\cdot g-a_j\le\norm v<\frac1{4R(x)}\le\frac1{R(g)}.
\]
This proves the inclusion in (c), even if the threshold system used
to define \(Q^s\) also contains redundant inequalities.
\end{proof}

Together with the paired ridge coefficients derived above, part (c)
gives
\[
 |\surf(Q)|\le C\E\left[R(G)^3
              \ind\{G\in Q^{1/R(G)}\setminus Q\}\right].
\]
The power three serves both dimensions: integration along a facet
uses one power of the tube width, leaving \(R^2\), while integration
over a ridge sector uses two, leaving \(R\).

\usedin{\PE{appx-an:adaptive-tubes}{C.29} in
\PR{Lemma}{appx-an:lem:surface}{C.5}. After setting
\(Q'\coloneqq (Q-V_0)/\sigma\), the event is
\(V_0+\sigma G\in Q^{\sigma/R(G)}\setminus Q\).
Conditional on \(G\), its width is fixed; concentration of \(V_0\)
and Fact~\ref{tb:projection-moments} give the averaged surface bound.}
\emph{Source.} The facet construction is Nazarov's outward-normal
argument~\TBc{Naz03}; see the detailed proof in \TBc{CCK17N},
Lemmas~1 and~2. Disjointness over arbitrary faces follows from the
classical projection characterization in \TBc{RW98}, Proposition~6.17.
The variable widths, ridge-sector integrals, and weighted estimates
needed here are derived above. For the related use of facet and ridge
tubes in Gaussian approximation, see \TBc{10}.

\subsection{The coarea formula and the first variation of a Gaussian
integral}\label{tb:sec:coarea}

The logistic mixture moves all thresholds together. To differentiate
its Gaussian integral, we integrate first over the level sets of the
maximum constraint. The coarea formula provides this change of variables.

\begin{fact}[Coarea formula]\label{tb:coarea}
Let \(\psi:\R^n\to\R\) be Lipschitz, so that \(\nabla\psi\) exists almost
everywhere. For every measurable \(f\ge0\),
\[
 \int_{\R^n}f\,\norm{\nabla\psi}\,dx=\int_{-\infty}^\infty
 \Bigl(\int_{\{\psi=s\}}f\,d\Haus^{n-1}\Bigr)ds ,
\]
For signed \(f\), the same identity holds if
\(\int |f|\norm{\nabla\psi}\,dx<\infty\), in particular if
\(f\in L^1(\R^n)\), since \(\psi\) is Lipschitz.
\end{fact}

\emph{Source.} \TBc{EG15}, Section~3.4. The proofs use the formula only for
a maximum of finitely many affine functions with gradients of equal length, and
for these Fact~\ref{tb:variation} gives a direct proof from Fubini's theorem.

For the finite maximum used here, the coarea formula can be proved
directly by partitioning the space according to the maximizing face.
This proof also identifies the exceptional levels and explains why the
independent Gamma shift permits differentiation at every fixed threshold.

\begin{fact}[First variation of a polyhedral sublevel integral]\label{tb:variation}
Let \(u_1,\ldots,u_m\) be unit vectors in \(\R^n\) (\(m\ge1\)), \(\kappa>0\),
\(b_1,\ldots,b_m\in\R\), and
\[
 \psi(g)\coloneqq \max_{j\le m}\{\kappa\,u_j\cdot g-b_j\},\qquad Q_s\coloneqq \{g:\psi(g)\le s\}.
\]
(a) For \(f\ge0\) measurable, or \(f\) integrable, and \(s\in\R\),
\[
 \int_{Q_s}f\,dg=\kappa^{-1}\int_{-\infty}^s\Bigl(\int_{\{\psi=s'\}}f\,d\Haus^{n-1}
 \Bigr)ds'.
\]
(b) If \(f\) is integrable, then \(\Lambda(s)\coloneqq \int_{Q_s}f\,dg\) is absolutely
continuous, and for almost every \(s\)
\[
 \Lambda'(s)=\kappa^{-1}\int_{\{\psi=s\}}f\,d\Haus^{n-1},\qquad
 |\Lambda'(s)|\le\kappa^{-1}\sum_{j=1}^m\int_{\{\kappa u_j\cdot g-b_j=s\}}|f|\,d\Haus^{n-1}.
\]
(c) If \(s>\inf\psi\), then \(Q_s\) is a proper, full-dimensional polyhedron
with face normals \(u_j\) and offsets \((s+b_j)/\kappa\), \(\{\psi=s\}=\partial Q_s\),
and \(\int_{\{\psi=s\}}f\,d\Haus^{n-1}\) is the surface integral of \(f\) over
\(\partial Q_s\) of Subsection~\ref{tb:sec:surface}. For \(s<\inf\psi\) the set
\(Q_s\) is empty, and \(s=\inf\psi\) is the only other level at which \(Q_s\)
may fail to be full-dimensional.

(d) For \(\Lambda\) from (b), assume in addition that it is
Lipschitz. If \(S\) is a real random variable with a density, then \(s\mapsto\E\Lambda(s+S)\) is differentiable with derivative
\(\E\Lambda'(s+S)\), where \(\Lambda'\) is defined almost everywhere.
\end{fact}

\begin{proof}
(a) Let \(C_j\) be the set where \(j\) is the smallest index attaining the
maximum defining \(\psi\); the Borel sets \(C_1,\ldots,C_m\) partition
\(\R^n\). On \(C_j\), \(\psi(g)=\kappa v-b_j\) in the coordinates
\(g=v u_j+y\), \(y\perp u_j\). By Fact~\ref{tb:fubini}, in these coordinates,
\[
 \int_{C_j\cap Q_s}f\,dg=\int_{u_j^\perp}\int_{-\infty}^{(s+b_j)/\kappa}
 (f\ind_{C_j})(v u_j+y)\,dv\,dy .
\]
The substitution \(s'=\kappa v-b_j\) turns the inner integral into
\(\kappa^{-1}\int_{-\infty}^s(f\ind_{C_j})(v u_j+y)\,ds'\) with
\(v=(s'+b_j)/\kappa\), and a further application of Fubini's theorem
exchanges the integrals. For fixed \(s'\), the set of \(v u_j+y\) in \(C_j\)
with this \(v\) is the slice \(C_j\cap\{\psi=s'\}\), and on the hyperplane
containing it \(\Haus^{n-1}\) is Lebesgue measure \(dy\)
(Fact~\ref{tb:hausdorff}(a)). Hence
\(\int_{C_j\cap Q_s}f\,dg=\kappa^{-1}\int_{-\infty}^s\int_{C_j\cap\{\psi=s'\}}
f\,d\Haus^{n-1}\,ds'\). Summing over \(j\) gives (a), since the slices
\(C_j\cap\{\psi=s'\}\) partition \(\{\psi=s'\}\).

(b) Applying (a) to \(|f|\) and letting \(s\to\infty\) shows that
\(I(s')\coloneqq \kappa^{-1}\int_{\{\psi=s'\}}f\,d\Haus^{n-1}\) is integrable on
\(\R\), and (a) says \(\Lambda(s)=\int_{-\infty}^sI(s')\,ds'\). Hence \(\Lambda\) is
absolutely continuous, and \(\Lambda'=I\) almost everywhere by Lebesgue's
differentiation theorem. The bound follows because
\(\{\psi=s\}\subseteq\bigcup_j\{\kappa u_j\cdot g-b_j=s\}\), a union of \(m\)
hyperplanes.

(c) The set \(Q_s\) is the polyhedron \(\bigcap_j\{u_j\cdot g\le(s+b_j)/\kappa\}\),
and \(\psi\) is convex and continuous. If \(\psi(g)=s>\inf\psi\), choose \(g'\)
with \(\psi(g')<s\); for \(\tau>0\), \(g=(g_\tau+\tau g')/(1+\tau)\) with
\(g_\tau\coloneqq g+\tau(g-g')\), so convexity gives
\(\psi(g_\tau)\ge s+\tau(s-\psi(g'))>s\). Thus \(g\) is a limit of points
outside \(Q_s\) and \(g\in\partial Q_s\); the open set \(\{\psi<s\}\) is
nonempty and lies in \(Q_s\), which is therefore full-dimensional, and it
is a proper subset because \(\psi(g)\ge\kappa u_1\cdot g-b_1\to\infty\) along
\(u_1\). A point with \(\psi<s\) is interior, so \(\partial Q_s=\{\psi=s\}\). After
the redundant and repeated inequalities are deleted, \(\partial Q_s=\{\psi=s\}\)
is the union of the sets \(Q_s\cap\{\kappa u_j\cdot g-b_j=s\}\) of the
remaining inequalities, which are faces of \(Q_s\) or empty. A face of
dimension \(n-1\) is a facet (Fact~\ref{tb:facets}(a)) and any other face is
\(\Haus^{n-1}\)-null, so the integral over \(\{\psi=s\}\) is the sum over the
facets of the irredundant description.

(d) The function \(\Lambda\) is bounded by \(\int|f|\), so the expectation
exists. A Lipschitz function is absolutely continuous and differentiable outside a
Lebesgue null set \(\mathcal N\). The random variable \(s+S\) avoids
\(\mathcal N\) almost surely, because \(S\) has a density, so the difference
quotients of \(\Lambda(s+S)\) converge almost surely and are bounded by the
Lipschitz constant. Dominated convergence gives the claim.
\end{proof}

\usedin{the first-variation paragraph in the proof of
\PR{Lemma}{appx-an:lem:matrix}{C.4}, immediately after
\PE{appx-an:gamma-mixture}{C.25}. Here
\(\psi(g)=\max_j\{u_j\cdot(V_0+\sigma g)-a_j\}\),
\(\kappa=\sigma\), \(b_j=a_j-u_j\cdot V_0\), and
\(f(g)=(\norm g^2-r)\varphi_r(g)\). Parts (a)--(c) give
\(\sigma^{-1}\surf((Q(a+s\mathbf1)-V_0)/\sigma)\) for almost every
\(s\). The independent shift \(-h\log T_\ell\), with
\(T_\ell\sim\operatorname{Gamma}(\ell,1)\), has a density, so part (d)
permits averaging these derivatives. A uniform bound
\(Cm(r+1)/\sigma\) follows from (b) and the Gaussian integral over
each supporting hyperplane; it justifies the interchange and does
not enter the final dimension-free bound.}

\clearpage
\subsection{Where each fact is used}\label{tb:sec:index}

The table follows the order of the proof. Each row identifies the
operation for which the corresponding facts are needed.

\medskip
\begin{center}
\small
\renewcommand{\arraystretch}{1.18}
\begin{tabular}{@{}p{0.27\textwidth}p{0.68\textwidth}@{}}
Facts~\ref{tb:fubini}, \ref{tb:taylor} &
Conditioning, Taylor remainders, and exchange of expectations with
convergent series in Appendices B--C.\\
Facts~\ref{tb:bernoulli}, \ref{tb:heat} &
Mask differentiation, the heat equation, and the signed Hessian
comparison in Subsection~\PN{appx-an:sec:path}{C.2}.\\
Facts~\ref{tb:faa}, \ref{tb:logistic}, \ref{tb:covariance} &
Derivative envelopes in Appendix~\PN{appx-sec:smoothing}{B}; early
bounds in \PN{appx-an:sec:early}{C.3} and coefficient matrices in
\PR{Lemma}{appx-an:lem:restoration}{C.2}.\\
Facts~\ref{tb:stein}, \ref{tb:nuclear} &
Two Gaussian integrations by parts and matrix contraction in
\PR{Lemma}{appx-an:lem:restoration}{C.2}.\\
Facts~\ref{tb:gauss-coord}, \ref{tb:logsum} &
Reference-face mean bounds and summation of reference masses in
\PR{Lemmas}{appx-an:lem:mean}{C.3} and~\PN{appx-an:lem:matrix}{C.4}.\\
Facts~\ref{tb:prekopa}--\ref{tb:lc-weights} &
Logistic-tilt covariance contraction in
\PR{Lemma}{appx-an:lem:matrix}{C.4}.\\
Facts~\ref{tb:facets}--\ref{tb:divergence} &
Surface measure and the facet--ridge identity in
\PR{Lemma}{appx-an:lem:surface}{C.5}.\\
Facts~\ref{tb:tubes}, \ref{tb:projection-moments} &
Adaptive normal tubes and Gaussian projection moments in
\PR{Lemma}{appx-an:lem:surface}{C.5}.\\
Facts~\ref{tb:coarea}, \ref{tb:variation} &
First variation of the Gamma mixture in
\PR{Lemma}{appx-an:lem:matrix}{C.4}.\\
\end{tabular}
\end{center}

\subsection{Sources}\label{tb:sec:sources}

\begin{description}\footnotesize\setlength{\itemsep}{0pt}\setlength{\parskip}{0pt}
\item[{[BL76]}] H.~J. Brascamp and E.~H. Lieb (1976). On extensions of the
Brunn--Minkowski and Pr\'ekopa--Leindler theorems, including inequalities
for log concave functions, and with an application to the diffusion
equation. \emph{Journal of Functional Analysis} \textbf{22}, 366--389.
\item[{[BV04]}] S. Boyd and L. Vandenberghe (2004).
\emph{Convex Optimization}. Cambridge University Press, Cambridge.
\item[{[CCK17N]}] V. Chernozhukov, D. Chetverikov, and K. Kato (2017).
Detailed proof of Nazarov's inequality. arXiv:1711.10696.
\item[{[Com74]}] L. Comtet (1974). \emph{Advanced Combinatorics}. Reidel,
Dordrecht.
\item[{[CT06]}] T.~M. Cover and J.~A. Thomas (2006). \emph{Elements of
Information Theory}, 2nd edition. Wiley.
\item[{[EG15]}] L.~C. Evans and R.~F. Gariepy (2015). \emph{Measure Theory
and Fine Properties of Functions}, revised edition. CRC Press.
\item[{[Fol99]}] G.~B. Folland (1999). \emph{Real Analysis: Modern
Techniques and Their Applications}, 2nd edition. Wiley.
\item[{[Gar02]}] R.~J. Gardner (2002). The Brunn--Minkowski inequality.
\emph{Bulletin of the American Mathematical Society} \textbf{39}, 355--405.
\item[{[10]}] Z. Gao and W. Zhang (2026). Cubic-root Gaussian approximation
under unrestricted covariance. arXiv:2608.30221v1.
\item[{[Lei72]}] L. Leindler (1972). On a certain converse of H\"older's
inequality II. \emph{Acta Scientiarum Mathematicarum (Szeged)} \textbf{33},
217--223.
\item[{[Naz03]}] F. Nazarov (2003). On the maximal perimeter of a convex
set in \(\R^n\) with respect to a Gaussian measure. In V.~D. Milman
and G. Schechtman (eds.), \emph{Geometric Aspects of Functional Analysis},
Lecture Notes in Mathematics \textbf{1807}, 169--187. Springer.
\item[{[Pre71]}] A. Pr\'ekopa (1971). Logarithmic concave measures with
application to stochastic programming. \emph{Acta Scientiarum Mathematicarum
(Szeged)} \textbf{32}, 301--316.
\item[{[Pre73]}] A. Pr\'ekopa (1973). On logarithmic concave measures and
functions. \emph{Acta Scientiarum Mathematicarum (Szeged)} \textbf{34},
335--343.
\item[{[Roc70]}] R.~T. Rockafellar (1970). \emph{Convex Analysis}. Princeton
University Press.
\item[{[RW98]}] R.~T. Rockafellar and R.~J.-B. Wets (1998). \emph{Variational
Analysis}. Springer, Berlin.
\item[{[Ste81]}] C.~M. Stein (1981). Estimation of the mean of a multivariate
normal distribution. \emph{The Annals of Statistics} \textbf{9}, 1135--1151.
\item[{[Zie95]}] G.~M. Ziegler (1995). \emph{Lectures on Polytopes}.
Springer, New York.
\end{description}



% Fallback paper numbers below are checked against the v3 auxiliary file.

\makeatletter
\providecommand{\E}{\mathbb E}
\providecommand{\Pp}{\mathbb P}
\providecommand{\R}{\mathbb R}
\providecommand{\tr}{\operatorname{tr}}
\providecommand{\op}{\mathrm{op}}
\providecommand{\norm}[1]{\left\lVert#1\right\rVert}
\providecommand{\ind}{\mathbf 1}
\providecommand{\dd}{\,d}
\providecommand{\Hh}{H_2}
\providecommand{\surf}{\mathcal J}
\providecommand{\relint}{\operatorname{relint}}
\providecommand{\aff}{\operatorname{aff}}
\providecommand{\cone}{\operatorname{cone}}
\providecommand{\dist}{\operatorname{dist}}
\providecommand{\prj}{\operatorname{pr}}
\providecommand{\Haus}{\mathcal H}
\providecommand{\TBc}[1]{[#1]}
% \PN{label}{C.5} prints a number of the paper (a \ref in the paper);
% \PR{Lemma}{label}{C.5} prints "Lemma C.5"; \PE{label}{C.21} an equation.
\providecommand{\PN}[2]{\ifdefined\toolboxinpaper\ref{#1}\else#2\fi}
\providecommand{\PR}[3]{#1~\PN{#2}{#3}}
\providecommand{\PE}[2]{(\PN{#1}{#2})}
\providecommand{\usedin}[1]{\par\smallskip\noindent{\small\emph{Used in:} #1}\par\smallskip}
\@ifundefined{c@fact}{\newtheorem{fact}{Fact}[section]}{}
\makeatother

\section{Toolbox: facts from analysis and geometry}\label{app:toolbox}

The late interpolation estimate reduces high derivatives to two objects:
a Gaussian matrix weighted by the logistic approximation, and a signed
Gaussian integral over a polyhedral boundary. This appendix supplies the
analysis and geometry behind those reductions. The calculus facts also
justify the early estimate and the smoothing argument in
Appendices~\PN{appx-sec:smoothing}{B} and~\PN{appx-sec:support}{C}.
Each fact is stated with the assumptions needed for its use, followed by
a proof or a standard reference and a precise pointer to the manuscript.

We first treat differentiation, Gaussian integration by parts, and
covariance contraction. Log-concavity then controls the Gaussian law
tilted by a logistic weight. The geometric facts explain how to pair the
two contributions at each ridge and compare the resulting boundary
integral with the mass of short outward tubes. Because the faces are
flat, the required area and tube formulas follow from orthogonal
coordinates and Fubini's theorem. Finally, the first-variation formula
connects the logistic mixture to the surface estimate, including the
exceptional levels at which the polyhedron loses dimension.
Subsection~\ref{tb:sec:index} records these dependencies in one table.

\subsection{Notation}\label{tb:sec:notation}

The smoothing and interpolation arguments take place in the covariance
support \(\R^r\). They use a fixed list of \(m=2p\) unit face normals
\(u_1,\ldots,u_m\), which may repeat, and the threshold sets
\[
 A_a\coloneqq\{x\in\R^r:u_\alpha\cdot x\le a_\alpha,
                          \ \alpha=1,\ldots,m\},\qquad a\in\R^m.
\]
All thresholds are finite; inconsistent inequalities give the empty set.
We keep the full list when defining the logistic approximation, even if
some inequalities are redundant. We delete redundant inequalities only
when describing the facets of a fixed full-dimensional polyhedron.
For a specified system of unit-normal inequalities, \(Q^s\) increases
each threshold by \(s\); this notation depends on the system and does
not denote a Euclidean neighborhood. We write \(L\coloneqq\log(ep)\).

In the general facts below, \(\R^n\) carries the inner product
\(x\cdot y\) and Euclidean norm \(\norm x\). This dimension is \(r\)
in the support, or \(r-1\) when we work within a facet. We write \(dx\) for Lebesgue measure, \(\varphi_n\) and
\(\Phi_n\) for the standard Gaussian density and law on \(\R^n\), and \(G\)
for a standard Gaussian vector. For a linear subspace \(L\), \(\varphi_L\)
is the standard Gaussian density on \(L\) with respect to Lebesgue measure
on \(L\). In dimension zero, Lebesgue and Gaussian measure both assign
mass one to the origin, and \(\varphi_0(0)\coloneqq1\). This convention
includes the tangential calculation at an endpoint when \(r=1\).
For a nonempty set \(S\), \(\ind_S\) is its indicator function, \(\aff S\)
is the smallest affine subspace containing it, the dimension of \(S\) is
that of \(\aff S\), and \(\relint S\) is the interior of \(S\) relative to
\(\aff S\). The distance from a point \(y\) to a set \(S\) is
\(\dist(y,S)\coloneqq \inf_{x\in S}\norm{y-x}\).

For an integer \(q\ge1\) and \(f\in C^q(\R^n)\), the derivative of order \(q\) in the directions
\(z_1,\ldots,z_q\in\R^n\) is
\[
 D^qf(x)[z_1,\ldots,z_q]
 \coloneqq \partial_{s_1}\cdots\partial_{s_q}f(x+s_1z_1+\cdots+s_qz_q)\Big|_{s=0}
 =\sum_{i_1,\ldots,i_q}\partial_{i_1\cdots i_q}f(x)\,
   z_{1,i_1}\cdots z_{q,i_q},
\]
which is symmetric and multilinear in the directions. We write
\(D^qf(x)[z^{\otimes q}]\) when all directions equal \(z\), with
\(D^0f\coloneqq f\). The first derivative satisfies
\(Df(x)[z]=\nabla f(x)\cdot z\), and we identify the second derivative
with its Hessian matrix. The Loewner order \(M\preceq N\) means that
\(N-M\) is positive semidefinite, \(\norm\Xi_\op\) is the operator norm
of a matrix \(\Xi\), and \(\norm M_*\) is the \emph{nuclear norm}, which for
a symmetric matrix \(M\) with eigenvalues \(\mu_i\) equals \(\sum_i|\mu_i|\).

\subsection{Integration and calculus}\label{tb:sec:calc}

We repeatedly condition on an independent component and then average the
resulting bound. The following theorem specifies when these operations,
and the corresponding exchanges of integrals or sums, are valid.

\begin{fact}[Fubini and Tonelli]\label{tb:fubini}
Let \((S_1,\mu_1)\) and \((S_2,\mu_2)\) be \(\sigma\)-finite measure
spaces and let \(f\) be a measurable function on \(S_1\times S_2\). If
\(f\ge0\), then \(s_1\mapsto\int f(s_1,s_2)\,\mu_2(ds_2)\) is measurable
and
\[
 \int\Bigl(\int f\,d\mu_2\Bigr)d\mu_1
 =\int\Bigl(\int f\,d\mu_1\Bigr)d\mu_2
 =\int f\,d(\mu_1\otimes\mu_2)\in[0,\infty].
\]
If \(\int|f|\,d(\mu_1\otimes\mu_2)<\infty\), the same equalities hold, and
the inner integrals are finite almost everywhere.
\end{fact}

Three special cases occur in the proofs. If \(X\) and \(Y\) are
independent random elements and either \(f\ge0\) or \(f(X,Y)\) is
integrable, then
\[
 \E f(X,Y)=\E g(X),\qquad g(x)\coloneqq \E f(x,Y),
\]
because the joint law is the product of the marginal laws. This is the
formal content of every argument that is carried out for a fixed value of
an independent component such as \(V_0\) and then averaged over it. With
\(\mu_2\) the counting measure, the theorem says that
\(\E\sum_k\xi_k=\sum_k\E\xi_k\) for nonnegative terms, and for terms with
\(\sum_k\E|\xi_k|<\infty\). For an orthogonal decomposition
\(\R^n=L\oplus L^\perp\), Lebesgue measure on \(\R^n\) is the product of the
Lebesgue measures on \(L\) and on \(L^\perp\), so that
\(\int_{\R^n}f\,dx=\int_{L^\perp}\int_Lf(l+w)\,dl\,dw\).

The final averaging in \PR{Lemma}{appx-an:lem:surface}{C.5} is an
instance of the first case. Let \(V_0\in\R^r\) and
\(G\sim N(0,I_r)\) be independent, and let \(\sigma>0\).
For a fixed threshold set \(Q\), it gives
\[
 \E\!\left[R(G)^3\ind\{V_0+\sigma G\in
                   Q^{\sigma/R(G)}\setminus Q\}\right]
 =\E_G\!\left[R(G)^3\Pp_{V_0}\{V_0+\sigma G\in
                   Q^{\sigma/R(G)}\setminus Q\}\right].
\]
Here \(R\) is the projection radius in
Fact~\ref{tb:projection-moments}, and \(Q^s\) shifts every
unit-normal threshold outward by \(s\). Conditional on \(G\), both
the width \(\sigma/R(G)\) and the translation \(\sigma G\) are fixed.
Concentration of \(V_0\) can therefore be applied directly. Tonelli
requires no integrability check at this stage because the integrand
is nonnegative.

\usedin{conditioning on independent components in Appendices~\PN{appx-sec:smoothing}{B}
and~\PN{appx-sec:support}{C}; absolute convergence of the translated-weight
series in \PR{Lemma}{appx-an:lem:restoration}{C.2}; final averaging in
\PR{Lemma}{appx-an:lem:surface}{C.5}; conditioning on independent pairs
in the retained proof in \PR{Appendix}{appx-app:concentration}{D}.}
\emph{Source.} \TBc{Fol99}, Section~2.5.

To differentiate the interpolation, we first justify moving a derivative
inside an expectation and then expand along a line segment. An integrable
bound is needed for the former operation; the latter is ordinary
one-dimensional Taylor expansion applied to a multivariate function.

\begin{fact}[Differentiation under the integral sign; Taylor's formula]
\label{tb:taylor}
(a) Let \(I\) be an open interval and let \(f(t,\omega)\) be
measurable and integrable in \(\omega\) for each \(t\in I\).
Suppose that, outside one null set, it is differentiable in \(t\).
On each compact subinterval \(J\subset I\), assume
\(|\partial_t f(t,\omega)|\le g_J(\omega)\) for all \(t\in J\),
where \(g_J\) is integrable. Then \(\frac d{dt}\E f(t,\cdot)
=\E\,\partial_t f(t,\cdot)\) on \(I\).

(b) For an integer \(q\ge1\), \(f\in C^q(\R^n)\), and \(x,v\in\R^n\),
\[
 f(x+v)=\sum_{k=0}^{q-1}\frac1{k!}D^kf(x)[v^{\otimes k}]
 +\int_0^1\frac{(1-s)^{q-1}}{(q-1)!}D^qf(x+sv)[v^{\otimes q}]\,ds.
\]
For \(q=1\) this is \(f(x+v)-f(x)=\int_0^1Df(x+sv)[v]\,ds\).
\end{fact}

\begin{proof}
For (a), fix an interior time and a compact neighborhood \(J\).
The mean value theorem bounds nearby difference quotients by \(g_J\),
and they converge pointwise, so dominated convergence applies. For (b),
the function \(s\mapsto f(x+sv)\) is of class \(C^q\) with \(k\)th derivative
\(D^kf(x+sv)[v^{\otimes k}]\); apply the one-dimensional Taylor formula with
integral remainder at \(s=0\) and evaluate it at \(s=1\).
\end{proof}

\usedin{the signed Hessian comparison and its Taylor expansion in
Subsection~\PN{appx-an:sec:path}{C.2}; the restoration of the cubic term
to the full path in \PR{Lemma}{appx-an:lem:identity}{C.1}; the segment
integral in \PR{Lemma}{appx-lem:endpoint-gap}{B.3}.}
\emph{Source.} \TBc{Fol99}, Section~2.3.

Differentiating the thinning probability changes one mask at a time.
Here a mask is a variable that retains a summand when it equals one
and deletes it when it equals zero. The following finite-product identity
isolates the change in retention probability; changes in amplitude and
Gaussian variance are separate contributions.

\begin{fact}[Differentiating independent Bernoulli masks]\label{tb:bernoulli}
Let \(B_1,\ldots,B_N\) be independent Bernoulli variables with parameter
\(q\in(0,1)\). For a function \(H:\{0,1\}^N\to\R\), write
\(B^{i,b}\) for the vector obtained by setting its \(i\)th entry to
\(b\in\{0,1\}\). Then
\[
 \frac{d}{dq}\E_q H(B)
 =\sum_{i=1}^N\E_q\{H(B^{i,1})-H(B^{i,0})\}.
\]
If \(q=q(t)\) and \(H=H_t\) are differentiable, then
\[
 \frac{d}{dt}\E_{q(t)}H_t(B)
 =\E_{q(t)}\partial_t H_t(B)
 +q'(t)\sum_i\E_{q(t)}
       \{H_t(B^{i,1})-H_t(B^{i,0})\}.
\]
\end{fact}

\begin{proof}
Temporarily give the masks separate parameters \(q_1,\ldots,q_N\).
Conditioning on all masks except the \(i\)th makes their expectation
affine in \(q_i\), with derivative equal to the insertion--deletion
difference in the first display. Differentiating along
\(q_1=\cdots=q_N=q\) gives the sum. The second formula is the chain
rule. Independent random arguments may be included by conditioning
and then applying Fact~\ref{tb:taylor}(a).
\end{proof}

\usedin{Subsection~\PN{appx-an:sec:path}{C.2}, before
\PR{Lemma}{appx-an:lem:identity}{C.1}, to differentiate the masks in
the Bernoulli--Gauss path.}

In the paper, the amplitude and retention probability are
\(\gamma=1+t\) and \(\varpi=(1-t)/(1+t)^2\), for \(0<t<1\).
The support-coordinate summands \(\zeta_i^0\) are centered, with
\(\sum_i\E\zeta_i^0(\zeta_i^0)^\top=I_r\). Write
\(Y_i\coloneqq\gamma\mathbb B_{i,t}\zeta_i^0\) and
\(U_t\coloneqq\sum_iY_i+\sqrt t\,G\), where all masks and \(G\)
are mutually independent and independent of the summands. To expand the
contribution of summand \(i\), remove it and put \(H_i\coloneqq U_t-Y_i\).
Let \(\zeta_i\) be an independent copy of \(\zeta_i^0\), with all copies
independent of each other and of the path variables. The signed Hessian
comparison in~\PE{appx-an:signed-hessian}{C.6} has scalar weight
\[
 \nu_t(s)\coloneqq \frac{4-2(3-t)s}{1+t},\qquad 0\le s\le1.
\]
Direct integration gives
\[
 \int_0^1\nu_t(s)\,ds=1,\qquad
 \int_0^1s\nu_t(s)\,ds=\frac{2t}{3(1+t)}.
\]
For a twice continuously differentiable scalar function \(g\),
Taylor's formula and the first identity give
\[
 g(0)-\int_0^1\nu_t(s)g(s)\,ds
 =-\frac{2t}{3(1+t)}g'(0)
   -\int_0^1\nu_t(s)\int_0^s(s-u)g''(u)\,du\,ds.
\]
For the logistic approximation \(F=F_{a,h}\) defined below, take
\(g(s)\coloneqq D^2F(H_i+s\gamma\zeta_i)[\zeta_i,\zeta_i]\).
Then \(g'(0)=\gamma D^3F(H_i)[\zeta_i^{\otimes3}]\), so
the factor \(1/2\) in the signed Hessian comparison turns the first
term into
\(-\tfrac12\{2t/(3\gamma)\}\gamma
 D^3F(H_i)[\zeta_i^{\otimes3}]
 =-tD^3F(H_i)[\zeta_i^{\otimes3}]/3\). This explains the
cubic coefficient before the derivatives are restored to the full path.
The weight \(\nu_t\) is signed; these identities use its
integrals, not a probability interpretation.

Repeated differentiation of the logistic approximation produces products
of face probabilities. The partition form of the higher-order chain
rule records which directions differentiate the same inner factor.

\begin{fact}[Fa\`a di Bruno's formula]\label{tb:faa}
Let \(\phi\) be a smooth function on an open interval of \(\R\) and let
\(g:\R^n\to\R\) be smooth with values in that interval. For \(q\ge1\), let
\(\Pi_q\) be the set of partitions \(\mathcal P\) of \(\{1,\ldots,q\}\) into
nonempty blocks, let \(|\mathcal P|\) be the number of blocks, and for a
block \(J\) let \(z_J\) be the list of directions \(z_a\) with
\(a\in J\). Then
\[
 D^q(\phi\circ g)(x)[z_1,\ldots,z_q]
 =\sum_{\mathcal P\in\Pi_q}\phi^{(|\mathcal P|)}(g(x))
  \prod_{J\in\mathcal P}D^{|J|}g(x)[z_J].
\]
\end{fact}

\begin{proof}
For \(q=1\) this is the chain rule. Differentiate the right side for \(q\)
in the direction \(z_{q+1}\). By the product rule, the derivative either
falls on \(\phi^{(|\mathcal P|)}(g(x))\), which gives
\(\phi^{(|\mathcal P|+1)}(g(x))\,Dg(x)[z_{q+1}]\) and adds the block
\(\{q+1\}\) to \(\mathcal P\), or falls on one factor
\(D^{|J|}g(x)[z_J]\), which adds \(q+1\) to the block \(J\). Every partition
of \(\{1,\ldots,q+1\}\) arises exactly once in this way, by removing
\(q+1\) from its block.
\end{proof}

\emph{Source.} \TBc{Com74}; the induction above proves the form used here.

For a finite list of face normals \(u_1,\ldots,u_m\in\R^n\), thresholds
\(a_1,\ldots,a_m\in\R\), and a width \(h>0\), put
\[
 q_\alpha(x)\coloneqq(u_\alpha\cdot x-a_\alpha)/h,\qquad
 w_\alpha(x)\coloneqq e^{q_\alpha(x)},\qquad T(x)\coloneqq\sum_\alpha w_\alpha(x),
\]
\[
 F(x)\coloneqq\frac1{1+T(x)},\qquad
 \pi_\alpha(x)\coloneqq F(x)w_\alpha(x),\qquad
 \theta(x)\coloneqq1-F(x)=\sum_\alpha\pi_\alpha(x).
\]
Here \(F\) is the logistic approximation to the indicator of the
threshold set. The baseline outcome has weight one and probability
\(F\); face \(\alpha\) has probability \(\pi_\alpha\). These
probabilities describe the smoothing formula, not the original event.
One and two differentiations give
\begin{align*}
 DF(x)[z]&=-h^{-1}\sum_\alpha F\pi_\alpha(u_\alpha\cdot z),\\
 D^2F(x)[z_1,z_2]
 &=h^{-2}\left\{2\sum_{\alpha,\beta}F\pi_\alpha\pi_\beta
       (u_\alpha\cdot z_1)(u_\beta\cdot z_2)
       -\sum_\alpha F\pi_\alpha
       (u_\alpha\cdot z_1)(u_\alpha\cdot z_2)\right\}.
\end{align*}
All weights are evaluated at \(x\). In the first formula their sum is
\(F\theta\), which is small away from the smoothed boundary.
The next result retains this factor at every derivative order.

\begin{fact}[Derivatives of the logistic approximation]\label{tb:logistic}
With the notation above, let \(q\ge1\).
For directions \(z_1,\ldots,z_q\), with the blocks \(J_1,\ldots,J_{|\mathcal P|}\)
of \(\mathcal P\) listed in any fixed order,
\[
 D^qF(x)[z_1,\ldots,z_q]
 =h^{-q}\sum_{\mathcal P\in\Pi_q}(-1)^{|\mathcal P|}|\mathcal P|!
 \sum_{\alpha_1,\ldots,\alpha_{|\mathcal P|}}
 F(x)\prod_{b=1}^{|\mathcal P|}\pi_{\alpha_b}(x)
 \prod_{b=1}^{|\mathcal P|}\prod_{a\in J_b}(u_{\alpha_b}\cdot z_a),
\]
Moreover, \(\sum_{\mathcal P\in\Pi_q}|\mathcal P|!\le q!\,2^{q-1}\), and hence
\[
 |D^qF(x)[z_1,\ldots,z_q]|
 \le q!\,2^{q-1}h^{-q}F(x)\theta(x)
       \prod_{a=1}^q\max_\alpha|u_\alpha\cdot z_a|.
\]
\end{fact}

\begin{proof}
Apply Fact~\ref{tb:faa} with \(\phi(\tau)\coloneqq (1+\tau)^{-1}\) and \(g=T\). Then
\(\phi^{(k)}(\tau)=(-1)^kk!\,(1+\tau)^{-k-1}\), and since each \(w_\alpha\) is
an exponential of an affine function,
\(D^lT(x)[z_1,\ldots,z_l]=h^{-l}\sum_\alpha w_\alpha(x)\prod_{a\le l}
(u_\alpha\cdot z_a)\). Substituting, and using
\((1+T)^{-|\mathcal P|-1}\prod_bw_{\alpha_b}=F\prod_b\pi_{\alpha_b}\) and
\(\sum_b|J_b|=q\), gives the formula. The number
\(\sum_{\mathcal P}|\mathcal P|!\) counts ordered partitions
\((J_1,\ldots,J_k)\) of \(\{1,\ldots,q\}\). Writing the elements of
\(J_1\), then of \(J_2\), and so on, in increasing order inside each block,
gives a permutation together with the set of positions at which a new block
starts, which is a subset of the \(q-1\) gaps between consecutive entries.
Every ordered partition arises from at least one of these
\(q!\,2^{q-1}\) pairs. For a partition with \(k\ge1\) blocks, summing the probability
weights gives \(F\theta^k\le F\theta\). Each product of projections
is bounded by \(\prod_a\max_\alpha|u_\alpha\cdot z_a|\).
These two observations prove the stated envelope without a factor
depending on the dimension.
\end{proof}

The early estimate uses this formula for orders three and four. After
the two Gaussian integrations by parts in the late estimate, the explicit
first and second derivatives above suffice. In both cases the factor
\(F\theta\) connects the derivative estimate to concentration near
the boundary.

\usedin{the general derivative expansion and local envelope in
Subsection~\PN{appx-sec:mixture}{B.2}, namely
\PE{appx-eq:partition-derivative}{B.1}--\PE{appx-eq:local-envelope}{B.3};
the explicit first and second derivatives in
\PR{Lemma}{appx-an:lem:restoration}{C.2}.}
\emph{Source.} The formula follows from the partition chain rule in
Fact~\ref{tb:faa}; the proof above is self-contained.

\subsection{Gaussian identities}\label{tb:sec:gauss}

The mean estimate separates a normal coordinate from the tangential
coordinates of a face. Orthogonality makes these Gaussian components
independent; a second elementary identity controls the density change
when we move outward from the face.

\begin{fact}[Gaussian density in orthogonal coordinates]\label{tb:gauss-coord}
(a) Let \(\R^n=L\oplus L^\perp\) be an orthogonal decomposition. Then
\(\varphi_n(l+w)=\varphi_L(l)\varphi_{L^\perp}(w)\) for \(l\in L\) and
\(w\in L^\perp\), so the projections of a standard Gaussian vector onto
\(L\) and \(L^\perp\) are independent standard Gaussian vectors. In
particular, for a unit vector \(u\), \(a\in\R\), and \(y\perp u\),
\(\varphi_n(au+y)=\varphi_1(a)\varphi_{n-1}(y)\).

(b) For \(x,v\in\R^n\),
\(\varphi_n(x+v)=\varphi_n(x)\exp\{-x\cdot v-\tfrac12\norm v^2\}\).
\end{fact}

\begin{proof}
Both follow from \(\norm{l+w}^2=\norm l^2+\norm w^2\) and from expanding
\(\norm{x+v}^2\).
\end{proof}

\usedin{the normal--tangential decomposition in
\PR{Lemma}{appx-an:lem:mean}{C.3}; the facet density factorization
and density comparison along normal tubes in
\PR{Lemma}{appx-an:lem:surface}{C.5}.}

Gaussian integration by parts replaces two directional derivatives by
a centered quadratic function of the Gaussian vector. The factor
\(\sigma^{-2}\) below records the Gaussian noise scale exactly.

\begin{fact}[Gaussian integration by parts]\label{tb:stein}
Let \(G\sim N(0,I_n)\). In (a), assume \(f\in C^1(\R^n)\) and
that \(f\) and its first derivatives have polynomial growth.
In (b), assume in addition that \(f\in C^2\) and its second
derivatives have polynomial growth. Then
\begin{enumerate}
\item[(a)] \(\E[Gf(G)]=\E\nabla f(G)\);
\item[(b)] \(\E[(GG^\top-I_n)f(G)]=\E D^2f(G)\); equivalently, for \(z\in\R^n\),
\(\E D^2f(G)[z,z]=\E\bigl[\tr\{(GG^\top-I_n)zz^\top\}f(G)\bigr]\), where
\(\tr\{(GG^\top-I_n)zz^\top\}=(z\cdot G)^2-\norm z^2\);
\item[(c)] if \(s\ge0\) is an integer, \(v\in\R^n\), \(\sigma>0\),
\(U\coloneqq v+\sigma G\), and
\(F\in C^{s+2}(\R^n)\) has derivatives of polynomial growth, then for
\(z,z_1,\ldots,z_s\in\R^n\),
\[
 \E D^{s+2}F(U)[z,z,z_1,\ldots,z_s]
 =\sigma^{-2}\E\bigl[\tr\{(GG^\top-I_n)zz^\top\}\,
 D^sF(U)[z_1,\ldots,z_s]\bigr].
\]
\end{enumerate}
\end{fact}

\begin{proof}
For (a) in one coordinate, \(x\varphi_1(x)=-\varphi_1'(x)\), so integration by
parts gives
\[\int xf(x)\varphi_1(x)\,dx=\int f'(x)\varphi_1(x)\,dx,\]
and the boundary terms vanish because of the growth condition. The other coordinates
are handled by Fact~\ref{tb:fubini}. For (b), apply (a) to \(\partial_if\) and
to \(G_kf\): \(\E[G_iG_kf(G)]=\E\partial_i(G_kf)(G)=\delta_{ik}\E f(G)
+\E[G_k\partial_if(G)]=\delta_{ik}\E f(G)+\E\partial_{ik}f(G)\). For (c), the
function \(f(g)\coloneqq D^sF(v+\sigma g)[z_1,\ldots,z_s]\) has
\(D^2f(g)[z,z]=\sigma^2D^{s+2}F(v+\sigma g)[z,z,z_1,\ldots,z_s]\); apply (b).
\end{proof}

\usedin{Subsection~\PN{sec:outline-late}{2.4} and
\PR{Lemma}{appx-an:lem:restoration}{C.2}, conditional on \(\mathcal H\)
with \(v=V_0+v_i\) and \(z=Z_i\); the identity
\(\E[Z\cdot\nabla\Psi(Z)]=\E\tr(\Sigma D^2\Psi(Z))\)
in \PR{Lemma}{appx-lem:endpoint-gap}{B.3}.}
The identity also holds conditionally when the base and directions are
measurable with respect to a sigma-field independent of \(G\), provided
the conditional growth assumptions hold. Averaging that identity further
requires integrability of both sides. In the paper the fixed width
\(h>0\), bounded face components, and conditional covariance bound
supply this integrability.

\emph{Source.} Stein~\TBc{Ste81}; the proof above is complete.

Changing the Gaussian variance contributes a Hessian trace to the path
derivative. We give the formula for a general covariance and then use
its standard-Gaussian form in the support coordinates.

\begin{fact}[Heat equation for the Gaussian component]\label{tb:heat}
Let \(Z\sim N(0,\Sigma)\) in \(\R^p\), where \(\Sigma\) may be singular, and
let \(f\in C^2(\R^p)\) be bounded with bounded first and second derivatives.
Then, for \(x\in\R^p\) and \(t>0\),
\[
 \frac d{dt}\E f(x+\sqrt t Z)
 =\tfrac12\E\tr\{\Sigma D^2f(x+\sqrt t Z)\}.
\]
\end{fact}

\begin{proof}
Write \(Z=\Gamma G\) with \(\Gamma\Gamma^\top=\Sigma\) and \(G\) standard
Gaussian. By Fact~\ref{tb:taylor}(a), which applies because the difference
quotients near \(t_0>0\) are bounded by
\(\norm\nabla f_\infty\norm{\Gamma G}/(2\sqrt{t_0/2})\)
when the times lie in \([t_0/2,3t_0/2]\), the derivative equals
\((2\sqrt t)^{-1}\E[G\cdot\Gamma^\top\nabla f(x+\sqrt t\Gamma G)]\).
Fact~\ref{tb:stein}(a) applied to the function \(g\mapsto
\Gamma^\top\nabla f(x+\sqrt t\Gamma g)\) turns this into
\((2\sqrt t)^{-1}\sqrt t\,\E\tr\{\Gamma^\top D^2f\,\Gamma\}
=\frac12\E\tr\{\Sigma D^2f\}\).
\end{proof}

In the support \(\R^r\), the covariance is \(I_r\), so this gives
\(\frac d{dt}\E_G F(v+\sqrt tG)=\frac12\E_G\tr D^2F(v+\sqrt tG)\).
The base \(v\) is held fixed for this differentiation. If its law also
depends on time, its change contributes separately.

\usedin{Subsection~\PN{sec:outline-path}{2.2} and the derivation of
\PE{appx-an:signed-hessian}{C.6} in
Subsection~\PN{appx-an:sec:path}{C.2}. The Gaussian component supplies
the covariance term that cancels the constant term of the Hessian
expansion.}

The surface estimate needs the size of a few normal projections,
not the Euclidean length of the full Gaussian vector. The following
bound explains why only a logarithm of the number of faces enters.

\begin{fact}[Maxima of one- and two-dimensional Gaussian projections]
\label{tb:projection-moments}
Let \(u_1,\ldots,u_m\) be unit vectors in \(\R^r\), with \(m\ge1\).
For linearly independent \(u_j,u_k\), let \(\Pi_{jk}\) be orthogonal
projection onto their span, and put
\[
 \mathcal M(g)\coloneqq \max\left\{\max_j|u_j\cdot g|,
   \max_{j,k:\,u_j,u_k\text{ independent}}\norm{\Pi_{jk}g}\right\},
 \qquad R(g)\coloneqq 1+\mathcal M(g).
\]
An empty second maximum is omitted. Then \(R\) is \(1\)-Lipschitz.
For \(G\sim N(0,I_r)\) and \(q\in\{2,3\}\),
\[
 \E R(G)^q\le C_q\{\log(em)\}^{q/2}.
\]
\end{fact}

\begin{proof}
Every projection is a contraction, so each function in the maximum
is \(1\)-Lipschitz, as is their maximum. A scalar standard normal
satisfies \(\Pp(|Z|>s)\le2e^{-s^2/2}\); the length of a
two-dimensional standard normal satisfies
\(\Pp(\norm{G_2}>s)=e^{-s^2/2}\). A union bound gives
\[
 \Pp\{\mathcal M(G)>s\}\le\min\{1,2(m+m^2)e^{-s^2/2}\}.
\]
The projections need not be independent. Integrate
\(q\int_0^\infty s^{q-1}\Pp\{\mathcal M(G)>s\}\,ds\), splitting
at \(\sqrt{8\log(em)}\), and use
\((1+x)^q\le2^{q-1}(1+x^q)\).
\end{proof}

\usedin{\PE{appx-an:projection-moments}{C.28} in
\PR{Lemma}{appx-an:lem:surface}{C.5}. For the full face list \(m=2p\), and also for its subsystems, these moments
are of orders \(L\) and \(L^{3/2}\), where \(L=\log(ep)\).}

\subsection{Matrix and entropy inequalities}\label{tb:sec:matrix}

Both time ranges use the same observation: the summed second moments
control a pair of projections even when other random coefficients
remain inside the expectation.

\begin{fact}[Conditional covariance contraction]\label{tb:covariance}
Let \(Z_1,\ldots,Z_N\in\R^r\) have finite second moments, and let
\(\mathcal H\) be a sigma-field such that
\(\sum_i\E[Z_iZ_i^\top\mid\mathcal H]\preceq I_r\) almost surely.
For \(\mathcal H\)-measurable unit vectors \(u,v\),
\[
 \sum_i\E[|(u\cdot Z_i)(v\cdot Z_i)|\mid\mathcal H]\le1.
\]
If real random coefficients satisfy \(|c_i|\le a\) almost surely,
where \(a\ge0\) is constant, then
\[
 \left\|\sum_i\E[c_iZ_iZ_i^\top\mid\mathcal H]\right\|_\op\le a.
\]
No independence of the coefficients and vectors is required.
\end{fact}

\begin{proof}
Conditional Cauchy--Schwarz, followed by Cauchy--Schwarz over the
index \(i\), bounds the first sum by
\[
 \left(\sum_i\E[(u\cdot Z_i)^2\mid\mathcal H]\right)^{1/2}
 \left(\sum_i\E[(v\cdot Z_i)^2\mid\mathcal H]\right)^{1/2}\le1.
\]
For the second claim, each unit vector \(u\) gives
\(\bigl|\sum_i\E[c_i(u\cdot Z_i)^2\mid\mathcal H]\bigr|\le a\).
The matrix is symmetric, so taking the supremum of its absolute
quadratic forms proves the operator-norm bound.
\end{proof}

\usedin{\PE{appx-an:pair-contraction}{C.2} and the early estimates in
Subsection~\PN{appx-an:sec:early}{C.3}; the coefficient matrices in
\PR{Lemma}{appx-an:lem:restoration}{C.2}. Bounded face projections
control the remaining directions by powers of \(\delta\).}

After conditioning, the coefficient matrix has bounded operator norm.
The following duality transfers that bound to the Gaussian matrix, whose
nuclear norm we estimate using its positive and negative parts.

\begin{fact}[Duality between the operator and nuclear norms]\label{tb:nuclear}
(a) For every square matrix \(\Xi\) and every symmetric matrix \(M\),
\(|\tr(\Xi M)|\le\norm\Xi_\op\norm M_*\). Moreover
\(\norm M_*=\max\{\tr(\Xi M):\norm\Xi_\op\le1\}\), so that \(\norm\cdot_*\) is
a norm on symmetric matrices.

(b) If \(M=P-N\) with \(P\succeq0\) and \(N\succeq0\), then
\(\norm M_*\le\tr P+\tr N=2\tr P-\tr M\).
\end{fact}

\begin{proof}
For (a), let \(M=\sum_i\mu_ie_ie_i^\top\) be a spectral decomposition. Then
\(\tr(\Xi M)=\sum_i\mu_i\,e_i^\top\Xi e_i\) and \(|e_i^\top\Xi e_i|\le
\norm\Xi_\op\), which gives the inequality. Equality is attained at
\(\Xi\coloneqq \sum_i\operatorname{sgn}(\mu_i)e_ie_i^\top\). A supremum of linear
functionals satisfies the triangle inequality. For (b), the norm of a
positive semidefinite matrix is its trace, so the triangle inequality gives
\(\norm M_*\le\tr P+\tr N\), and \(\tr N=\tr P-\tr M\).
\end{proof}

\usedin{the final contraction in \PR{Lemma}{appx-an:lem:restoration}{C.2};
the reduction \PE{appx-an:matrix-trace-reduction}{C.22} in
\PR{Lemma}{appx-an:lem:matrix}{C.4}, with
\(P=p_\alpha\mu_\alpha\mu_\alpha^\top\) and
\(N=p_\alpha(I_r-K_\alpha)\).}

An individual logistic weight can have arbitrarily small mass, so its
logarithmic bound should be summed before the mass is estimated. The
log-sum inequality performs that summation; its integral version then
allows us to average over the independent base.

\begin{fact}[Log-sum inequality and a concave function]\label{tb:logsum}
(a) For positive numbers \(y_1,\ldots,y_l\) and \(z_1,\ldots,z_l\) with sums
\(y\) and \(z\),
\(\sum_jy_j\log(z_j/y_j)\le y\log(z/y)\).

(b) The function \((x,y)\mapsto x\log(y/x)\) is concave on \((0,\infty)^2\).
Hence, for positive integrable random variables \(X,Y\),
\(\E[X\log(Y/X)]\le\E X\,\log(\E Y/\E X)\).
The expectation on the left may equal \(-\infty\); its positive
part is integrable because \(x\log(y/x)\le y/e\).
\end{fact}

\begin{proof}
For (a), Jensen's inequality for the concave function \(\log\), with the
weights \(y_j/y\), gives \(\sum_j(y_j/y)\log(z_j/y_j)\le\log\sum_j(y_j/y)(z_j/y_j)
=\log(z/y)\). For (b), the function \(f(x,y)\coloneqq x\log y-x\log x\) has Hessian
\(\begin{pmatrix}-1/x&1/y\\1/y&-x/y^2\end{pmatrix}\), whose determinant is
zero and whose trace is negative, so the Hessian is negative semidefinite.
For the expectation bound, put \(\mu\coloneqq\E X\) and
\(\eta\coloneqq\E Y\). The tangent-plane inequality for this concave
function gives
\[
 f(X,Y)\le f(\mu,\eta)
   +f_x(\mu,\eta)(X-\mu)+f_y(\mu,\eta)(Y-\eta).
\]
The right side is integrable. Taking expectations proves the claim,
including when the negative part of the left side has infinite mean.
\end{proof}

\usedin{the reference-mass estimate in
\PR{Lemma}{appx-an:lem:matrix}{C.4}: first sum over tuples sharing
their first face, then average over \(V_0\). Both the reference mass
and the logistic mass contain the factor \(h\sqrt L\), which cancels
inside the logarithm.}
\emph{Source.} \TBc{CT06}, Theorem~2.7.1, for (a).

\subsection{Log-concavity}\label{tb:sec:logconcave}

A function \(f:\R^n\to[0,\infty)\) is \emph{log-concave} if
\(f((1-\alpha)x+\alpha y)\ge f(x)^{1-\alpha}f(y)^\alpha\) for all \(x,y\) and \(\alpha\in[0,1]\);
when \(f>0\) this says that \(\log f\) is concave. At the endpoints
\(\alpha=0,1\) the inequality means equality, without a convention
for \(0^0\). We need the fact that integrating out coordinates
preserves log-concavity: it will turn a Gaussian convolution into a
concave log normalizing constant.

\begin{fact}[Pr\'ekopa--Leindler inequality and Pr\'ekopa's theorem]
\label{tb:prekopa}
(a) Let \(0<\alpha<1\) and let \(f,g,h:\R^n\to[0,\infty)\) be bounded,
continuous, and integrable, with
\(h((1-\alpha)x+\alpha y)\ge f(x)^{1-\alpha}g(y)^\alpha\) for all \(x,y\). Then
\(\int h\ge(\int f)^{1-\alpha}(\int g)^\alpha\).

(b) Let \(f:\R^a\times\R^b\to[0,\infty)\) be bounded, continuous, and
log-concave, with \(f(z,\cdot)\) integrable for every \(z\). Then
\(z\mapsto\int f(z,y)\,dy\) is log-concave on \(\R^a\).
\end{fact}

\begin{proof}
Part (a) is the Pr\'ekopa--Leindler inequality; we use the standard
integrable-function theorem in \TBc{Gar02}, Theorem~7.1.
To deduce (b), fix \(z_0,z_1\) and \(0<\alpha<1\), and put
\(z_\alpha\coloneqq(1-\alpha)z_0+\alpha z_1\).
Log-concavity gives
\[
 f(z_\alpha,(1-\alpha)y_0+\alpha y_1)
 \ge f(z_0,y_0)^{1-\alpha}f(z_1,y_1)^\alpha.
\]
The three sections are bounded, continuous, and integrable by hypothesis,
so (a) applies to them. Its conclusion is exactly the log-concavity of
\(z\mapsto\int f(z,y)\,dy\). The endpoint cases are immediate.
\end{proof}

\emph{Source.} Pr\'ekopa~\TBc{Pre71,Pre73} and Leindler~\TBc{Lei72}.
The marginalization result is Theorem~6 of \TBc{Pre73};
\TBc{Gar02}, Section~7, treats the integral inequality.

We now apply marginal log-concavity to a Gaussian convolution. Its
logarithmic Hessian is the covariance of the tilted law minus the
identity, so concavity gives the matrix inequality needed in the proof.

\begin{fact}[Gaussian laws tilted by a log-concave weight]\label{tb:tilt}
Let \(w:\R^r\to[0,\infty)\) be bounded, continuous, and log-concave, with
\(p\coloneqq \E w(G)>0\) for \(G\sim N(0,I_r)\). Let \(\mu\) and \(K\) be the mean and
covariance matrix of the probability law with density \(w(g)\varphi_r(g)/p\).
Then \(K\preceq I_r\), and
\(\E[(GG^\top-I_r)w(G)]=p(K-I_r+\mu\mu^\top)\).
\end{fact}

\begin{proof}
The function \((z,g)\mapsto e^{-\norm{g-z}^2/2}w(g)\) is log-concave, since
\(\norm{g-z}^2\) is jointly convex in \((z,g)\), and it is bounded and
continuous, with an integrable \(g\)-section for every \(z\).
The integral is positive because \(p>0\) and the Gaussian kernel is
strictly positive. By Fact~\ref{tb:prekopa}(b), \(\ell(z)\coloneqq \log\int e^{-\norm{g-z}^2/2}
w(g)\,dg\) is concave. Since \(\norm{g-z}^2=\norm g^2-2z\cdot g+\norm z^2\),
\[
 \ell(z)=\text{const}-\tfrac12\norm z^2+\log M(z),\qquad
 M(z)\coloneqq \int e^{z\cdot g}w(g)\varphi_r(g)\,dg=p\,\E_\nu e^{z\cdot g},
\]
where \(\nu\) is the tilted law. The function \(M\) is finite and smooth: on every bounded set of
\(z\)'s, each derivative of the integrand is bounded by a constant
times \((1+\norm g)^k e^{c\norm g}\varphi_r(g)\), an integrable function.
Thus differentiation under the integral is valid, and the Hessian of \(\log M\) at \(0\) is the
covariance \(K\) of \(\nu\). A concave function of class \(C^2\) has a
negative semidefinite Hessian, so \(K-I_r\preceq0\). The identity is the
decomposition of the second moment \(\E_\nu gg^\top=K+\mu\mu^\top\).
\end{proof}

\usedin{the Gaussian-tilt calculation in
\PR{Lemma}{appx-an:lem:matrix}{C.4}, where \(w=w_\alpha\) and
\(K=K_\alpha\), and its explanation in
Subsection~\PN{sec:outline-late}{2.4}.}
\emph{Source.} The covariance contraction is a classical consequence of
Pr\'ekopa's marginalization theorem~\TBc{Pre73}, Theorem~6, as the
preceding proof shows. For smooth positive weights, it also follows
from the Brascamp--Lieb variance inequality~\TBc{BL76}, applied to
linear functions under the potential \(\norm g^2/2-\log w(g)\), whose
Hessian is at least \(I_r\). The same convolution argument covers
indicators of convex sets by the measurable form of Pr\'ekopa's
theorem. This convex-cell specialization is also used in \TBc{10}.

It remains to check that every weight produced by the derivative
expansion meets the preceding hypothesis. Its logarithm has a particularly
simple form: an affine function minus a multiple of a log-sum-exponential.

\begin{fact}[Log-concavity of the logistic weights]\label{tb:lc-weights}
With the notation of Fact~\ref{tb:logistic}, for every \(k\ge0\) and every
tuple \((\alpha_1,\ldots,\alpha_k)\), the function
\(F\pi_{\alpha_1}\cdots\pi_{\alpha_k}=\exp\{\sum_{a\le k}q_{\alpha_a}\}
(1+T)^{-k-1}\), where \(q_\alpha\coloneqq (u_\alpha\cdot x-a_\alpha)/h\), is
log-concave on \(\R^n\), and so is its composition
\(g\mapsto F\pi_{\alpha_1}\cdots\pi_{\alpha_k}(v+\sigma g)\) with any affine map.
\end{fact}

\begin{proof}
The logarithm is a sum of the affine functions \(q_{\alpha_a}\) and
\(-(k+1)\log(1+\sum_\alpha e^{q_\alpha})\). The function
\(\log(1+\sum_\alpha e^{q_\alpha})\) is convex. To include the baseline
term explicitly, put \(q_0\coloneqq0\). H\"older's inequality gives
\[
 \sum_{\alpha=0}^m e^{(1-s)q_\alpha(x)+s q_\alpha(y)}
 \le\left(\sum_{\alpha=0}^m e^{q_\alpha(x)}\right)^{1-s}
     \left(\sum_{\alpha=0}^m e^{q_\alpha(y)}\right)^s,
 \qquad 0<s<1.
\]
Taking logarithms proves the assertion. Composition
with an affine map preserves concavity.
\end{proof}

\usedin{the application of covariance contraction to every
\(w_\alpha=F\prod_a\pi_{\alpha_a}\) in
\PR{Lemma}{appx-an:lem:matrix}{C.4}.}
\emph{Source.} Convexity of log-sum-exp is standard; see
Boyd and Vandenberghe~\TBc{BV04}, Section~3.1.5, pp.~72--74.
The baseline term \(1=e^{q_0}\), with \(q_0=0\), is included in that sum.

\subsection{Polyhedra: faces, normal cones, and nearest points}\label{tb:sec:poly}

The surface estimate treats a fixed polyhedron geometrically. Unlike
the smoothing formula, its boundary integral depends only on the set,
so repeated and redundant inequalities may now be removed.
A \emph{polyhedron} is a set
\[
 Q\coloneqq \bigcap_{j=1}^m\{x\in\R^n:u_j\cdot x\le a_j\},
\]
with unit vectors \(u_j\), called \emph{face normals}, and numbers
\(a_j\), called \emph{offsets}. Each face normal points outward.
The polyhedron is closed and convex. It is
\emph{full-dimensional} if it has an interior point, and \emph{proper} if it
is neither empty nor \(\R^n\). The description is \emph{irredundant} if
deleting any one inequality changes the set; deleting redundant or repeated
inequalities one at a time turns every description into an irredundant one
of the same set. For \(j\le m\), put \(H_j\coloneqq \{x:u_j\cdot x=a_j\}\) and
\(F_j\coloneqq Q\cap H_j\). A \emph{face} of \(Q\) is \(Q\) itself or a set
\(Q\cap\{u\cdot x=c\}\) with \(u\ne0\) and \(Q\subseteq\{u\cdot x\le c\}\).
A \emph{facet} is a face of dimension \(n-1\), and a \emph{ridge}, when \(n\ge2\), is a face
of dimension \(n-2\). In the plane the facets are the sides of a polygon
and the ridges its corners, in \(\R^3\) the facets are the faces of a solid
and the ridges its edges, and on the line the facets are the endpoints of an
interval, which has no ridges. We use these standard properties of faces:
the relative interiors of the faces of \(Q\) are pairwise disjoint and cover
\(Q\); a face of a face of \(Q\) is a face of \(Q\); and a face \(G'\)
of a face \(G\) with \(G'\ne G\) lies in the relative boundary
\(G\setminus\relint G\) of \(G\) (\TBc{Roc70}, Section~18; \TBc{Zie95}, Chapter~2).

The next fact identifies the pieces that contribute to the boundary
calculation. In particular, every ridge belongs to exactly two facets;
this allows us to combine their coefficients before taking absolute values.

\begin{fact}[Facets and ridges of a polyhedron]\label{tb:facets}
Let \(Q\) be a proper, full-dimensional polyhedron in \(\R^n\) with an
irredundant description. Statements about ridges are omitted when \(n=1\).
\begin{enumerate}
\item[(a)] The face normals \(u_1,\ldots,u_m\) are pairwise distinct, each
\(F_j\) is a facet, \(Q\) has no other facets, and
\(\partial Q=F_1\cup\cdots\cup F_m\).
\item[(b)] For \(j\ne k\), \(F_j\cap F_k\subseteq H_j\cap H_k\), and
\(H_j\cap H_k\) is empty if \(u_k=-u_j\), and an affine subspace of
dimension \(n-2\) otherwise. In particular \(F_j\cap F_k\) lies in a
finite union of affine subspaces of dimension at most \(n-2\).
\item[(c)] Each \(F_j\) is a full-dimensional polyhedron in the hyperplane
\(H_j\). Write the points of \(H_j\) as \(a_ju_j+y\) with \(y\perp u_j\), and,
for \(k\ne j\) with \(u_k\ne\pm u_j\), put
\(\alpha_{jk}\coloneqq \arccos(u_j\cdot u_k)\in(0,\pi)\),
\[
 \nu_{k\mid j}\coloneqq \frac{u_k-(u_j\cdot u_k)u_j}{\sin\alpha_{jk}},\qquad
 d_{k\mid j}\coloneqq \frac{a_k-(u_j\cdot u_k)a_j}{\sin\alpha_{jk}} .
\]
Then \(F_j\) is described in the coordinate \(y\) by the inequalities
\(\nu_{k\mid j}\cdot y\le d_{k\mid j}\), and its facets are exactly the
ridges \(E_{jk}\coloneqq F_j\cap F_k\) of \(Q\) of dimension \(n-2\); the outward
unit normal of \(F_j\) at \(E_{jk}\) is \(\nu_{k\mid j}\).
\item[(d)] Every ridge of \(Q\) lies in exactly two facets, so the ridges
correspond to the unordered pairs \(\{j,k\}\) with
\(\dim(F_j\cap F_k)=n-2\); for such a pair \(u_k\ne\pm u_j\), and the
only inequalities active at a point of \(\relint E_{jk}\) are \(j\) and \(k\).
\end{enumerate}
\end{fact}

\begin{proof}
(a) Fix \(j\). Irredundancy gives a point \(x_j\) with \(u_j\cdot x_j>a_j\)
and \(u_k\cdot x_j\le a_k\) for \(k\ne j\). Let \(x_0\) be an interior point of
\(Q\), and let \(x_*\) be the point where the segment from \(x_0\) to \(x_j\)
crosses \(H_j\). For \(k\ne j\), \(u_k\cdot x_*<a_k\), because \(x_*\) is a
convex combination of \(x_0\), where the inequality is strict, with positive
weight, and of \(x_j\), where it holds weakly. Hence a neighborhood of
\(x_*\) in \(H_j\) lies in \(F_j\), and \(\dim F_j=n-1\). Two inequalities
with the same normal cannot both be irredundant, since the one with the
larger offset is implied by the other. A point of \(Q\) at
which no inequality is active is interior, and at a point where
inequality \(j\) is active the points \(x+\varepsilon u_j\) lie outside \(Q\),
so \(\partial Q=\bigcup_jF_j\). If \(G=Q\cap\{u\cdot x=c\}\) is a facet and
\(x\in\relint G\), some inequality \(j\) is active at \(x\); on \(G\) the
affine function \(u_j\cdot x-a_j\le0\) attains its maximum at a relative
interior point and is therefore zero, so \(G\subseteq H_j\), and since \(G\) has
dimension \(n-1\), \(G=Q\cap H_j=F_j\).

(b) If \(u_k=-u_j\) and \(H_j\cap H_k\ne\emptyset\), then \(Q\) would lie in
the hyperplane \(H_j\), contradicting full dimension. If \(u_k\ne\pm u_j\),
the two equations are independent.

(c) Suppose \(n\ge2\). The set \(F_j\) is the polyhedron in \(H_j\) cut out by the remaining
inequalities, which take the stated form in the coordinate \(y\) because
\(u_k\cdot(a_ju_j+y)=a_ju_j\cdot u_k+(u_k-(u_j\cdot u_k)u_j)\cdot y\); an
inequality with \(u_k=-u_j\) holds automatically on \(H_j\) since
\(F_j\ne\emptyset\). After removing inequalities redundant within \(H_j\), part (a)
applied there shows that the facets of \(F_j\) are
the faces of \(F_j\) of dimension \(n-2\), hence faces of \(Q\) of that
dimension, that is, ridges of \(Q\) in \(F_j\); each equals
\(F_j\cap H_k=F_j\cap F_k\) for the inequality \(k\) that defines it.

(d) Let \(E\) be a ridge and \(y\in\relint E\), and let \(J\) be the set of
inequalities active at \(y\). On \(E\) each affine function
\(u_i\cdot x-a_i\), \(i\in J\), is at most zero and vanishes at the relative
interior point \(y\), so it vanishes on \(E\); thus \(E\subseteq F_i\) for
\(i\in J\) and \(\dim\operatorname{span}\{u_i:i\in J\}\le2\). A single active
inequality is impossible, since then \(y\) would lie in the relative interior
of \(F_j\) and the only face of \(F_j\) containing it is \(F_j\) itself, of
dimension \(n-1\). Opposite face normals \(u_i=-u_j\) are excluded by (b). Hence
the active face normals span a plane \(\Pi\). Near \(y\), \(Q\) coincides with
\(y\) plus the product of \(\Pi^\perp\) and a planar cone \(K\subseteq\Pi\)
with nonempty interior, defined by the active inequalities and contained in
a wedge of opening \(\pi-\alpha_{jk}<\pi\); this cone has exactly two
boundary rays. Each active inequality \(i\) has \(\dim F_i=n-1\), so
\(F_i\) meets every neighborhood of \(y\) in a set of dimension \(n-1\), and
the line \(\{u_i\cdot v=0\}\) must therefore carry a boundary ray of \(K\).
With distinct face normals this leaves exactly two active inequalities.
\end{proof}

\usedin{the facet--ridge decomposition leading to
\PE{appx-an:surface-identity}{C.30} and the pairing of the two terms
attached to each ridge in \PR{Lemma}{appx-an:lem:surface}{C.5}.}
\emph{Source.} These are standard facts about faces of polyhedra; see
\TBc{Roc70}, Sections~18 and~19, and \TBc{Zie95}, Chapter~2.
The planar-cone argument above explains why precisely two facets
meet at a ridge, including for unbounded polyhedra.

To attach a tube to a face, we must know which outward directions
keep that face as the nearest part of the polyhedron. These directions
form its normal cone, determined by the active inequalities.

\begin{fact}[Normal cones]\label{tb:normalcone}
Let \(Q\) be a polyhedron as in Fact~\ref{tb:facets}. For \(x\in Q\) let
\(J(x)\coloneqq \{j:u_j\cdot x=a_j\}\), and let
\[
 N_Q(x)\coloneqq \{v\in\R^n:v\cdot(q-x)\le0\text{ for all }q\in Q\}
\]
be the \emph{normal cone} of \(Q\) at \(x\). Then
\(N_Q(x)=\{\sum_{j\in J(x)}c_ju_j:c_j\ge0\}\). In particular
\(N_Q(x)=\{0\}\) at interior points, \(N_Q(x)=\{cu_j:c\ge0\}\) for
\(x\in\relint F_j\), and \(N_Q(x)=\{c_ju_j+c_ku_k:c_j,c_k\ge0\}\), a sector of
angle \(\alpha_{jk}\) in the plane spanned by \(u_j\) and \(u_k\), for
\(x\in\relint E_{jk}\).
\end{fact}

\begin{proof}
Put \(T\coloneqq \{d:u_j\cdot d\le0\text{ for }j\in J(x)\}\). A vector \(v\) lies in
\(N_Q(x)\) if and only if \(v\cdot d\le0\) for all \(d\in T\): if \(d\in T\)
then \(x+\varepsilon d\in Q\) for small \(\varepsilon>0\), because the
inactive inequalities are strict at \(x\), and conversely \(q-x\in T\) for
\(q\in Q\). By Farkas's lemma, which says that \(v\cdot d\le0\) for all
\(d\) with \(u_j\cdot d\le0\) (\(j\in J\)) if and only if \(v\) is a
nonnegative combination of the \(u_j\), \(j\in J\), this is the stated
cone. For \(x\in\relint F_j\), \(J(x)=\{j\}\), since an index \(k\ne j\) in
\(J(x)\) would put \(x\) in \(F_j\cap F_k\), a face of \(F_j\) of dimension
at most \(n-2\), and so in the relative boundary of \(F_j\). For
\(x\in\relint E_{jk}\), \(J(x)=\{j,k\}\) by Fact~\ref{tb:facets}(d).
\end{proof}

\usedin{the choice of the normal interval over a facet and the normal
sector over a ridge in \PR{Lemma}{appx-an:lem:surface}{C.5}.}
\emph{Source.} \TBc{Roc70}, Sections~22 and~23, for Farkas's lemma and normal
cones; \TBc{RW98}, Theorem~6.46.

The nearest-point characterization makes the outward tubes disjoint.
It applies to any nonempty closed convex set; for such a set we use
\(N_Q(x)=\{v:v\cdot(q-x)\le0\text{ for all }q\in Q\}\).

\begin{fact}[Nearest point]\label{tb:projection}
Let \(Q\subseteq\R^n\) be nonempty, closed, and convex. For each
\(y\in\R^n\) there is a unique point \(\prj_Q(y)\in Q\) with
\(\norm{y-\prj_Q(y)}=\dist(y,Q)\), and \(x=\prj_Q(y)\) if and only if
\(x\in Q\) and \(y-x\in N_Q(x)\). Consequently every \(y\in\R^n\) can be
written in exactly one way as \(y=x+v\) with \(x\in Q\) and \(v\in N_Q(x)\),
and then \(\norm v=\dist(y,Q)\). A direction \(v\) lies in \(N_Q(x)\) if and
only if \(x\) is the nearest point of \(x+s v\) in \(Q\) for every \(s>0\).
\end{fact}

\begin{proof}
For \(x,q\in Q\), \(\norm{y-q}^2=\norm{y-x}^2+\norm{x-q}^2-2(y-x)\cdot(q-x)\).
If \(y-x\in N_Q(x)\), the last term is nonnegative, so
\(\norm{y-q}^2\ge\norm{y-x}^2+\norm{x-q}^2>\norm{y-x}^2\) for \(q\ne x\); thus
\(x\) is the unique nearest point. Conversely, if \((y-x)\cdot(q-x)>0\) for
some \(q\in Q\), then moving from \(x\) toward \(q\) decreases the distance
to \(y\) for small steps, by the displayed identity applied to
\(x+\varepsilon(q-x)\in Q\), so \(x\) is not nearest. For existence, choose \(q_0\in Q\). A minimizing sequence may be
restricted to the closed ball centered at \(y\) with radius
\(\norm{y-q_0}+1\); it therefore has a convergent subsequence in
finite dimension. Closedness puts its limit in \(Q\), and continuity
of the norm makes that limit a minimizer. The last statement follows because \(N_Q(x)\) is a
cone.
\end{proof}

\usedin{disjointness of the variable-width tubes in
\PR{Lemma}{appx-an:lem:surface}{C.5}; each tube point has its nearest
point in the relative interior of exactly one face.}
\emph{Source.} \TBc{RW98}, Example~2.25, for existence and uniqueness,
and Proposition~6.17 for the characterization by the normal cone.

\subsection{Surface measure, Gaussian surface integrals, and the
divergence theorem}\label{tb:sec:surface}

We integrate each facet with its ordinary Euclidean area measure.
Hausdorff measure gives a common notation for facets and ridges in
all dimensions and shows why their lower-dimensional overlaps do not
affect a surface integral.

\begin{fact}[Hausdorff measure on flat pieces]\label{tb:hausdorff}
For \(0\le k\le n\), the \(k\)-dimensional Hausdorff measure \(\Haus^k\) on
\(\R^n\) is obtained from the covering formula
\[
 \Haus^k(S)\coloneqq \lim_{\epsilon\downarrow0}\inf\Bigl\{\sum_i\omega_k
 \Bigl(\tfrac12\operatorname{diam}C_i\Bigr)^k:S\subseteq\bigcup_iC_i,\
 \operatorname{diam}C_i\le\epsilon\Bigr\},
\]
where the infimum runs over countable covers and \(\omega_k\) is
the volume of the unit ball of \(\R^k\).
For \(k=0\), every nonempty covering set contributes one, including
a singleton. The formula defines an outer measure, whose restriction
to Borel sets is the measure used here. (a) On an
affine subspace of dimension \(k\), \(\Haus^k\) is \(k\)-dimensional Lebesgue
measure; in particular \(\Haus^0\) is counting measure and \(\Haus^n\) is
Lebesgue measure on \(\R^n\). (b) \(\Haus^k\) is invariant under isometries.
(c) A set contained in a finite union of affine subspaces of dimension less
than \(k\) has \(\Haus^k\)-measure zero.
\end{fact}

Part (c) follows from (a), since each such subspace lies in a
\(k\)-dimensional affine subspace in which it is a Lebesgue null set. For a
proper, full-dimensional polyhedron \(Q\) with an irredundant description,
and for a function \(f\ge0\) or \(f\) integrable on \(\partial Q\), the
\emph{surface integral} is
\[
 \int_{\partial Q}f\,d\Haus^{n-1}=\sum_{j=1}^m\int_{F_j}f\,d\Haus^{n-1},
\]
where on each facet \(\Haus^{n-1}\) is Lebesgue measure on the hyperplane
\(H_j\) by part~(a). By Fact~\ref{tb:facets}(b) and part~(c), distinct
facets overlap in sets of \(\Haus^{n-1}\)-measure zero, so the sum is the
integral of \(f\) over the set \(\partial Q\) with respect to
\(\Haus^{n-1}\). It depends on the set \(Q\) and not on its description,
which justifies deleting redundant inequalities before computing it.
The Gaussian surface integral \(\surf(Q)=\int_{\partial Q}\Hh\,\varphi_r\,
d\Haus^{r-1}\) of Subsection~\PN{sec:outline-late}{2.4} is the case
\(n=r\) and \(f=\Hh\varphi_r\), with \(\Hh(g)\coloneqq \norm g^2-r\).
The integral is signed; we take its absolute value only after summing
the facet integrals. We set \(\surf(\varnothing)=\surf(\R^r)\coloneqq0\).
Lower-dimensional polyhedra are excluded from this definition and are
handled as exceptional levels in Fact~\ref{tb:variation}.

\usedin{the definition of \(\surf\) in
Subsection~\PN{sec:outline-late}{2.4} and
\PR{Lemma}{appx-an:lem:surface}{C.5}.}
\emph{Source.} \TBc{EG15}, Chapter~2.

Integrating a tangential divergence over a facet converts its interior
integral into contributions from its ridges. We prove the needed divergence
formula by slicing, then use a cutoff to allow an unbounded polyhedron
under precisely the stated decay assumptions.

\begin{samepage}
\begin{fact}[Gauss--Green formula (divergence theorem) for a polyhedron]\label{tb:divergence}
Let \(Q\subseteq\R^n\) (\(n\ge1\)) be a proper, full-dimensional polyhedron
with an irredundant description, facets \(F_1,\ldots,F_m\), and unit
face normals \(u_1,\ldots,u_m\). Let \(X:\R^n\to\R^n\) be of class \(C^1\), with
\(\norm{X(x)}+|\operatorname{div}X(x)|\le C(1+\norm x)^k\varphi_n(x)\) for
some \(C\) and \(k\). Then
\[
 \int_Q\operatorname{div}X\,dx=\sum_{j=1}^m\int_{F_j}X\cdot u_j\,d\Haus^{n-1}.
\]
For \(n=1\) the facets are the endpoints of an interval, \(\Haus^0\) is
counting measure, and the identity is the fundamental theorem of calculus.
\end{fact}
\end{samepage}

\begin{proof}
We first prove the formula for a compactly supported \(C^1\) vector
field. For a unit vector \(e\) and a compactly supported
\(C^1\) scalar function \(f\), we claim that
\begin{equation}
 \int_Q\partial_ef\,dx=\sum_j(e\cdot u_j)\int_{F_j}f\,d\Haus^{n-1},
 \tag{\(*\)}
\end{equation}
since the theorem follows by taking \(e=e_i\) and \(f=X_i\) and summing over
\(i\). Write \(x=w+s e\) with \(w\in e^\perp\). By Fubini's theorem,
\(\int_Q\partial_ef\,dx=\int_{e^\perp}\int_{I(w)}\partial_sf(w+s e)\,ds\,dw\),
where \(I(w)\coloneqq \{s:w+s e\in Q\}\) is a closed interval, possibly empty, a point,
or unbounded, because \(Q\) is convex. Since \(f\) tends to zero along every
line, the inner integral equals \(f(w+s_+(w)e)-f(w+s_-(w)e)\), where
\(s_\pm(w)\) are the endpoints of a nonempty \(I(w)\).
An empty interval contributes zero, and a term with an infinite
endpoint is read as zero.

Consider the upper endpoints \(x=w+s_+(w)e\). If \(x\in\relint F_j\), then
only the inequality \(j\) is active at \(x\) (Fact~\ref{tb:normalcone}), so
\(x+s e\in Q\) for small \(s>0\) unless \(e\cdot u_j>0\), and \(x\) is an upper
endpoint if and only if \(e\cdot u_j>0\). The upper endpoints that lie in no
\(\relint F_j\) lie in \(\bigcup_{j\ne k}F_j\cap F_k\), which is a subset of a
finite union of affine subspaces of dimension at most \(n-2\)
(Fact~\ref{tb:facets}(b)); its projection along \(e\) onto \(e^\perp\) is a
Lebesgue null set of \(e^\perp\). For fixed \(j\) with \(e\cdot u_j>0\), the
map \(x\mapsto x-(x\cdot e)e\) is injective on the set of upper endpoints in
\(\relint F_j\), since a line meets \(Q\) in one interval, and it is the
restriction to \(H_j\) of a linear map whose Jacobian with respect to
\(\Haus^{n-1}\) on \(H_j\) is \(e\cdot u_j\): if \(b_1,\ldots,b_{n-1}\) is an
orthonormal basis of \(u_j^\perp\), the volume of the projection along \(e\)
of the cube they span equals the volume of the parallelepiped spanned by
\(b_1,\ldots,b_{n-1},e\), which is \(|e\cdot u_j|\). Therefore
\(\int f(w+s_+(w)e)\,dw=\sum_{j:\,e\cdot u_j>0}(e\cdot u_j)\int_{F_j}f\,d\Haus^{n-1}\),
the sets \(F_j\setminus\relint F_j\) being \(\Haus^{n-1}\)-null. The same
argument for the lower endpoints gives
\(\int f(w+s_-(w)e)\,dw
=\sum_{j:\,e\cdot u_j<0}|e\cdot u_j|\int_{F_j}f\,d\Haus^{n-1}\).
Subtracting this lower-endpoint integral from the upper-endpoint
integral proves \((*)\). Facets with \(e\cdot u_j=0\) contribute
nothing. Summing over coordinate directions proves the compactly
supported vector-field formula.

For the stated field \(X\), choose smooth cutoffs \(\chi_R\) equal to
one on the ball of radius \(R\), zero outside the ball of radius
\(2R\), with \(0\le\chi_R\le1\) and
\(\norm{\nabla\chi_R}\le C/R\). Apply the formula to \(\chi_RX\).
The product rule gives
\[
 \operatorname{div}(\chi_RX)
 =\chi_R\operatorname{div}X+\nabla\chi_R\cdot X.
\]
The first term converges in integral by the assumed integrable bound;
the second has integral in absolute value at most
\((C/R)\int_{\{\norm x\ge R\}}\norm{X(x)}\,dx\to0\).
On each of the finitely many facets, polynomial times Gaussian decay
is integrable, so dominated convergence also applies to the boundary
integrals. This proves the formula without requiring separate bounds
on the component derivatives of \(X\).
\end{proof}

\usedin{the divergence calculation inside a facet in
\PR{Lemma}{appx-an:lem:surface}{C.5}. The vector field is
\(X(y)\coloneqq y\varphi_{r-1}(y)\), with divergence
\(\{r-1-\norm y^2\}\varphi_{r-1}(y)\); the boundary terms are the
ridge contributions in \PE{appx-an:surface-identity}{C.30}.}
\emph{Source.} This is the classical Gauss--Green formula specialized
to a polyhedron; see \TBc{EG15}, Chapter~5. The cutoff argument above
extends the compact-support formula to the stated Gaussian-decay
hypothesis on an unbounded polyhedron.

For the surface integral used in the paper, this divergence
calculation produces coefficients whose separate expressions contain
\(1/\sin\alpha_{jk}\). We combine the two contributions of a ridge
before estimating them. Let
\(s_j\coloneqq \int_{F_j}\varphi_r\,d\Haus^{r-1}\),
\(\rho_{jk}\coloneqq \int_{E_{jk}}\varphi_r\,d\Haus^{r-2}\), and let
\(z_{jk}\) be the projection of a ridge point onto
\(\operatorname{span}(u_j,u_k)\). Applying the divergence formula in
each facet gives
\[
 \surf(Q)=\sum_j(a_j^2-1)s_j
       -\sum_{\{j,k\}}(d_{k\mid j}+d_{j\mid k})\rho_{jk}.
\]
The offsets \(d_{k\mid j}\) from Fact~\ref{tb:facets}(c) satisfy
\[
 d_{k\mid j}+d_{j\mid k}
 =2\sin(\alpha_{jk}/2)\,
       \frac{u_j+u_k}{\norm{u_j+u_k}}\cdot z_{jk},\qquad
 |d_{k\mid j}+d_{j\mid k}|\le\alpha_{jk}\norm{z_{jk}}.
\]
Indeed, their sum is
\((1-\cos\alpha_{jk})(a_j+a_k)/\sin\alpha_{jk}\), and
\(a_j+a_k=(u_j+u_k)\cdot z_{jk}\).
The inverse sine cancels. The remaining factor \(\alpha_{jk}\)
is precisely the angle of the normal sector whose mass we use next.

\subsection{Tubes around a polyhedron}\label{tb:sec:tubes}

We now compare the paired boundary terms with Gaussian mass outside
the polyhedron. The Gaussian density changes rapidly when a boundary point has a large
normal projection. We therefore shorten the outward tube at that point.
Its width may vary along a face: the required integration formula still
follows from orthogonal coordinates and Fubini's theorem.

Let \(Q\) be a proper, full-dimensional polyhedron in \(\R^n\) with
facets \(F_j\), ridges \(E_{jk}\), and an irredundant description by
unit face normals \(u_j\). For a ridge, put
\(L_{jk}\coloneqq \operatorname{span}(u_j,u_k)\) and
\(z_{jk}\coloneqq \Pi_{jk}x\), for any \(x\in E_{jk}\). This vector is constant
on the ridge because
\(u_j\cdot z_{jk}=a_j\) and \(u_k\cdot z_{jk}=a_k\).
Let \(\eta(x)>0\) be a measurable function on the facets and ridges.
We use the same value of \(\eta\) at a point whenever it belongs to
more than one closed face; only relative interiors enter the tubes.
The outward tubes are
\begin{align*}
 \mathcal T_j&\coloneqq \{x+s u_j:x\in\relint F_j,\ 0<s<\eta(x)\},\\
 \mathcal T_{jk}&\coloneqq \{x+v:x\in\relint E_{jk},\ v\in N_Q(x),\
                         0<\norm v<\eta(x)\}.
\end{align*}
For \(x\in\relint E_{jk}\), the normal cone is a sector of angle
\(\alpha_{jk}\) in \(L_{jk}\). Write
\(S_{jk}(a)\coloneqq N_Q(x)\cap\{v:0<\norm v<a\}\); its area is
\(\alpha_{jk}a^2/2\). For \(Q=\bigcap_j\{u_j\cdot x\le a_j\}\),
the notation \(Q^s\) increases every threshold by \(s\); it does not
denote a Euclidean neighborhood.

\begin{fact}[Variable-width normal tubes]\label{tb:tubes}
(a) The tubes are measurable, and tubes over distinct facets and
ridges are pairwise disjoint.
Each tube point \(g=x+v\) has nearest point \(x\) in \(Q\), and
\(\dist(g,Q)=\norm v\).

(b) For a nonnegative measurable function \(f\),
\begin{align*}
 \int_{\mathcal T_j}f(g)\,dg
 &=\int_{F_j}\int_0^{\eta(x)} f(x+s u_j)\,ds\,d\Haus^{n-1}(x),\\
 \int_{\mathcal T_{jk}}f(g)\,dg
 &=\int_{E_{jk}}\int_{S_{jk}(\eta(x))}f(x+v)\,dv\,d\Haus^{n-2}(x).
\end{align*}
The same formulas hold whenever \(f\) is integrable on the relevant tube.

(c) Use the function \(R\) in Fact~\ref{tb:projection-moments}, built
from the face normals of \(Q\), and choose \(\eta(x)\coloneqq 1/(4R(x))\).
Then all these tubes lie in \(\{g:g\in Q^{1/R(g)}\setminus Q\}\).
For an absolute constant \(c>0\), their normal integrals obey
\begin{align*}
 \int_0^{\eta(x)}R(x+s u_j)^3\varphi_n(x+s u_j)\,ds
 &\ge c\{1+a_j^2\}\varphi_n(x),\qquad x\in F_j,\\
 \int_{S_{jk}(\eta(x))}R(x+v)^3\varphi_n(x+v)\,dv
 &\ge c\alpha_{jk}\norm{z_{jk}}\varphi_n(x),\qquad x\in E_{jk}.
\end{align*}
For \(n=1\), there are no ridge tubes.
\end{fact}

\begin{proof}
(a) Facts~\ref{tb:normalcone} and~\ref{tb:projection} identify \(x\)
as the unique nearest point of \(x+v\). The relative interiors of
distinct faces are disjoint, so the tubes are disjoint.

(b) The map \((x,s)\mapsto x+s u_j\) is an affine isometry from
\(H_j\times\R\) onto \(\R^n\). Likewise, \((x,v)\mapsto x+v\)
is an affine isometry from \((z_{jk}+L_{jk}^\perp)\times L_{jk}\)
onto \(\R^n\). The domains are measurable in these coordinates because \(\eta\)
is measurable; their images under the affine isometries are therefore
measurable as well. The Jacobians are one. Apply Fubini to the indicator
of the domain \(0<s<\eta(x)\), or \(0<\norm v<\eta(x)\) in the
normal sector. No derivative of \(\eta\) is involved, and no
smoothness of \(\eta\) is needed.

(c) For a facet displacement \(v=s u_j\), and for a ridge displacement
\(v\in L_{jk}\), respectively,
\[
 |x\cdot v|\le\mathcal M(x)\norm v\le\tfrac14.
\]
The ridge case uses \(x\cdot v=z_{jk}\cdot v\). Since
\(\norm v\le1/4\), Fact~\ref{tb:gauss-coord}(b) gives
\(\varphi_n(x+v)\ge e^{-1/4-1/32}\varphi_n(x)\).
The Lipschitz property and \(R\ge1\) also give
\[
 \tfrac34R(x)\le R(x+v)\le\tfrac54R(x).
\]
The facet integral is therefore at least
\(cR(x)^3\eta(x)\varphi_n(x)=c'R(x)^2\varphi_n(x)\), which
dominates \(c''(1+a_j^2)\varphi_n(x)\) because
\(R(x)\ge1+|u_j\cdot x|=1+|a_j|\). The ridge integral is at least
\(c\alpha_{jk}R(x)^3\eta(x)^2\varphi_n(x)
=c'\alpha_{jk}R(x)\varphi_n(x)\); use
\(R(x)\ge\norm{z_{jk}}\).

Finally, a tube point \(g=x+v\) lies outside \(Q\), and each
unit-normal inequality defining \(Q\) satisfies
\[
 u_j\cdot g-a_j\le\norm v<\frac1{4R(x)}\le\frac1{R(g)}.
\]
This proves the inclusion in (c), even if the threshold system used
to define \(Q^s\) also contains redundant inequalities.
\end{proof}

Together with the paired ridge coefficients derived above, part (c)
gives
\[
 |\surf(Q)|\le C\E\left[R(G)^3
              \ind\{G\in Q^{1/R(G)}\setminus Q\}\right].
\]
The power three serves both dimensions: integration along a facet
uses one power of the tube width, leaving \(R^2\), while integration
over a ridge sector uses two, leaving \(R\).

\usedin{\PE{appx-an:adaptive-tubes}{C.29} in
\PR{Lemma}{appx-an:lem:surface}{C.5}. After setting
\(Q'\coloneqq (Q-V_0)/\sigma\), the event is
\(V_0+\sigma G\in Q^{\sigma/R(G)}\setminus Q\).
Conditional on \(G\), its width is fixed; concentration of \(V_0\)
and Fact~\ref{tb:projection-moments} give the averaged surface bound.}
\emph{Source.} The facet construction is Nazarov's outward-normal
argument~\TBc{Naz03}; see the detailed proof in \TBc{CCK17N},
Lemmas~1 and~2. Disjointness over arbitrary faces follows from the
classical projection characterization in \TBc{RW98}, Proposition~6.17.
The variable widths, ridge-sector integrals, and weighted estimates
needed here are derived above. For the related use of facet and ridge
tubes in Gaussian approximation, see \TBc{10}.

\subsection{The coarea formula and the first variation of a Gaussian
integral}\label{tb:sec:coarea}

The logistic mixture moves all thresholds together. To differentiate
its Gaussian integral, we integrate first over the level sets of the
maximum constraint. The coarea formula provides this change of variables.

\begin{fact}[Coarea formula]\label{tb:coarea}
Let \(\psi:\R^n\to\R\) be Lipschitz, so that \(\nabla\psi\) exists almost
everywhere. For every measurable \(f\ge0\),
\[
 \int_{\R^n}f\,\norm{\nabla\psi}\,dx=\int_{-\infty}^\infty
 \Bigl(\int_{\{\psi=s\}}f\,d\Haus^{n-1}\Bigr)ds ,
\]
For signed \(f\), the same identity holds if
\(\int |f|\norm{\nabla\psi}\,dx<\infty\), in particular if
\(f\in L^1(\R^n)\), since \(\psi\) is Lipschitz.
\end{fact}

\emph{Source.} \TBc{EG15}, Section~3.4. The proofs use the formula only for
a maximum of finitely many affine functions with gradients of equal length, and
for these Fact~\ref{tb:variation} gives a direct proof from Fubini's theorem.

For the finite maximum used here, the coarea formula can be proved
directly by partitioning the space according to the maximizing face.
This proof also identifies the exceptional levels and explains why the
independent Gamma shift permits differentiation at every fixed threshold.

\begin{fact}[First variation of a polyhedral sublevel integral]\label{tb:variation}
Let \(u_1,\ldots,u_m\) be unit vectors in \(\R^n\) (\(m\ge1\)), \(\kappa>0\),
\(b_1,\ldots,b_m\in\R\), and
\[
 \psi(g)\coloneqq \max_{j\le m}\{\kappa\,u_j\cdot g-b_j\},\qquad Q_s\coloneqq \{g:\psi(g)\le s\}.
\]
(a) For \(f\ge0\) measurable, or \(f\) integrable, and \(s\in\R\),
\[
 \int_{Q_s}f\,dg=\kappa^{-1}\int_{-\infty}^s\Bigl(\int_{\{\psi=s'\}}f\,d\Haus^{n-1}
 \Bigr)ds'.
\]
(b) If \(f\) is integrable, then \(\Lambda(s)\coloneqq \int_{Q_s}f\,dg\) is absolutely
continuous, and for almost every \(s\)
\[
 \Lambda'(s)=\kappa^{-1}\int_{\{\psi=s\}}f\,d\Haus^{n-1},\qquad
 |\Lambda'(s)|\le\kappa^{-1}\sum_{j=1}^m\int_{\{\kappa u_j\cdot g-b_j=s\}}|f|\,d\Haus^{n-1}.
\]
(c) If \(s>\inf\psi\), then \(Q_s\) is a proper, full-dimensional polyhedron
with face normals \(u_j\) and offsets \((s+b_j)/\kappa\), \(\{\psi=s\}=\partial Q_s\),
and \(\int_{\{\psi=s\}}f\,d\Haus^{n-1}\) is the surface integral of \(f\) over
\(\partial Q_s\) of Subsection~\ref{tb:sec:surface}. For \(s<\inf\psi\) the set
\(Q_s\) is empty, and \(s=\inf\psi\) is the only other level at which \(Q_s\)
may fail to be full-dimensional.

(d) For \(\Lambda\) from (b), assume in addition that it is
Lipschitz. If \(S\) is a real random variable with a density, then \(s\mapsto\E\Lambda(s+S)\) is differentiable with derivative
\(\E\Lambda'(s+S)\), where \(\Lambda'\) is defined almost everywhere.
\end{fact}

\begin{proof}
(a) Let \(C_j\) be the set where \(j\) is the smallest index attaining the
maximum defining \(\psi\); the Borel sets \(C_1,\ldots,C_m\) partition
\(\R^n\). On \(C_j\), \(\psi(g)=\kappa v-b_j\) in the coordinates
\(g=v u_j+y\), \(y\perp u_j\). By Fact~\ref{tb:fubini}, in these coordinates,
\[
 \int_{C_j\cap Q_s}f\,dg=\int_{u_j^\perp}\int_{-\infty}^{(s+b_j)/\kappa}
 (f\ind_{C_j})(v u_j+y)\,dv\,dy .
\]
The substitution \(s'=\kappa v-b_j\) turns the inner integral into
\(\kappa^{-1}\int_{-\infty}^s(f\ind_{C_j})(v u_j+y)\,ds'\) with
\(v=(s'+b_j)/\kappa\), and a further application of Fubini's theorem
exchanges the integrals. For fixed \(s'\), the set of \(v u_j+y\) in \(C_j\)
with this \(v\) is the slice \(C_j\cap\{\psi=s'\}\), and on the hyperplane
containing it \(\Haus^{n-1}\) is Lebesgue measure \(dy\)
(Fact~\ref{tb:hausdorff}(a)). Hence
\(\int_{C_j\cap Q_s}f\,dg=\kappa^{-1}\int_{-\infty}^s\int_{C_j\cap\{\psi=s'\}}
f\,d\Haus^{n-1}\,ds'\). Summing over \(j\) gives (a), since the slices
\(C_j\cap\{\psi=s'\}\) partition \(\{\psi=s'\}\).

(b) Applying (a) to \(|f|\) and letting \(s\to\infty\) shows that
\(I(s')\coloneqq \kappa^{-1}\int_{\{\psi=s'\}}f\,d\Haus^{n-1}\) is integrable on
\(\R\), and (a) says \(\Lambda(s)=\int_{-\infty}^sI(s')\,ds'\). Hence \(\Lambda\) is
absolutely continuous, and \(\Lambda'=I\) almost everywhere by Lebesgue's
differentiation theorem. The bound follows because
\(\{\psi=s\}\subseteq\bigcup_j\{\kappa u_j\cdot g-b_j=s\}\), a union of \(m\)
hyperplanes.

(c) The set \(Q_s\) is the polyhedron \(\bigcap_j\{u_j\cdot g\le(s+b_j)/\kappa\}\),
and \(\psi\) is convex and continuous. If \(\psi(g)=s>\inf\psi\), choose \(g'\)
with \(\psi(g')<s\); for \(\tau>0\), \(g=(g_\tau+\tau g')/(1+\tau)\) with
\(g_\tau\coloneqq g+\tau(g-g')\), so convexity gives
\(\psi(g_\tau)\ge s+\tau(s-\psi(g'))>s\). Thus \(g\) is a limit of points
outside \(Q_s\) and \(g\in\partial Q_s\); the open set \(\{\psi<s\}\) is
nonempty and lies in \(Q_s\), which is therefore full-dimensional, and it
is a proper subset because \(\psi(g)\ge\kappa u_1\cdot g-b_1\to\infty\) along
\(u_1\). A point with \(\psi<s\) is interior, so \(\partial Q_s=\{\psi=s\}\). After
the redundant and repeated inequalities are deleted, \(\partial Q_s=\{\psi=s\}\)
is the union of the sets \(Q_s\cap\{\kappa u_j\cdot g-b_j=s\}\) of the
remaining inequalities, which are faces of \(Q_s\) or empty. A face of
dimension \(n-1\) is a facet (Fact~\ref{tb:facets}(a)) and any other face is
\(\Haus^{n-1}\)-null, so the integral over \(\{\psi=s\}\) is the sum over the
facets of the irredundant description.

(d) The function \(\Lambda\) is bounded by \(\int|f|\), so the expectation
exists. A Lipschitz function is absolutely continuous and differentiable outside a
Lebesgue null set \(\mathcal N\). The random variable \(s+S\) avoids
\(\mathcal N\) almost surely, because \(S\) has a density, so the difference
quotients of \(\Lambda(s+S)\) converge almost surely and are bounded by the
Lipschitz constant. Dominated convergence gives the claim.
\end{proof}

\usedin{the first-variation paragraph in the proof of
\PR{Lemma}{appx-an:lem:matrix}{C.4}, immediately after
\PE{appx-an:gamma-mixture}{C.25}. Here
\(\psi(g)=\max_j\{u_j\cdot(V_0+\sigma g)-a_j\}\),
\(\kappa=\sigma\), \(b_j=a_j-u_j\cdot V_0\), and
\(f(g)=(\norm g^2-r)\varphi_r(g)\). Parts (a)--(c) give
\(\sigma^{-1}\surf((Q(a+s\mathbf1)-V_0)/\sigma)\) for almost every
\(s\). The independent shift \(-h\log T_\ell\), with
\(T_\ell\sim\operatorname{Gamma}(\ell,1)\), has a density, so part (d)
permits averaging these derivatives. A uniform bound
\(Cm(r+1)/\sigma\) follows from (b) and the Gaussian integral over
each supporting hyperplane; it justifies the interchange and does
not enter the final dimension-free bound.}

\clearpage
\subsection{Where each fact is used}\label{tb:sec:index}

The table follows the order of the proof. Each row identifies the
operation for which the corresponding facts are needed.

\medskip
\begin{center}
\small
\renewcommand{\arraystretch}{1.18}
\begin{tabular}{@{}p{0.27\textwidth}p{0.68\textwidth}@{}}
Facts~\ref{tb:fubini}, \ref{tb:taylor} &
Conditioning, Taylor remainders, and exchange of expectations with
convergent series in Appendices B--C.\\
Facts~\ref{tb:bernoulli}, \ref{tb:heat} &
Mask differentiation, the heat equation, and the signed Hessian
comparison in Subsection~\PN{appx-an:sec:path}{C.2}.\\
Facts~\ref{tb:faa}, \ref{tb:logistic}, \ref{tb:covariance} &
Derivative envelopes in Appendix~\PN{appx-sec:smoothing}{B}; early
bounds in \PN{appx-an:sec:early}{C.3} and coefficient matrices in
\PR{Lemma}{appx-an:lem:restoration}{C.2}.\\
Facts~\ref{tb:stein}, \ref{tb:nuclear} &
Two Gaussian integrations by parts and matrix contraction in
\PR{Lemma}{appx-an:lem:restoration}{C.2}.\\
Facts~\ref{tb:gauss-coord}, \ref{tb:logsum} &
Reference-face mean bounds and summation of reference masses in
\PR{Lemmas}{appx-an:lem:mean}{C.3} and~\PN{appx-an:lem:matrix}{C.4}.\\
Facts~\ref{tb:prekopa}--\ref{tb:lc-weights} &
Logistic-tilt covariance contraction in
\PR{Lemma}{appx-an:lem:matrix}{C.4}.\\
Facts~\ref{tb:facets}--\ref{tb:divergence} &
Surface measure and the facet--ridge identity in
\PR{Lemma}{appx-an:lem:surface}{C.5}.\\
Facts~\ref{tb:tubes}, \ref{tb:projection-moments} &
Adaptive normal tubes and Gaussian projection moments in
\PR{Lemma}{appx-an:lem:surface}{C.5}.\\
Facts~\ref{tb:coarea}, \ref{tb:variation} &
First variation of the Gamma mixture in
\PR{Lemma}{appx-an:lem:matrix}{C.4}.\\
\end{tabular}
\end{center}

\subsection{Sources}\label{tb:sec:sources}

\begin{description}\footnotesize\setlength{\itemsep}{0pt}\setlength{\parskip}{0pt}
\item[{[BL76]}] H.~J. Brascamp and E.~H. Lieb (1976). On extensions of the
Brunn--Minkowski and Pr\'ekopa--Leindler theorems, including inequalities
for log concave functions, and with an application to the diffusion
equation. \emph{Journal of Functional Analysis} \textbf{22}, 366--389.
\item[{[BV04]}] S. Boyd and L. Vandenberghe (2004).
\emph{Convex Optimization}. Cambridge University Press, Cambridge.
\item[{[CCK17N]}] V. Chernozhukov, D. Chetverikov, and K. Kato (2017).
Detailed proof of Nazarov's inequality. arXiv:1711.10696.
\item[{[Com74]}] L. Comtet (1974). \emph{Advanced Combinatorics}. Reidel,
Dordrecht.
\item[{[CT06]}] T.~M. Cover and J.~A. Thomas (2006). \emph{Elements of
Information Theory}, 2nd edition. Wiley.
\item[{[EG15]}] L.~C. Evans and R.~F. Gariepy (2015). \emph{Measure Theory
and Fine Properties of Functions}, revised edition. CRC Press.
\item[{[Fol99]}] G.~B. Folland (1999). \emph{Real Analysis: Modern
Techniques and Their Applications}, 2nd edition. Wiley.
\item[{[Gar02]}] R.~J. Gardner (2002). The Brunn--Minkowski inequality.
\emph{Bulletin of the American Mathematical Society} \textbf{39}, 355--405.
\item[{[10]}] Z. Gao and W. Zhang (2026). Cubic-root Gaussian approximation
under unrestricted covariance. arXiv:2608.30221v1.
\item[{[Lei72]}] L. Leindler (1972). On a certain converse of H\"older's
inequality II. \emph{Acta Scientiarum Mathematicarum (Szeged)} \textbf{33},
217--223.
\item[{[Naz03]}] F. Nazarov (2003). On the maximal perimeter of a convex
set in \(\R^n\) with respect to a Gaussian measure. In V.~D. Milman
and G. Schechtman (eds.), \emph{Geometric Aspects of Functional Analysis},
Lecture Notes in Mathematics \textbf{1807}, 169--187. Springer.
\item[{[Pre71]}] A. Pr\'ekopa (1971). Logarithmic concave measures with
application to stochastic programming. \emph{Acta Scientiarum Mathematicarum
(Szeged)} \textbf{32}, 301--316.
\item[{[Pre73]}] A. Pr\'ekopa (1973). On logarithmic concave measures and
functions. \emph{Acta Scientiarum Mathematicarum (Szeged)} \textbf{34},
335--343.
\item[{[Roc70]}] R.~T. Rockafellar (1970). \emph{Convex Analysis}. Princeton
University Press.
\item[{[RW98]}] R.~T. Rockafellar and R.~J.-B. Wets (1998). \emph{Variational
Analysis}. Springer, Berlin.
\item[{[Ste81]}] C.~M. Stein (1981). Estimation of the mean of a multivariate
normal distribution. \emph{The Annals of Statistics} \textbf{9}, 1135--1151.
\item[{[Zie95]}] G.~M. Ziegler (1995). \emph{Lectures on Polytopes}.
Springer, New York.
\end{description}



% Fallback paper numbers below are checked against the v3 auxiliary file.

\makeatletter
\providecommand{\E}{\mathbb E}
\providecommand{\Pp}{\mathbb P}
\providecommand{\R}{\mathbb R}
\providecommand{\tr}{\operatorname{tr}}
\providecommand{\op}{\mathrm{op}}
\providecommand{\norm}[1]{\left\lVert#1\right\rVert}
\providecommand{\ind}{\mathbf 1}
\providecommand{\dd}{\,d}
\providecommand{\Hh}{H_2}
\providecommand{\surf}{\mathcal J}
\providecommand{\relint}{\operatorname{relint}}
\providecommand{\aff}{\operatorname{aff}}
\providecommand{\cone}{\operatorname{cone}}
\providecommand{\dist}{\operatorname{dist}}
\providecommand{\prj}{\operatorname{pr}}
\providecommand{\Haus}{\mathcal H}
\providecommand{\TBc}[1]{[#1]}
% \PN{label}{C.5} prints a number of the paper (a \ref in the paper);
% \PR{Lemma}{label}{C.5} prints "Lemma C.5"; \PE{label}{C.21} an equation.
\providecommand{\PN}[2]{\ifdefined\toolboxinpaper\ref{#1}\else#2\fi}
\providecommand{\PR}[3]{#1~\PN{#2}{#3}}
\providecommand{\PE}[2]{(\PN{#1}{#2})}
\providecommand{\usedin}[1]{\par\smallskip\noindent{\small\emph{Used in:} #1}\par\smallskip}
\@ifundefined{c@fact}{\newtheorem{fact}{Fact}[section]}{}
\makeatother

\section{Toolbox: facts from analysis and geometry}\label{app:toolbox}

The late interpolation estimate reduces high derivatives to two objects:
a Gaussian matrix weighted by the logistic approximation, and a signed
Gaussian integral over a polyhedral boundary. This appendix supplies the
analysis and geometry behind those reductions. The calculus facts also
justify the early estimate and the smoothing argument in
Appendices~\PN{appx-sec:smoothing}{B} and~\PN{appx-sec:support}{C}.
Each fact is stated with the assumptions needed for its use, followed by
a proof or a standard reference and a precise pointer to the manuscript.

We first treat differentiation, Gaussian integration by parts, and
covariance contraction. Log-concavity then controls the Gaussian law
tilted by a logistic weight. The geometric facts explain how to pair the
two contributions at each ridge and compare the resulting boundary
integral with the mass of short outward tubes. Because the faces are
flat, the required area and tube formulas follow from orthogonal
coordinates and Fubini's theorem. Finally, the first-variation formula
connects the logistic mixture to the surface estimate, including the
exceptional levels at which the polyhedron loses dimension.
Subsection~\ref{tb:sec:index} records these dependencies in one table.

\subsection{Notation}\label{tb:sec:notation}

The smoothing and interpolation arguments take place in the covariance
support \(\R^r\). They use a fixed list of \(m=2p\) unit face normals
\(u_1,\ldots,u_m\), which may repeat, and the threshold sets
\[
 A_a\coloneqq\{x\in\R^r:u_\alpha\cdot x\le a_\alpha,
                          \ \alpha=1,\ldots,m\},\qquad a\in\R^m.
\]
All thresholds are finite; inconsistent inequalities give the empty set.
We keep the full list when defining the logistic approximation, even if
some inequalities are redundant. We delete redundant inequalities only
when describing the facets of a fixed full-dimensional polyhedron.
For a specified system of unit-normal inequalities, \(Q^s\) increases
each threshold by \(s\); this notation depends on the system and does
not denote a Euclidean neighborhood. We write \(L\coloneqq\log(ep)\).

In the general facts below, \(\R^n\) carries the inner product
\(x\cdot y\) and Euclidean norm \(\norm x\). This dimension is \(r\)
in the support, or \(r-1\) when we work within a facet. We write \(dx\) for Lebesgue measure, \(\varphi_n\) and
\(\Phi_n\) for the standard Gaussian density and law on \(\R^n\), and \(G\)
for a standard Gaussian vector. For a linear subspace \(L\), \(\varphi_L\)
is the standard Gaussian density on \(L\) with respect to Lebesgue measure
on \(L\). In dimension zero, Lebesgue and Gaussian measure both assign
mass one to the origin, and \(\varphi_0(0)\coloneqq1\). This convention
includes the tangential calculation at an endpoint when \(r=1\).
For a nonempty set \(S\), \(\ind_S\) is its indicator function, \(\aff S\)
is the smallest affine subspace containing it, the dimension of \(S\) is
that of \(\aff S\), and \(\relint S\) is the interior of \(S\) relative to
\(\aff S\). The distance from a point \(y\) to a set \(S\) is
\(\dist(y,S)\coloneqq \inf_{x\in S}\norm{y-x}\).

For an integer \(q\ge1\) and \(f\in C^q(\R^n)\), the derivative of order \(q\) in the directions
\(z_1,\ldots,z_q\in\R^n\) is
\[
 D^qf(x)[z_1,\ldots,z_q]
 \coloneqq \partial_{s_1}\cdots\partial_{s_q}f(x+s_1z_1+\cdots+s_qz_q)\Big|_{s=0}
 =\sum_{i_1,\ldots,i_q}\partial_{i_1\cdots i_q}f(x)\,
   z_{1,i_1}\cdots z_{q,i_q},
\]
which is symmetric and multilinear in the directions. We write
\(D^qf(x)[z^{\otimes q}]\) when all directions equal \(z\), with
\(D^0f\coloneqq f\). The first derivative satisfies
\(Df(x)[z]=\nabla f(x)\cdot z\), and we identify the second derivative
with its Hessian matrix. The Loewner order \(M\preceq N\) means that
\(N-M\) is positive semidefinite, \(\norm\Xi_\op\) is the operator norm
of a matrix \(\Xi\), and \(\norm M_*\) is the \emph{nuclear norm}, which for
a symmetric matrix \(M\) with eigenvalues \(\mu_i\) equals \(\sum_i|\mu_i|\).

\subsection{Integration and calculus}\label{tb:sec:calc}

We repeatedly condition on an independent component and then average the
resulting bound. The following theorem specifies when these operations,
and the corresponding exchanges of integrals or sums, are valid.

\begin{fact}[Fubini and Tonelli]\label{tb:fubini}
Let \((S_1,\mu_1)\) and \((S_2,\mu_2)\) be \(\sigma\)-finite measure
spaces and let \(f\) be a measurable function on \(S_1\times S_2\). If
\(f\ge0\), then \(s_1\mapsto\int f(s_1,s_2)\,\mu_2(ds_2)\) is measurable
and
\[
 \int\Bigl(\int f\,d\mu_2\Bigr)d\mu_1
 =\int\Bigl(\int f\,d\mu_1\Bigr)d\mu_2
 =\int f\,d(\mu_1\otimes\mu_2)\in[0,\infty].
\]
If \(\int|f|\,d(\mu_1\otimes\mu_2)<\infty\), the same equalities hold, and
the inner integrals are finite almost everywhere.
\end{fact}

Three special cases occur in the proofs. If \(X\) and \(Y\) are
independent random elements and either \(f\ge0\) or \(f(X,Y)\) is
integrable, then
\[
 \E f(X,Y)=\E g(X),\qquad g(x)\coloneqq \E f(x,Y),
\]
because the joint law is the product of the marginal laws. This is the
formal content of every argument that is carried out for a fixed value of
an independent component such as \(V_0\) and then averaged over it. With
\(\mu_2\) the counting measure, the theorem says that
\(\E\sum_k\xi_k=\sum_k\E\xi_k\) for nonnegative terms, and for terms with
\(\sum_k\E|\xi_k|<\infty\). For an orthogonal decomposition
\(\R^n=L\oplus L^\perp\), Lebesgue measure on \(\R^n\) is the product of the
Lebesgue measures on \(L\) and on \(L^\perp\), so that
\(\int_{\R^n}f\,dx=\int_{L^\perp}\int_Lf(l+w)\,dl\,dw\).

The final averaging in \PR{Lemma}{appx-an:lem:surface}{C.5} is an
instance of the first case. Let \(V_0\in\R^r\) and
\(G\sim N(0,I_r)\) be independent, and let \(\sigma>0\).
For a fixed threshold set \(Q\), it gives
\[
 \E\!\left[R(G)^3\ind\{V_0+\sigma G\in
                   Q^{\sigma/R(G)}\setminus Q\}\right]
 =\E_G\!\left[R(G)^3\Pp_{V_0}\{V_0+\sigma G\in
                   Q^{\sigma/R(G)}\setminus Q\}\right].
\]
Here \(R\) is the projection radius in
Fact~\ref{tb:projection-moments}, and \(Q^s\) shifts every
unit-normal threshold outward by \(s\). Conditional on \(G\), both
the width \(\sigma/R(G)\) and the translation \(\sigma G\) are fixed.
Concentration of \(V_0\) can therefore be applied directly. Tonelli
requires no integrability check at this stage because the integrand
is nonnegative.

\usedin{conditioning on independent components in Appendices~\PN{appx-sec:smoothing}{B}
and~\PN{appx-sec:support}{C}; absolute convergence of the translated-weight
series in \PR{Lemma}{appx-an:lem:restoration}{C.2}; final averaging in
\PR{Lemma}{appx-an:lem:surface}{C.5}; conditioning on independent pairs
in the retained proof in \PR{Appendix}{appx-app:concentration}{D}.}
\emph{Source.} \TBc{Fol99}, Section~2.5.

To differentiate the interpolation, we first justify moving a derivative
inside an expectation and then expand along a line segment. An integrable
bound is needed for the former operation; the latter is ordinary
one-dimensional Taylor expansion applied to a multivariate function.

\begin{fact}[Differentiation under the integral sign; Taylor's formula]
\label{tb:taylor}
(a) Let \(I\) be an open interval and let \(f(t,\omega)\) be
measurable and integrable in \(\omega\) for each \(t\in I\).
Suppose that, outside one null set, it is differentiable in \(t\).
On each compact subinterval \(J\subset I\), assume
\(|\partial_t f(t,\omega)|\le g_J(\omega)\) for all \(t\in J\),
where \(g_J\) is integrable. Then \(\frac d{dt}\E f(t,\cdot)
=\E\,\partial_t f(t,\cdot)\) on \(I\).

(b) For an integer \(q\ge1\), \(f\in C^q(\R^n)\), and \(x,v\in\R^n\),
\[
 f(x+v)=\sum_{k=0}^{q-1}\frac1{k!}D^kf(x)[v^{\otimes k}]
 +\int_0^1\frac{(1-s)^{q-1}}{(q-1)!}D^qf(x+sv)[v^{\otimes q}]\,ds.
\]
For \(q=1\) this is \(f(x+v)-f(x)=\int_0^1Df(x+sv)[v]\,ds\).
\end{fact}

\begin{proof}
For (a), fix an interior time and a compact neighborhood \(J\).
The mean value theorem bounds nearby difference quotients by \(g_J\),
and they converge pointwise, so dominated convergence applies. For (b),
the function \(s\mapsto f(x+sv)\) is of class \(C^q\) with \(k\)th derivative
\(D^kf(x+sv)[v^{\otimes k}]\); apply the one-dimensional Taylor formula with
integral remainder at \(s=0\) and evaluate it at \(s=1\).
\end{proof}

\usedin{the signed Hessian comparison and its Taylor expansion in
Subsection~\PN{appx-an:sec:path}{C.2}; the restoration of the cubic term
to the full path in \PR{Lemma}{appx-an:lem:identity}{C.1}; the segment
integral in \PR{Lemma}{appx-lem:endpoint-gap}{B.3}.}
\emph{Source.} \TBc{Fol99}, Section~2.3.

Differentiating the thinning probability changes one mask at a time.
Here a mask is a variable that retains a summand when it equals one
and deletes it when it equals zero. The following finite-product identity
isolates the change in retention probability; changes in amplitude and
Gaussian variance are separate contributions.

\begin{fact}[Differentiating independent Bernoulli masks]\label{tb:bernoulli}
Let \(B_1,\ldots,B_N\) be independent Bernoulli variables with parameter
\(q\in(0,1)\). For a function \(H:\{0,1\}^N\to\R\), write
\(B^{i,b}\) for the vector obtained by setting its \(i\)th entry to
\(b\in\{0,1\}\). Then
\[
 \frac{d}{dq}\E_q H(B)
 =\sum_{i=1}^N\E_q\{H(B^{i,1})-H(B^{i,0})\}.
\]
If \(q=q(t)\) and \(H=H_t\) are differentiable, then
\[
 \frac{d}{dt}\E_{q(t)}H_t(B)
 =\E_{q(t)}\partial_t H_t(B)
 +q'(t)\sum_i\E_{q(t)}
       \{H_t(B^{i,1})-H_t(B^{i,0})\}.
\]
\end{fact}

\begin{proof}
Temporarily give the masks separate parameters \(q_1,\ldots,q_N\).
Conditioning on all masks except the \(i\)th makes their expectation
affine in \(q_i\), with derivative equal to the insertion--deletion
difference in the first display. Differentiating along
\(q_1=\cdots=q_N=q\) gives the sum. The second formula is the chain
rule. Independent random arguments may be included by conditioning
and then applying Fact~\ref{tb:taylor}(a).
\end{proof}

\usedin{Subsection~\PN{appx-an:sec:path}{C.2}, before
\PR{Lemma}{appx-an:lem:identity}{C.1}, to differentiate the masks in
the Bernoulli--Gauss path.}

In the paper, the amplitude and retention probability are
\(\gamma=1+t\) and \(\varpi=(1-t)/(1+t)^2\), for \(0<t<1\).
The support-coordinate summands \(\zeta_i^0\) are centered, with
\(\sum_i\E\zeta_i^0(\zeta_i^0)^\top=I_r\). Write
\(Y_i\coloneqq\gamma\mathbb B_{i,t}\zeta_i^0\) and
\(U_t\coloneqq\sum_iY_i+\sqrt t\,G\), where all masks and \(G\)
are mutually independent and independent of the summands. To expand the
contribution of summand \(i\), remove it and put \(H_i\coloneqq U_t-Y_i\).
Let \(\zeta_i\) be an independent copy of \(\zeta_i^0\), with all copies
independent of each other and of the path variables. The signed Hessian
comparison in~\PE{appx-an:signed-hessian}{C.6} has scalar weight
\[
 \nu_t(s)\coloneqq \frac{4-2(3-t)s}{1+t},\qquad 0\le s\le1.
\]
Direct integration gives
\[
 \int_0^1\nu_t(s)\,ds=1,\qquad
 \int_0^1s\nu_t(s)\,ds=\frac{2t}{3(1+t)}.
\]
For a twice continuously differentiable scalar function \(g\),
Taylor's formula and the first identity give
\[
 g(0)-\int_0^1\nu_t(s)g(s)\,ds
 =-\frac{2t}{3(1+t)}g'(0)
   -\int_0^1\nu_t(s)\int_0^s(s-u)g''(u)\,du\,ds.
\]
For the logistic approximation \(F=F_{a,h}\) defined below, take
\(g(s)\coloneqq D^2F(H_i+s\gamma\zeta_i)[\zeta_i,\zeta_i]\).
Then \(g'(0)=\gamma D^3F(H_i)[\zeta_i^{\otimes3}]\), so
the factor \(1/2\) in the signed Hessian comparison turns the first
term into
\(-\tfrac12\{2t/(3\gamma)\}\gamma
 D^3F(H_i)[\zeta_i^{\otimes3}]
 =-tD^3F(H_i)[\zeta_i^{\otimes3}]/3\). This explains the
cubic coefficient before the derivatives are restored to the full path.
The weight \(\nu_t\) is signed; these identities use its
integrals, not a probability interpretation.

Repeated differentiation of the logistic approximation produces products
of face probabilities. The partition form of the higher-order chain
rule records which directions differentiate the same inner factor.

\begin{fact}[Fa\`a di Bruno's formula]\label{tb:faa}
Let \(\phi\) be a smooth function on an open interval of \(\R\) and let
\(g:\R^n\to\R\) be smooth with values in that interval. For \(q\ge1\), let
\(\Pi_q\) be the set of partitions \(\mathcal P\) of \(\{1,\ldots,q\}\) into
nonempty blocks, let \(|\mathcal P|\) be the number of blocks, and for a
block \(J\) let \(z_J\) be the list of directions \(z_a\) with
\(a\in J\). Then
\[
 D^q(\phi\circ g)(x)[z_1,\ldots,z_q]
 =\sum_{\mathcal P\in\Pi_q}\phi^{(|\mathcal P|)}(g(x))
  \prod_{J\in\mathcal P}D^{|J|}g(x)[z_J].
\]
\end{fact}

\begin{proof}
For \(q=1\) this is the chain rule. Differentiate the right side for \(q\)
in the direction \(z_{q+1}\). By the product rule, the derivative either
falls on \(\phi^{(|\mathcal P|)}(g(x))\), which gives
\(\phi^{(|\mathcal P|+1)}(g(x))\,Dg(x)[z_{q+1}]\) and adds the block
\(\{q+1\}\) to \(\mathcal P\), or falls on one factor
\(D^{|J|}g(x)[z_J]\), which adds \(q+1\) to the block \(J\). Every partition
of \(\{1,\ldots,q+1\}\) arises exactly once in this way, by removing
\(q+1\) from its block.
\end{proof}

\emph{Source.} \TBc{Com74}; the induction above proves the form used here.

For a finite list of face normals \(u_1,\ldots,u_m\in\R^n\), thresholds
\(a_1,\ldots,a_m\in\R\), and a width \(h>0\), put
\[
 q_\alpha(x)\coloneqq(u_\alpha\cdot x-a_\alpha)/h,\qquad
 w_\alpha(x)\coloneqq e^{q_\alpha(x)},\qquad T(x)\coloneqq\sum_\alpha w_\alpha(x),
\]
\[
 F(x)\coloneqq\frac1{1+T(x)},\qquad
 \pi_\alpha(x)\coloneqq F(x)w_\alpha(x),\qquad
 \theta(x)\coloneqq1-F(x)=\sum_\alpha\pi_\alpha(x).
\]
Here \(F\) is the logistic approximation to the indicator of the
threshold set. The baseline outcome has weight one and probability
\(F\); face \(\alpha\) has probability \(\pi_\alpha\). These
probabilities describe the smoothing formula, not the original event.
One and two differentiations give
\begin{align*}
 DF(x)[z]&=-h^{-1}\sum_\alpha F\pi_\alpha(u_\alpha\cdot z),\\
 D^2F(x)[z_1,z_2]
 &=h^{-2}\left\{2\sum_{\alpha,\beta}F\pi_\alpha\pi_\beta
       (u_\alpha\cdot z_1)(u_\beta\cdot z_2)
       -\sum_\alpha F\pi_\alpha
       (u_\alpha\cdot z_1)(u_\alpha\cdot z_2)\right\}.
\end{align*}
All weights are evaluated at \(x\). In the first formula their sum is
\(F\theta\), which is small away from the smoothed boundary.
The next result retains this factor at every derivative order.

\begin{fact}[Derivatives of the logistic approximation]\label{tb:logistic}
With the notation above, let \(q\ge1\).
For directions \(z_1,\ldots,z_q\), with the blocks \(J_1,\ldots,J_{|\mathcal P|}\)
of \(\mathcal P\) listed in any fixed order,
\[
 D^qF(x)[z_1,\ldots,z_q]
 =h^{-q}\sum_{\mathcal P\in\Pi_q}(-1)^{|\mathcal P|}|\mathcal P|!
 \sum_{\alpha_1,\ldots,\alpha_{|\mathcal P|}}
 F(x)\prod_{b=1}^{|\mathcal P|}\pi_{\alpha_b}(x)
 \prod_{b=1}^{|\mathcal P|}\prod_{a\in J_b}(u_{\alpha_b}\cdot z_a),
\]
Moreover, \(\sum_{\mathcal P\in\Pi_q}|\mathcal P|!\le q!\,2^{q-1}\), and hence
\[
 |D^qF(x)[z_1,\ldots,z_q]|
 \le q!\,2^{q-1}h^{-q}F(x)\theta(x)
       \prod_{a=1}^q\max_\alpha|u_\alpha\cdot z_a|.
\]
\end{fact}

\begin{proof}
Apply Fact~\ref{tb:faa} with \(\phi(\tau)\coloneqq (1+\tau)^{-1}\) and \(g=T\). Then
\(\phi^{(k)}(\tau)=(-1)^kk!\,(1+\tau)^{-k-1}\), and since each \(w_\alpha\) is
an exponential of an affine function,
\(D^lT(x)[z_1,\ldots,z_l]=h^{-l}\sum_\alpha w_\alpha(x)\prod_{a\le l}
(u_\alpha\cdot z_a)\). Substituting, and using
\((1+T)^{-|\mathcal P|-1}\prod_bw_{\alpha_b}=F\prod_b\pi_{\alpha_b}\) and
\(\sum_b|J_b|=q\), gives the formula. The number
\(\sum_{\mathcal P}|\mathcal P|!\) counts ordered partitions
\((J_1,\ldots,J_k)\) of \(\{1,\ldots,q\}\). Writing the elements of
\(J_1\), then of \(J_2\), and so on, in increasing order inside each block,
gives a permutation together with the set of positions at which a new block
starts, which is a subset of the \(q-1\) gaps between consecutive entries.
Every ordered partition arises from at least one of these
\(q!\,2^{q-1}\) pairs. For a partition with \(k\ge1\) blocks, summing the probability
weights gives \(F\theta^k\le F\theta\). Each product of projections
is bounded by \(\prod_a\max_\alpha|u_\alpha\cdot z_a|\).
These two observations prove the stated envelope without a factor
depending on the dimension.
\end{proof}

The early estimate uses this formula for orders three and four. After
the two Gaussian integrations by parts in the late estimate, the explicit
first and second derivatives above suffice. In both cases the factor
\(F\theta\) connects the derivative estimate to concentration near
the boundary.

\usedin{the general derivative expansion and local envelope in
Subsection~\PN{appx-sec:mixture}{B.2}, namely
\PE{appx-eq:partition-derivative}{B.1}--\PE{appx-eq:local-envelope}{B.3};
the explicit first and second derivatives in
\PR{Lemma}{appx-an:lem:restoration}{C.2}.}
\emph{Source.} The formula follows from the partition chain rule in
Fact~\ref{tb:faa}; the proof above is self-contained.

\subsection{Gaussian identities}\label{tb:sec:gauss}

The mean estimate separates a normal coordinate from the tangential
coordinates of a face. Orthogonality makes these Gaussian components
independent; a second elementary identity controls the density change
when we move outward from the face.

\begin{fact}[Gaussian density in orthogonal coordinates]\label{tb:gauss-coord}
(a) Let \(\R^n=L\oplus L^\perp\) be an orthogonal decomposition. Then
\(\varphi_n(l+w)=\varphi_L(l)\varphi_{L^\perp}(w)\) for \(l\in L\) and
\(w\in L^\perp\), so the projections of a standard Gaussian vector onto
\(L\) and \(L^\perp\) are independent standard Gaussian vectors. In
particular, for a unit vector \(u\), \(a\in\R\), and \(y\perp u\),
\(\varphi_n(au+y)=\varphi_1(a)\varphi_{n-1}(y)\).

(b) For \(x,v\in\R^n\),
\(\varphi_n(x+v)=\varphi_n(x)\exp\{-x\cdot v-\tfrac12\norm v^2\}\).
\end{fact}

\begin{proof}
Both follow from \(\norm{l+w}^2=\norm l^2+\norm w^2\) and from expanding
\(\norm{x+v}^2\).
\end{proof}

\usedin{the normal--tangential decomposition in
\PR{Lemma}{appx-an:lem:mean}{C.3}; the facet density factorization
and density comparison along normal tubes in
\PR{Lemma}{appx-an:lem:surface}{C.5}.}

Gaussian integration by parts replaces two directional derivatives by
a centered quadratic function of the Gaussian vector. The factor
\(\sigma^{-2}\) below records the Gaussian noise scale exactly.

\begin{fact}[Gaussian integration by parts]\label{tb:stein}
Let \(G\sim N(0,I_n)\). In (a), assume \(f\in C^1(\R^n)\) and
that \(f\) and its first derivatives have polynomial growth.
In (b), assume in addition that \(f\in C^2\) and its second
derivatives have polynomial growth. Then
\begin{enumerate}
\item[(a)] \(\E[Gf(G)]=\E\nabla f(G)\);
\item[(b)] \(\E[(GG^\top-I_n)f(G)]=\E D^2f(G)\); equivalently, for \(z\in\R^n\),
\(\E D^2f(G)[z,z]=\E\bigl[\tr\{(GG^\top-I_n)zz^\top\}f(G)\bigr]\), where
\(\tr\{(GG^\top-I_n)zz^\top\}=(z\cdot G)^2-\norm z^2\);
\item[(c)] if \(s\ge0\) is an integer, \(v\in\R^n\), \(\sigma>0\),
\(U\coloneqq v+\sigma G\), and
\(F\in C^{s+2}(\R^n)\) has derivatives of polynomial growth, then for
\(z,z_1,\ldots,z_s\in\R^n\),
\[
 \E D^{s+2}F(U)[z,z,z_1,\ldots,z_s]
 =\sigma^{-2}\E\bigl[\tr\{(GG^\top-I_n)zz^\top\}\,
 D^sF(U)[z_1,\ldots,z_s]\bigr].
\]
\end{enumerate}
\end{fact}

\begin{proof}
For (a) in one coordinate, \(x\varphi_1(x)=-\varphi_1'(x)\), so integration by
parts gives
\[\int xf(x)\varphi_1(x)\,dx=\int f'(x)\varphi_1(x)\,dx,\]
and the boundary terms vanish because of the growth condition. The other coordinates
are handled by Fact~\ref{tb:fubini}. For (b), apply (a) to \(\partial_if\) and
to \(G_kf\): \(\E[G_iG_kf(G)]=\E\partial_i(G_kf)(G)=\delta_{ik}\E f(G)
+\E[G_k\partial_if(G)]=\delta_{ik}\E f(G)+\E\partial_{ik}f(G)\). For (c), the
function \(f(g)\coloneqq D^sF(v+\sigma g)[z_1,\ldots,z_s]\) has
\(D^2f(g)[z,z]=\sigma^2D^{s+2}F(v+\sigma g)[z,z,z_1,\ldots,z_s]\); apply (b).
\end{proof}

\usedin{Subsection~\PN{sec:outline-late}{2.4} and
\PR{Lemma}{appx-an:lem:restoration}{C.2}, conditional on \(\mathcal H\)
with \(v=V_0+v_i\) and \(z=Z_i\); the identity
\(\E[Z\cdot\nabla\Psi(Z)]=\E\tr(\Sigma D^2\Psi(Z))\)
in \PR{Lemma}{appx-lem:endpoint-gap}{B.3}.}
The identity also holds conditionally when the base and directions are
measurable with respect to a sigma-field independent of \(G\), provided
the conditional growth assumptions hold. Averaging that identity further
requires integrability of both sides. In the paper the fixed width
\(h>0\), bounded face components, and conditional covariance bound
supply this integrability.

\emph{Source.} Stein~\TBc{Ste81}; the proof above is complete.

Changing the Gaussian variance contributes a Hessian trace to the path
derivative. We give the formula for a general covariance and then use
its standard-Gaussian form in the support coordinates.

\begin{fact}[Heat equation for the Gaussian component]\label{tb:heat}
Let \(Z\sim N(0,\Sigma)\) in \(\R^p\), where \(\Sigma\) may be singular, and
let \(f\in C^2(\R^p)\) be bounded with bounded first and second derivatives.
Then, for \(x\in\R^p\) and \(t>0\),
\[
 \frac d{dt}\E f(x+\sqrt t Z)
 =\tfrac12\E\tr\{\Sigma D^2f(x+\sqrt t Z)\}.
\]
\end{fact}

\begin{proof}
Write \(Z=\Gamma G\) with \(\Gamma\Gamma^\top=\Sigma\) and \(G\) standard
Gaussian. By Fact~\ref{tb:taylor}(a), which applies because the difference
quotients near \(t_0>0\) are bounded by
\(\norm\nabla f_\infty\norm{\Gamma G}/(2\sqrt{t_0/2})\)
when the times lie in \([t_0/2,3t_0/2]\), the derivative equals
\((2\sqrt t)^{-1}\E[G\cdot\Gamma^\top\nabla f(x+\sqrt t\Gamma G)]\).
Fact~\ref{tb:stein}(a) applied to the function \(g\mapsto
\Gamma^\top\nabla f(x+\sqrt t\Gamma g)\) turns this into
\((2\sqrt t)^{-1}\sqrt t\,\E\tr\{\Gamma^\top D^2f\,\Gamma\}
=\frac12\E\tr\{\Sigma D^2f\}\).
\end{proof}

In the support \(\R^r\), the covariance is \(I_r\), so this gives
\(\frac d{dt}\E_G F(v+\sqrt tG)=\frac12\E_G\tr D^2F(v+\sqrt tG)\).
The base \(v\) is held fixed for this differentiation. If its law also
depends on time, its change contributes separately.

\usedin{Subsection~\PN{sec:outline-path}{2.2} and the derivation of
\PE{appx-an:signed-hessian}{C.6} in
Subsection~\PN{appx-an:sec:path}{C.2}. The Gaussian component supplies
the covariance term that cancels the constant term of the Hessian
expansion.}

The surface estimate needs the size of a few normal projections,
not the Euclidean length of the full Gaussian vector. The following
bound explains why only a logarithm of the number of faces enters.

\begin{fact}[Maxima of one- and two-dimensional Gaussian projections]
\label{tb:projection-moments}
Let \(u_1,\ldots,u_m\) be unit vectors in \(\R^r\), with \(m\ge1\).
For linearly independent \(u_j,u_k\), let \(\Pi_{jk}\) be orthogonal
projection onto their span, and put
\[
 \mathcal M(g)\coloneqq \max\left\{\max_j|u_j\cdot g|,
   \max_{j,k:\,u_j,u_k\text{ independent}}\norm{\Pi_{jk}g}\right\},
 \qquad R(g)\coloneqq 1+\mathcal M(g).
\]
An empty second maximum is omitted. Then \(R\) is \(1\)-Lipschitz.
For \(G\sim N(0,I_r)\) and \(q\in\{2,3\}\),
\[
 \E R(G)^q\le C_q\{\log(em)\}^{q/2}.
\]
\end{fact}

\begin{proof}
Every projection is a contraction, so each function in the maximum
is \(1\)-Lipschitz, as is their maximum. A scalar standard normal
satisfies \(\Pp(|Z|>s)\le2e^{-s^2/2}\); the length of a
two-dimensional standard normal satisfies
\(\Pp(\norm{G_2}>s)=e^{-s^2/2}\). A union bound gives
\[
 \Pp\{\mathcal M(G)>s\}\le\min\{1,2(m+m^2)e^{-s^2/2}\}.
\]
The projections need not be independent. Integrate
\(q\int_0^\infty s^{q-1}\Pp\{\mathcal M(G)>s\}\,ds\), splitting
at \(\sqrt{8\log(em)}\), and use
\((1+x)^q\le2^{q-1}(1+x^q)\).
\end{proof}

\usedin{\PE{appx-an:projection-moments}{C.28} in
\PR{Lemma}{appx-an:lem:surface}{C.5}. For the full face list \(m=2p\), and also for its subsystems, these moments
are of orders \(L\) and \(L^{3/2}\), where \(L=\log(ep)\).}

\subsection{Matrix and entropy inequalities}\label{tb:sec:matrix}

Both time ranges use the same observation: the summed second moments
control a pair of projections even when other random coefficients
remain inside the expectation.

\begin{fact}[Conditional covariance contraction]\label{tb:covariance}
Let \(Z_1,\ldots,Z_N\in\R^r\) have finite second moments, and let
\(\mathcal H\) be a sigma-field such that
\(\sum_i\E[Z_iZ_i^\top\mid\mathcal H]\preceq I_r\) almost surely.
For \(\mathcal H\)-measurable unit vectors \(u,v\),
\[
 \sum_i\E[|(u\cdot Z_i)(v\cdot Z_i)|\mid\mathcal H]\le1.
\]
If real random coefficients satisfy \(|c_i|\le a\) almost surely,
where \(a\ge0\) is constant, then
\[
 \left\|\sum_i\E[c_iZ_iZ_i^\top\mid\mathcal H]\right\|_\op\le a.
\]
No independence of the coefficients and vectors is required.
\end{fact}

\begin{proof}
Conditional Cauchy--Schwarz, followed by Cauchy--Schwarz over the
index \(i\), bounds the first sum by
\[
 \left(\sum_i\E[(u\cdot Z_i)^2\mid\mathcal H]\right)^{1/2}
 \left(\sum_i\E[(v\cdot Z_i)^2\mid\mathcal H]\right)^{1/2}\le1.
\]
For the second claim, each unit vector \(u\) gives
\(\bigl|\sum_i\E[c_i(u\cdot Z_i)^2\mid\mathcal H]\bigr|\le a\).
The matrix is symmetric, so taking the supremum of its absolute
quadratic forms proves the operator-norm bound.
\end{proof}

\usedin{\PE{appx-an:pair-contraction}{C.2} and the early estimates in
Subsection~\PN{appx-an:sec:early}{C.3}; the coefficient matrices in
\PR{Lemma}{appx-an:lem:restoration}{C.2}. Bounded face projections
control the remaining directions by powers of \(\delta\).}

After conditioning, the coefficient matrix has bounded operator norm.
The following duality transfers that bound to the Gaussian matrix, whose
nuclear norm we estimate using its positive and negative parts.

\begin{fact}[Duality between the operator and nuclear norms]\label{tb:nuclear}
(a) For every square matrix \(\Xi\) and every symmetric matrix \(M\),
\(|\tr(\Xi M)|\le\norm\Xi_\op\norm M_*\). Moreover
\(\norm M_*=\max\{\tr(\Xi M):\norm\Xi_\op\le1\}\), so that \(\norm\cdot_*\) is
a norm on symmetric matrices.

(b) If \(M=P-N\) with \(P\succeq0\) and \(N\succeq0\), then
\(\norm M_*\le\tr P+\tr N=2\tr P-\tr M\).
\end{fact}

\begin{proof}
For (a), let \(M=\sum_i\mu_ie_ie_i^\top\) be a spectral decomposition. Then
\(\tr(\Xi M)=\sum_i\mu_i\,e_i^\top\Xi e_i\) and \(|e_i^\top\Xi e_i|\le
\norm\Xi_\op\), which gives the inequality. Equality is attained at
\(\Xi\coloneqq \sum_i\operatorname{sgn}(\mu_i)e_ie_i^\top\). A supremum of linear
functionals satisfies the triangle inequality. For (b), the norm of a
positive semidefinite matrix is its trace, so the triangle inequality gives
\(\norm M_*\le\tr P+\tr N\), and \(\tr N=\tr P-\tr M\).
\end{proof}

\usedin{the final contraction in \PR{Lemma}{appx-an:lem:restoration}{C.2};
the reduction \PE{appx-an:matrix-trace-reduction}{C.22} in
\PR{Lemma}{appx-an:lem:matrix}{C.4}, with
\(P=p_\alpha\mu_\alpha\mu_\alpha^\top\) and
\(N=p_\alpha(I_r-K_\alpha)\).}

An individual logistic weight can have arbitrarily small mass, so its
logarithmic bound should be summed before the mass is estimated. The
log-sum inequality performs that summation; its integral version then
allows us to average over the independent base.

\begin{fact}[Log-sum inequality and a concave function]\label{tb:logsum}
(a) For positive numbers \(y_1,\ldots,y_l\) and \(z_1,\ldots,z_l\) with sums
\(y\) and \(z\),
\(\sum_jy_j\log(z_j/y_j)\le y\log(z/y)\).

(b) The function \((x,y)\mapsto x\log(y/x)\) is concave on \((0,\infty)^2\).
Hence, for positive integrable random variables \(X,Y\),
\(\E[X\log(Y/X)]\le\E X\,\log(\E Y/\E X)\).
The expectation on the left may equal \(-\infty\); its positive
part is integrable because \(x\log(y/x)\le y/e\).
\end{fact}

\begin{proof}
For (a), Jensen's inequality for the concave function \(\log\), with the
weights \(y_j/y\), gives \(\sum_j(y_j/y)\log(z_j/y_j)\le\log\sum_j(y_j/y)(z_j/y_j)
=\log(z/y)\). For (b), the function \(f(x,y)\coloneqq x\log y-x\log x\) has Hessian
\(\begin{pmatrix}-1/x&1/y\\1/y&-x/y^2\end{pmatrix}\), whose determinant is
zero and whose trace is negative, so the Hessian is negative semidefinite.
For the expectation bound, put \(\mu\coloneqq\E X\) and
\(\eta\coloneqq\E Y\). The tangent-plane inequality for this concave
function gives
\[
 f(X,Y)\le f(\mu,\eta)
   +f_x(\mu,\eta)(X-\mu)+f_y(\mu,\eta)(Y-\eta).
\]
The right side is integrable. Taking expectations proves the claim,
including when the negative part of the left side has infinite mean.
\end{proof}

\usedin{the reference-mass estimate in
\PR{Lemma}{appx-an:lem:matrix}{C.4}: first sum over tuples sharing
their first face, then average over \(V_0\). Both the reference mass
and the logistic mass contain the factor \(h\sqrt L\), which cancels
inside the logarithm.}
\emph{Source.} \TBc{CT06}, Theorem~2.7.1, for (a).

\subsection{Log-concavity}\label{tb:sec:logconcave}

A function \(f:\R^n\to[0,\infty)\) is \emph{log-concave} if
\(f((1-\alpha)x+\alpha y)\ge f(x)^{1-\alpha}f(y)^\alpha\) for all \(x,y\) and \(\alpha\in[0,1]\);
when \(f>0\) this says that \(\log f\) is concave. At the endpoints
\(\alpha=0,1\) the inequality means equality, without a convention
for \(0^0\). We need the fact that integrating out coordinates
preserves log-concavity: it will turn a Gaussian convolution into a
concave log normalizing constant.

\begin{fact}[Pr\'ekopa--Leindler inequality and Pr\'ekopa's theorem]
\label{tb:prekopa}
(a) Let \(0<\alpha<1\) and let \(f,g,h:\R^n\to[0,\infty)\) be bounded,
continuous, and integrable, with
\(h((1-\alpha)x+\alpha y)\ge f(x)^{1-\alpha}g(y)^\alpha\) for all \(x,y\). Then
\(\int h\ge(\int f)^{1-\alpha}(\int g)^\alpha\).

(b) Let \(f:\R^a\times\R^b\to[0,\infty)\) be bounded, continuous, and
log-concave, with \(f(z,\cdot)\) integrable for every \(z\). Then
\(z\mapsto\int f(z,y)\,dy\) is log-concave on \(\R^a\).
\end{fact}

\begin{proof}
Part (a) is the Pr\'ekopa--Leindler inequality; we use the standard
integrable-function theorem in \TBc{Gar02}, Theorem~7.1.
To deduce (b), fix \(z_0,z_1\) and \(0<\alpha<1\), and put
\(z_\alpha\coloneqq(1-\alpha)z_0+\alpha z_1\).
Log-concavity gives
\[
 f(z_\alpha,(1-\alpha)y_0+\alpha y_1)
 \ge f(z_0,y_0)^{1-\alpha}f(z_1,y_1)^\alpha.
\]
The three sections are bounded, continuous, and integrable by hypothesis,
so (a) applies to them. Its conclusion is exactly the log-concavity of
\(z\mapsto\int f(z,y)\,dy\). The endpoint cases are immediate.
\end{proof}

\emph{Source.} Pr\'ekopa~\TBc{Pre71,Pre73} and Leindler~\TBc{Lei72}.
The marginalization result is Theorem~6 of \TBc{Pre73};
\TBc{Gar02}, Section~7, treats the integral inequality.

We now apply marginal log-concavity to a Gaussian convolution. Its
logarithmic Hessian is the covariance of the tilted law minus the
identity, so concavity gives the matrix inequality needed in the proof.

\begin{fact}[Gaussian laws tilted by a log-concave weight]\label{tb:tilt}
Let \(w:\R^r\to[0,\infty)\) be bounded, continuous, and log-concave, with
\(p\coloneqq \E w(G)>0\) for \(G\sim N(0,I_r)\). Let \(\mu\) and \(K\) be the mean and
covariance matrix of the probability law with density \(w(g)\varphi_r(g)/p\).
Then \(K\preceq I_r\), and
\(\E[(GG^\top-I_r)w(G)]=p(K-I_r+\mu\mu^\top)\).
\end{fact}

\begin{proof}
The function \((z,g)\mapsto e^{-\norm{g-z}^2/2}w(g)\) is log-concave, since
\(\norm{g-z}^2\) is jointly convex in \((z,g)\), and it is bounded and
continuous, with an integrable \(g\)-section for every \(z\).
The integral is positive because \(p>0\) and the Gaussian kernel is
strictly positive. By Fact~\ref{tb:prekopa}(b), \(\ell(z)\coloneqq \log\int e^{-\norm{g-z}^2/2}
w(g)\,dg\) is concave. Since \(\norm{g-z}^2=\norm g^2-2z\cdot g+\norm z^2\),
\[
 \ell(z)=\text{const}-\tfrac12\norm z^2+\log M(z),\qquad
 M(z)\coloneqq \int e^{z\cdot g}w(g)\varphi_r(g)\,dg=p\,\E_\nu e^{z\cdot g},
\]
where \(\nu\) is the tilted law. The function \(M\) is finite and smooth: on every bounded set of
\(z\)'s, each derivative of the integrand is bounded by a constant
times \((1+\norm g)^k e^{c\norm g}\varphi_r(g)\), an integrable function.
Thus differentiation under the integral is valid, and the Hessian of \(\log M\) at \(0\) is the
covariance \(K\) of \(\nu\). A concave function of class \(C^2\) has a
negative semidefinite Hessian, so \(K-I_r\preceq0\). The identity is the
decomposition of the second moment \(\E_\nu gg^\top=K+\mu\mu^\top\).
\end{proof}

\usedin{the Gaussian-tilt calculation in
\PR{Lemma}{appx-an:lem:matrix}{C.4}, where \(w=w_\alpha\) and
\(K=K_\alpha\), and its explanation in
Subsection~\PN{sec:outline-late}{2.4}.}
\emph{Source.} The covariance contraction is a classical consequence of
Pr\'ekopa's marginalization theorem~\TBc{Pre73}, Theorem~6, as the
preceding proof shows. For smooth positive weights, it also follows
from the Brascamp--Lieb variance inequality~\TBc{BL76}, applied to
linear functions under the potential \(\norm g^2/2-\log w(g)\), whose
Hessian is at least \(I_r\). The same convolution argument covers
indicators of convex sets by the measurable form of Pr\'ekopa's
theorem. This convex-cell specialization is also used in \TBc{10}.

It remains to check that every weight produced by the derivative
expansion meets the preceding hypothesis. Its logarithm has a particularly
simple form: an affine function minus a multiple of a log-sum-exponential.

\begin{fact}[Log-concavity of the logistic weights]\label{tb:lc-weights}
With the notation of Fact~\ref{tb:logistic}, for every \(k\ge0\) and every
tuple \((\alpha_1,\ldots,\alpha_k)\), the function
\(F\pi_{\alpha_1}\cdots\pi_{\alpha_k}=\exp\{\sum_{a\le k}q_{\alpha_a}\}
(1+T)^{-k-1}\), where \(q_\alpha\coloneqq (u_\alpha\cdot x-a_\alpha)/h\), is
log-concave on \(\R^n\), and so is its composition
\(g\mapsto F\pi_{\alpha_1}\cdots\pi_{\alpha_k}(v+\sigma g)\) with any affine map.
\end{fact}

\begin{proof}
The logarithm is a sum of the affine functions \(q_{\alpha_a}\) and
\(-(k+1)\log(1+\sum_\alpha e^{q_\alpha})\). The function
\(\log(1+\sum_\alpha e^{q_\alpha})\) is convex. To include the baseline
term explicitly, put \(q_0\coloneqq0\). H\"older's inequality gives
\[
 \sum_{\alpha=0}^m e^{(1-s)q_\alpha(x)+s q_\alpha(y)}
 \le\left(\sum_{\alpha=0}^m e^{q_\alpha(x)}\right)^{1-s}
     \left(\sum_{\alpha=0}^m e^{q_\alpha(y)}\right)^s,
 \qquad 0<s<1.
\]
Taking logarithms proves the assertion. Composition
with an affine map preserves concavity.
\end{proof}

\usedin{the application of covariance contraction to every
\(w_\alpha=F\prod_a\pi_{\alpha_a}\) in
\PR{Lemma}{appx-an:lem:matrix}{C.4}.}
\emph{Source.} Convexity of log-sum-exp is standard; see
Boyd and Vandenberghe~\TBc{BV04}, Section~3.1.5, pp.~72--74.
The baseline term \(1=e^{q_0}\), with \(q_0=0\), is included in that sum.

\subsection{Polyhedra: faces, normal cones, and nearest points}\label{tb:sec:poly}

The surface estimate treats a fixed polyhedron geometrically. Unlike
the smoothing formula, its boundary integral depends only on the set,
so repeated and redundant inequalities may now be removed.
A \emph{polyhedron} is a set
\[
 Q\coloneqq \bigcap_{j=1}^m\{x\in\R^n:u_j\cdot x\le a_j\},
\]
with unit vectors \(u_j\), called \emph{face normals}, and numbers
\(a_j\), called \emph{offsets}. Each face normal points outward.
The polyhedron is closed and convex. It is
\emph{full-dimensional} if it has an interior point, and \emph{proper} if it
is neither empty nor \(\R^n\). The description is \emph{irredundant} if
deleting any one inequality changes the set; deleting redundant or repeated
inequalities one at a time turns every description into an irredundant one
of the same set. For \(j\le m\), put \(H_j\coloneqq \{x:u_j\cdot x=a_j\}\) and
\(F_j\coloneqq Q\cap H_j\). A \emph{face} of \(Q\) is \(Q\) itself or a set
\(Q\cap\{u\cdot x=c\}\) with \(u\ne0\) and \(Q\subseteq\{u\cdot x\le c\}\).
A \emph{facet} is a face of dimension \(n-1\), and a \emph{ridge}, when \(n\ge2\), is a face
of dimension \(n-2\). In the plane the facets are the sides of a polygon
and the ridges its corners, in \(\R^3\) the facets are the faces of a solid
and the ridges its edges, and on the line the facets are the endpoints of an
interval, which has no ridges. We use these standard properties of faces:
the relative interiors of the faces of \(Q\) are pairwise disjoint and cover
\(Q\); a face of a face of \(Q\) is a face of \(Q\); and a face \(G'\)
of a face \(G\) with \(G'\ne G\) lies in the relative boundary
\(G\setminus\relint G\) of \(G\) (\TBc{Roc70}, Section~18; \TBc{Zie95}, Chapter~2).

The next fact identifies the pieces that contribute to the boundary
calculation. In particular, every ridge belongs to exactly two facets;
this allows us to combine their coefficients before taking absolute values.

\begin{fact}[Facets and ridges of a polyhedron]\label{tb:facets}
Let \(Q\) be a proper, full-dimensional polyhedron in \(\R^n\) with an
irredundant description. Statements about ridges are omitted when \(n=1\).
\begin{enumerate}
\item[(a)] The face normals \(u_1,\ldots,u_m\) are pairwise distinct, each
\(F_j\) is a facet, \(Q\) has no other facets, and
\(\partial Q=F_1\cup\cdots\cup F_m\).
\item[(b)] For \(j\ne k\), \(F_j\cap F_k\subseteq H_j\cap H_k\), and
\(H_j\cap H_k\) is empty if \(u_k=-u_j\), and an affine subspace of
dimension \(n-2\) otherwise. In particular \(F_j\cap F_k\) lies in a
finite union of affine subspaces of dimension at most \(n-2\).
\item[(c)] Each \(F_j\) is a full-dimensional polyhedron in the hyperplane
\(H_j\). Write the points of \(H_j\) as \(a_ju_j+y\) with \(y\perp u_j\), and,
for \(k\ne j\) with \(u_k\ne\pm u_j\), put
\(\alpha_{jk}\coloneqq \arccos(u_j\cdot u_k)\in(0,\pi)\),
\[
 \nu_{k\mid j}\coloneqq \frac{u_k-(u_j\cdot u_k)u_j}{\sin\alpha_{jk}},\qquad
 d_{k\mid j}\coloneqq \frac{a_k-(u_j\cdot u_k)a_j}{\sin\alpha_{jk}} .
\]
Then \(F_j\) is described in the coordinate \(y\) by the inequalities
\(\nu_{k\mid j}\cdot y\le d_{k\mid j}\), and its facets are exactly the
ridges \(E_{jk}\coloneqq F_j\cap F_k\) of \(Q\) of dimension \(n-2\); the outward
unit normal of \(F_j\) at \(E_{jk}\) is \(\nu_{k\mid j}\).
\item[(d)] Every ridge of \(Q\) lies in exactly two facets, so the ridges
correspond to the unordered pairs \(\{j,k\}\) with
\(\dim(F_j\cap F_k)=n-2\); for such a pair \(u_k\ne\pm u_j\), and the
only inequalities active at a point of \(\relint E_{jk}\) are \(j\) and \(k\).
\end{enumerate}
\end{fact}

\begin{proof}
(a) Fix \(j\). Irredundancy gives a point \(x_j\) with \(u_j\cdot x_j>a_j\)
and \(u_k\cdot x_j\le a_k\) for \(k\ne j\). Let \(x_0\) be an interior point of
\(Q\), and let \(x_*\) be the point where the segment from \(x_0\) to \(x_j\)
crosses \(H_j\). For \(k\ne j\), \(u_k\cdot x_*<a_k\), because \(x_*\) is a
convex combination of \(x_0\), where the inequality is strict, with positive
weight, and of \(x_j\), where it holds weakly. Hence a neighborhood of
\(x_*\) in \(H_j\) lies in \(F_j\), and \(\dim F_j=n-1\). Two inequalities
with the same normal cannot both be irredundant, since the one with the
larger offset is implied by the other. A point of \(Q\) at
which no inequality is active is interior, and at a point where
inequality \(j\) is active the points \(x+\varepsilon u_j\) lie outside \(Q\),
so \(\partial Q=\bigcup_jF_j\). If \(G=Q\cap\{u\cdot x=c\}\) is a facet and
\(x\in\relint G\), some inequality \(j\) is active at \(x\); on \(G\) the
affine function \(u_j\cdot x-a_j\le0\) attains its maximum at a relative
interior point and is therefore zero, so \(G\subseteq H_j\), and since \(G\) has
dimension \(n-1\), \(G=Q\cap H_j=F_j\).

(b) If \(u_k=-u_j\) and \(H_j\cap H_k\ne\emptyset\), then \(Q\) would lie in
the hyperplane \(H_j\), contradicting full dimension. If \(u_k\ne\pm u_j\),
the two equations are independent.

(c) Suppose \(n\ge2\). The set \(F_j\) is the polyhedron in \(H_j\) cut out by the remaining
inequalities, which take the stated form in the coordinate \(y\) because
\(u_k\cdot(a_ju_j+y)=a_ju_j\cdot u_k+(u_k-(u_j\cdot u_k)u_j)\cdot y\); an
inequality with \(u_k=-u_j\) holds automatically on \(H_j\) since
\(F_j\ne\emptyset\). After removing inequalities redundant within \(H_j\), part (a)
applied there shows that the facets of \(F_j\) are
the faces of \(F_j\) of dimension \(n-2\), hence faces of \(Q\) of that
dimension, that is, ridges of \(Q\) in \(F_j\); each equals
\(F_j\cap H_k=F_j\cap F_k\) for the inequality \(k\) that defines it.

(d) Let \(E\) be a ridge and \(y\in\relint E\), and let \(J\) be the set of
inequalities active at \(y\). On \(E\) each affine function
\(u_i\cdot x-a_i\), \(i\in J\), is at most zero and vanishes at the relative
interior point \(y\), so it vanishes on \(E\); thus \(E\subseteq F_i\) for
\(i\in J\) and \(\dim\operatorname{span}\{u_i:i\in J\}\le2\). A single active
inequality is impossible, since then \(y\) would lie in the relative interior
of \(F_j\) and the only face of \(F_j\) containing it is \(F_j\) itself, of
dimension \(n-1\). Opposite face normals \(u_i=-u_j\) are excluded by (b). Hence
the active face normals span a plane \(\Pi\). Near \(y\), \(Q\) coincides with
\(y\) plus the product of \(\Pi^\perp\) and a planar cone \(K\subseteq\Pi\)
with nonempty interior, defined by the active inequalities and contained in
a wedge of opening \(\pi-\alpha_{jk}<\pi\); this cone has exactly two
boundary rays. Each active inequality \(i\) has \(\dim F_i=n-1\), so
\(F_i\) meets every neighborhood of \(y\) in a set of dimension \(n-1\), and
the line \(\{u_i\cdot v=0\}\) must therefore carry a boundary ray of \(K\).
With distinct face normals this leaves exactly two active inequalities.
\end{proof}

\usedin{the facet--ridge decomposition leading to
\PE{appx-an:surface-identity}{C.30} and the pairing of the two terms
attached to each ridge in \PR{Lemma}{appx-an:lem:surface}{C.5}.}
\emph{Source.} These are standard facts about faces of polyhedra; see
\TBc{Roc70}, Sections~18 and~19, and \TBc{Zie95}, Chapter~2.
The planar-cone argument above explains why precisely two facets
meet at a ridge, including for unbounded polyhedra.

To attach a tube to a face, we must know which outward directions
keep that face as the nearest part of the polyhedron. These directions
form its normal cone, determined by the active inequalities.

\begin{fact}[Normal cones]\label{tb:normalcone}
Let \(Q\) be a polyhedron as in Fact~\ref{tb:facets}. For \(x\in Q\) let
\(J(x)\coloneqq \{j:u_j\cdot x=a_j\}\), and let
\[
 N_Q(x)\coloneqq \{v\in\R^n:v\cdot(q-x)\le0\text{ for all }q\in Q\}
\]
be the \emph{normal cone} of \(Q\) at \(x\). Then
\(N_Q(x)=\{\sum_{j\in J(x)}c_ju_j:c_j\ge0\}\). In particular
\(N_Q(x)=\{0\}\) at interior points, \(N_Q(x)=\{cu_j:c\ge0\}\) for
\(x\in\relint F_j\), and \(N_Q(x)=\{c_ju_j+c_ku_k:c_j,c_k\ge0\}\), a sector of
angle \(\alpha_{jk}\) in the plane spanned by \(u_j\) and \(u_k\), for
\(x\in\relint E_{jk}\).
\end{fact}

\begin{proof}
Put \(T\coloneqq \{d:u_j\cdot d\le0\text{ for }j\in J(x)\}\). A vector \(v\) lies in
\(N_Q(x)\) if and only if \(v\cdot d\le0\) for all \(d\in T\): if \(d\in T\)
then \(x+\varepsilon d\in Q\) for small \(\varepsilon>0\), because the
inactive inequalities are strict at \(x\), and conversely \(q-x\in T\) for
\(q\in Q\). By Farkas's lemma, which says that \(v\cdot d\le0\) for all
\(d\) with \(u_j\cdot d\le0\) (\(j\in J\)) if and only if \(v\) is a
nonnegative combination of the \(u_j\), \(j\in J\), this is the stated
cone. For \(x\in\relint F_j\), \(J(x)=\{j\}\), since an index \(k\ne j\) in
\(J(x)\) would put \(x\) in \(F_j\cap F_k\), a face of \(F_j\) of dimension
at most \(n-2\), and so in the relative boundary of \(F_j\). For
\(x\in\relint E_{jk}\), \(J(x)=\{j,k\}\) by Fact~\ref{tb:facets}(d).
\end{proof}

\usedin{the choice of the normal interval over a facet and the normal
sector over a ridge in \PR{Lemma}{appx-an:lem:surface}{C.5}.}
\emph{Source.} \TBc{Roc70}, Sections~22 and~23, for Farkas's lemma and normal
cones; \TBc{RW98}, Theorem~6.46.

The nearest-point characterization makes the outward tubes disjoint.
It applies to any nonempty closed convex set; for such a set we use
\(N_Q(x)=\{v:v\cdot(q-x)\le0\text{ for all }q\in Q\}\).

\begin{fact}[Nearest point]\label{tb:projection}
Let \(Q\subseteq\R^n\) be nonempty, closed, and convex. For each
\(y\in\R^n\) there is a unique point \(\prj_Q(y)\in Q\) with
\(\norm{y-\prj_Q(y)}=\dist(y,Q)\), and \(x=\prj_Q(y)\) if and only if
\(x\in Q\) and \(y-x\in N_Q(x)\). Consequently every \(y\in\R^n\) can be
written in exactly one way as \(y=x+v\) with \(x\in Q\) and \(v\in N_Q(x)\),
and then \(\norm v=\dist(y,Q)\). A direction \(v\) lies in \(N_Q(x)\) if and
only if \(x\) is the nearest point of \(x+s v\) in \(Q\) for every \(s>0\).
\end{fact}

\begin{proof}
For \(x,q\in Q\), \(\norm{y-q}^2=\norm{y-x}^2+\norm{x-q}^2-2(y-x)\cdot(q-x)\).
If \(y-x\in N_Q(x)\), the last term is nonnegative, so
\(\norm{y-q}^2\ge\norm{y-x}^2+\norm{x-q}^2>\norm{y-x}^2\) for \(q\ne x\); thus
\(x\) is the unique nearest point. Conversely, if \((y-x)\cdot(q-x)>0\) for
some \(q\in Q\), then moving from \(x\) toward \(q\) decreases the distance
to \(y\) for small steps, by the displayed identity applied to
\(x+\varepsilon(q-x)\in Q\), so \(x\) is not nearest. For existence, choose \(q_0\in Q\). A minimizing sequence may be
restricted to the closed ball centered at \(y\) with radius
\(\norm{y-q_0}+1\); it therefore has a convergent subsequence in
finite dimension. Closedness puts its limit in \(Q\), and continuity
of the norm makes that limit a minimizer. The last statement follows because \(N_Q(x)\) is a
cone.
\end{proof}

\usedin{disjointness of the variable-width tubes in
\PR{Lemma}{appx-an:lem:surface}{C.5}; each tube point has its nearest
point in the relative interior of exactly one face.}
\emph{Source.} \TBc{RW98}, Example~2.25, for existence and uniqueness,
and Proposition~6.17 for the characterization by the normal cone.

\subsection{Surface measure, Gaussian surface integrals, and the
divergence theorem}\label{tb:sec:surface}

We integrate each facet with its ordinary Euclidean area measure.
Hausdorff measure gives a common notation for facets and ridges in
all dimensions and shows why their lower-dimensional overlaps do not
affect a surface integral.

\begin{fact}[Hausdorff measure on flat pieces]\label{tb:hausdorff}
For \(0\le k\le n\), the \(k\)-dimensional Hausdorff measure \(\Haus^k\) on
\(\R^n\) is obtained from the covering formula
\[
 \Haus^k(S)\coloneqq \lim_{\epsilon\downarrow0}\inf\Bigl\{\sum_i\omega_k
 \Bigl(\tfrac12\operatorname{diam}C_i\Bigr)^k:S\subseteq\bigcup_iC_i,\
 \operatorname{diam}C_i\le\epsilon\Bigr\},
\]
where the infimum runs over countable covers and \(\omega_k\) is
the volume of the unit ball of \(\R^k\).
For \(k=0\), every nonempty covering set contributes one, including
a singleton. The formula defines an outer measure, whose restriction
to Borel sets is the measure used here. (a) On an
affine subspace of dimension \(k\), \(\Haus^k\) is \(k\)-dimensional Lebesgue
measure; in particular \(\Haus^0\) is counting measure and \(\Haus^n\) is
Lebesgue measure on \(\R^n\). (b) \(\Haus^k\) is invariant under isometries.
(c) A set contained in a finite union of affine subspaces of dimension less
than \(k\) has \(\Haus^k\)-measure zero.
\end{fact}

Part (c) follows from (a), since each such subspace lies in a
\(k\)-dimensional affine subspace in which it is a Lebesgue null set. For a
proper, full-dimensional polyhedron \(Q\) with an irredundant description,
and for a function \(f\ge0\) or \(f\) integrable on \(\partial Q\), the
\emph{surface integral} is
\[
 \int_{\partial Q}f\,d\Haus^{n-1}=\sum_{j=1}^m\int_{F_j}f\,d\Haus^{n-1},
\]
where on each facet \(\Haus^{n-1}\) is Lebesgue measure on the hyperplane
\(H_j\) by part~(a). By Fact~\ref{tb:facets}(b) and part~(c), distinct
facets overlap in sets of \(\Haus^{n-1}\)-measure zero, so the sum is the
integral of \(f\) over the set \(\partial Q\) with respect to
\(\Haus^{n-1}\). It depends on the set \(Q\) and not on its description,
which justifies deleting redundant inequalities before computing it.
The Gaussian surface integral \(\surf(Q)=\int_{\partial Q}\Hh\,\varphi_r\,
d\Haus^{r-1}\) of Subsection~\PN{sec:outline-late}{2.4} is the case
\(n=r\) and \(f=\Hh\varphi_r\), with \(\Hh(g)\coloneqq \norm g^2-r\).
The integral is signed; we take its absolute value only after summing
the facet integrals. We set \(\surf(\varnothing)=\surf(\R^r)\coloneqq0\).
Lower-dimensional polyhedra are excluded from this definition and are
handled as exceptional levels in Fact~\ref{tb:variation}.

\usedin{the definition of \(\surf\) in
Subsection~\PN{sec:outline-late}{2.4} and
\PR{Lemma}{appx-an:lem:surface}{C.5}.}
\emph{Source.} \TBc{EG15}, Chapter~2.

Integrating a tangential divergence over a facet converts its interior
integral into contributions from its ridges. We prove the needed divergence
formula by slicing, then use a cutoff to allow an unbounded polyhedron
under precisely the stated decay assumptions.

\begin{samepage}
\begin{fact}[Gauss--Green formula (divergence theorem) for a polyhedron]\label{tb:divergence}
Let \(Q\subseteq\R^n\) (\(n\ge1\)) be a proper, full-dimensional polyhedron
with an irredundant description, facets \(F_1,\ldots,F_m\), and unit
face normals \(u_1,\ldots,u_m\). Let \(X:\R^n\to\R^n\) be of class \(C^1\), with
\(\norm{X(x)}+|\operatorname{div}X(x)|\le C(1+\norm x)^k\varphi_n(x)\) for
some \(C\) and \(k\). Then
\[
 \int_Q\operatorname{div}X\,dx=\sum_{j=1}^m\int_{F_j}X\cdot u_j\,d\Haus^{n-1}.
\]
For \(n=1\) the facets are the endpoints of an interval, \(\Haus^0\) is
counting measure, and the identity is the fundamental theorem of calculus.
\end{fact}
\end{samepage}

\begin{proof}
We first prove the formula for a compactly supported \(C^1\) vector
field. For a unit vector \(e\) and a compactly supported
\(C^1\) scalar function \(f\), we claim that
\begin{equation}
 \int_Q\partial_ef\,dx=\sum_j(e\cdot u_j)\int_{F_j}f\,d\Haus^{n-1},
 \tag{\(*\)}
\end{equation}
since the theorem follows by taking \(e=e_i\) and \(f=X_i\) and summing over
\(i\). Write \(x=w+s e\) with \(w\in e^\perp\). By Fubini's theorem,
\(\int_Q\partial_ef\,dx=\int_{e^\perp}\int_{I(w)}\partial_sf(w+s e)\,ds\,dw\),
where \(I(w)\coloneqq \{s:w+s e\in Q\}\) is a closed interval, possibly empty, a point,
or unbounded, because \(Q\) is convex. Since \(f\) tends to zero along every
line, the inner integral equals \(f(w+s_+(w)e)-f(w+s_-(w)e)\), where
\(s_\pm(w)\) are the endpoints of a nonempty \(I(w)\).
An empty interval contributes zero, and a term with an infinite
endpoint is read as zero.

Consider the upper endpoints \(x=w+s_+(w)e\). If \(x\in\relint F_j\), then
only the inequality \(j\) is active at \(x\) (Fact~\ref{tb:normalcone}), so
\(x+s e\in Q\) for small \(s>0\) unless \(e\cdot u_j>0\), and \(x\) is an upper
endpoint if and only if \(e\cdot u_j>0\). The upper endpoints that lie in no
\(\relint F_j\) lie in \(\bigcup_{j\ne k}F_j\cap F_k\), which is a subset of a
finite union of affine subspaces of dimension at most \(n-2\)
(Fact~\ref{tb:facets}(b)); its projection along \(e\) onto \(e^\perp\) is a
Lebesgue null set of \(e^\perp\). For fixed \(j\) with \(e\cdot u_j>0\), the
map \(x\mapsto x-(x\cdot e)e\) is injective on the set of upper endpoints in
\(\relint F_j\), since a line meets \(Q\) in one interval, and it is the
restriction to \(H_j\) of a linear map whose Jacobian with respect to
\(\Haus^{n-1}\) on \(H_j\) is \(e\cdot u_j\): if \(b_1,\ldots,b_{n-1}\) is an
orthonormal basis of \(u_j^\perp\), the volume of the projection along \(e\)
of the cube they span equals the volume of the parallelepiped spanned by
\(b_1,\ldots,b_{n-1},e\), which is \(|e\cdot u_j|\). Therefore
\(\int f(w+s_+(w)e)\,dw=\sum_{j:\,e\cdot u_j>0}(e\cdot u_j)\int_{F_j}f\,d\Haus^{n-1}\),
the sets \(F_j\setminus\relint F_j\) being \(\Haus^{n-1}\)-null. The same
argument for the lower endpoints gives
\(\int f(w+s_-(w)e)\,dw
=\sum_{j:\,e\cdot u_j<0}|e\cdot u_j|\int_{F_j}f\,d\Haus^{n-1}\).
Subtracting this lower-endpoint integral from the upper-endpoint
integral proves \((*)\). Facets with \(e\cdot u_j=0\) contribute
nothing. Summing over coordinate directions proves the compactly
supported vector-field formula.

For the stated field \(X\), choose smooth cutoffs \(\chi_R\) equal to
one on the ball of radius \(R\), zero outside the ball of radius
\(2R\), with \(0\le\chi_R\le1\) and
\(\norm{\nabla\chi_R}\le C/R\). Apply the formula to \(\chi_RX\).
The product rule gives
\[
 \operatorname{div}(\chi_RX)
 =\chi_R\operatorname{div}X+\nabla\chi_R\cdot X.
\]
The first term converges in integral by the assumed integrable bound;
the second has integral in absolute value at most
\((C/R)\int_{\{\norm x\ge R\}}\norm{X(x)}\,dx\to0\).
On each of the finitely many facets, polynomial times Gaussian decay
is integrable, so dominated convergence also applies to the boundary
integrals. This proves the formula without requiring separate bounds
on the component derivatives of \(X\).
\end{proof}

\usedin{the divergence calculation inside a facet in
\PR{Lemma}{appx-an:lem:surface}{C.5}. The vector field is
\(X(y)\coloneqq y\varphi_{r-1}(y)\), with divergence
\(\{r-1-\norm y^2\}\varphi_{r-1}(y)\); the boundary terms are the
ridge contributions in \PE{appx-an:surface-identity}{C.30}.}
\emph{Source.} This is the classical Gauss--Green formula specialized
to a polyhedron; see \TBc{EG15}, Chapter~5. The cutoff argument above
extends the compact-support formula to the stated Gaussian-decay
hypothesis on an unbounded polyhedron.

For the surface integral used in the paper, this divergence
calculation produces coefficients whose separate expressions contain
\(1/\sin\alpha_{jk}\). We combine the two contributions of a ridge
before estimating them. Let
\(s_j\coloneqq \int_{F_j}\varphi_r\,d\Haus^{r-1}\),
\(\rho_{jk}\coloneqq \int_{E_{jk}}\varphi_r\,d\Haus^{r-2}\), and let
\(z_{jk}\) be the projection of a ridge point onto
\(\operatorname{span}(u_j,u_k)\). Applying the divergence formula in
each facet gives
\[
 \surf(Q)=\sum_j(a_j^2-1)s_j
       -\sum_{\{j,k\}}(d_{k\mid j}+d_{j\mid k})\rho_{jk}.
\]
The offsets \(d_{k\mid j}\) from Fact~\ref{tb:facets}(c) satisfy
\[
 d_{k\mid j}+d_{j\mid k}
 =2\sin(\alpha_{jk}/2)\,
       \frac{u_j+u_k}{\norm{u_j+u_k}}\cdot z_{jk},\qquad
 |d_{k\mid j}+d_{j\mid k}|\le\alpha_{jk}\norm{z_{jk}}.
\]
Indeed, their sum is
\((1-\cos\alpha_{jk})(a_j+a_k)/\sin\alpha_{jk}\), and
\(a_j+a_k=(u_j+u_k)\cdot z_{jk}\).
The inverse sine cancels. The remaining factor \(\alpha_{jk}\)
is precisely the angle of the normal sector whose mass we use next.

\subsection{Tubes around a polyhedron}\label{tb:sec:tubes}

We now compare the paired boundary terms with Gaussian mass outside
the polyhedron. The Gaussian density changes rapidly when a boundary point has a large
normal projection. We therefore shorten the outward tube at that point.
Its width may vary along a face: the required integration formula still
follows from orthogonal coordinates and Fubini's theorem.

Let \(Q\) be a proper, full-dimensional polyhedron in \(\R^n\) with
facets \(F_j\), ridges \(E_{jk}\), and an irredundant description by
unit face normals \(u_j\). For a ridge, put
\(L_{jk}\coloneqq \operatorname{span}(u_j,u_k)\) and
\(z_{jk}\coloneqq \Pi_{jk}x\), for any \(x\in E_{jk}\). This vector is constant
on the ridge because
\(u_j\cdot z_{jk}=a_j\) and \(u_k\cdot z_{jk}=a_k\).
Let \(\eta(x)>0\) be a measurable function on the facets and ridges.
We use the same value of \(\eta\) at a point whenever it belongs to
more than one closed face; only relative interiors enter the tubes.
The outward tubes are
\begin{align*}
 \mathcal T_j&\coloneqq \{x+s u_j:x\in\relint F_j,\ 0<s<\eta(x)\},\\
 \mathcal T_{jk}&\coloneqq \{x+v:x\in\relint E_{jk},\ v\in N_Q(x),\
                         0<\norm v<\eta(x)\}.
\end{align*}
For \(x\in\relint E_{jk}\), the normal cone is a sector of angle
\(\alpha_{jk}\) in \(L_{jk}\). Write
\(S_{jk}(a)\coloneqq N_Q(x)\cap\{v:0<\norm v<a\}\); its area is
\(\alpha_{jk}a^2/2\). For \(Q=\bigcap_j\{u_j\cdot x\le a_j\}\),
the notation \(Q^s\) increases every threshold by \(s\); it does not
denote a Euclidean neighborhood.

\begin{fact}[Variable-width normal tubes]\label{tb:tubes}
(a) The tubes are measurable, and tubes over distinct facets and
ridges are pairwise disjoint.
Each tube point \(g=x+v\) has nearest point \(x\) in \(Q\), and
\(\dist(g,Q)=\norm v\).

(b) For a nonnegative measurable function \(f\),
\begin{align*}
 \int_{\mathcal T_j}f(g)\,dg
 &=\int_{F_j}\int_0^{\eta(x)} f(x+s u_j)\,ds\,d\Haus^{n-1}(x),\\
 \int_{\mathcal T_{jk}}f(g)\,dg
 &=\int_{E_{jk}}\int_{S_{jk}(\eta(x))}f(x+v)\,dv\,d\Haus^{n-2}(x).
\end{align*}
The same formulas hold whenever \(f\) is integrable on the relevant tube.

(c) Use the function \(R\) in Fact~\ref{tb:projection-moments}, built
from the face normals of \(Q\), and choose \(\eta(x)\coloneqq 1/(4R(x))\).
Then all these tubes lie in \(\{g:g\in Q^{1/R(g)}\setminus Q\}\).
For an absolute constant \(c>0\), their normal integrals obey
\begin{align*}
 \int_0^{\eta(x)}R(x+s u_j)^3\varphi_n(x+s u_j)\,ds
 &\ge c\{1+a_j^2\}\varphi_n(x),\qquad x\in F_j,\\
 \int_{S_{jk}(\eta(x))}R(x+v)^3\varphi_n(x+v)\,dv
 &\ge c\alpha_{jk}\norm{z_{jk}}\varphi_n(x),\qquad x\in E_{jk}.
\end{align*}
For \(n=1\), there are no ridge tubes.
\end{fact}

\begin{proof}
(a) Facts~\ref{tb:normalcone} and~\ref{tb:projection} identify \(x\)
as the unique nearest point of \(x+v\). The relative interiors of
distinct faces are disjoint, so the tubes are disjoint.

(b) The map \((x,s)\mapsto x+s u_j\) is an affine isometry from
\(H_j\times\R\) onto \(\R^n\). Likewise, \((x,v)\mapsto x+v\)
is an affine isometry from \((z_{jk}+L_{jk}^\perp)\times L_{jk}\)
onto \(\R^n\). The domains are measurable in these coordinates because \(\eta\)
is measurable; their images under the affine isometries are therefore
measurable as well. The Jacobians are one. Apply Fubini to the indicator
of the domain \(0<s<\eta(x)\), or \(0<\norm v<\eta(x)\) in the
normal sector. No derivative of \(\eta\) is involved, and no
smoothness of \(\eta\) is needed.

(c) For a facet displacement \(v=s u_j\), and for a ridge displacement
\(v\in L_{jk}\), respectively,
\[
 |x\cdot v|\le\mathcal M(x)\norm v\le\tfrac14.
\]
The ridge case uses \(x\cdot v=z_{jk}\cdot v\). Since
\(\norm v\le1/4\), Fact~\ref{tb:gauss-coord}(b) gives
\(\varphi_n(x+v)\ge e^{-1/4-1/32}\varphi_n(x)\).
The Lipschitz property and \(R\ge1\) also give
\[
 \tfrac34R(x)\le R(x+v)\le\tfrac54R(x).
\]
The facet integral is therefore at least
\(cR(x)^3\eta(x)\varphi_n(x)=c'R(x)^2\varphi_n(x)\), which
dominates \(c''(1+a_j^2)\varphi_n(x)\) because
\(R(x)\ge1+|u_j\cdot x|=1+|a_j|\). The ridge integral is at least
\(c\alpha_{jk}R(x)^3\eta(x)^2\varphi_n(x)
=c'\alpha_{jk}R(x)\varphi_n(x)\); use
\(R(x)\ge\norm{z_{jk}}\).

Finally, a tube point \(g=x+v\) lies outside \(Q\), and each
unit-normal inequality defining \(Q\) satisfies
\[
 u_j\cdot g-a_j\le\norm v<\frac1{4R(x)}\le\frac1{R(g)}.
\]
This proves the inclusion in (c), even if the threshold system used
to define \(Q^s\) also contains redundant inequalities.
\end{proof}

Together with the paired ridge coefficients derived above, part (c)
gives
\[
 |\surf(Q)|\le C\E\left[R(G)^3
              \ind\{G\in Q^{1/R(G)}\setminus Q\}\right].
\]
The power three serves both dimensions: integration along a facet
uses one power of the tube width, leaving \(R^2\), while integration
over a ridge sector uses two, leaving \(R\).

\usedin{\PE{appx-an:adaptive-tubes}{C.29} in
\PR{Lemma}{appx-an:lem:surface}{C.5}. After setting
\(Q'\coloneqq (Q-V_0)/\sigma\), the event is
\(V_0+\sigma G\in Q^{\sigma/R(G)}\setminus Q\).
Conditional on \(G\), its width is fixed; concentration of \(V_0\)
and Fact~\ref{tb:projection-moments} give the averaged surface bound.}
\emph{Source.} The facet construction is Nazarov's outward-normal
argument~\TBc{Naz03}; see the detailed proof in \TBc{CCK17N},
Lemmas~1 and~2. Disjointness over arbitrary faces follows from the
classical projection characterization in \TBc{RW98}, Proposition~6.17.
The variable widths, ridge-sector integrals, and weighted estimates
needed here are derived above. For the related use of facet and ridge
tubes in Gaussian approximation, see \TBc{10}.

\subsection{The coarea formula and the first variation of a Gaussian
integral}\label{tb:sec:coarea}

The logistic mixture moves all thresholds together. To differentiate
its Gaussian integral, we integrate first over the level sets of the
maximum constraint. The coarea formula provides this change of variables.

\begin{fact}[Coarea formula]\label{tb:coarea}
Let \(\psi:\R^n\to\R\) be Lipschitz, so that \(\nabla\psi\) exists almost
everywhere. For every measurable \(f\ge0\),
\[
 \int_{\R^n}f\,\norm{\nabla\psi}\,dx=\int_{-\infty}^\infty
 \Bigl(\int_{\{\psi=s\}}f\,d\Haus^{n-1}\Bigr)ds ,
\]
For signed \(f\), the same identity holds if
\(\int |f|\norm{\nabla\psi}\,dx<\infty\), in particular if
\(f\in L^1(\R^n)\), since \(\psi\) is Lipschitz.
\end{fact}

\emph{Source.} \TBc{EG15}, Section~3.4. The proofs use the formula only for
a maximum of finitely many affine functions with gradients of equal length, and
for these Fact~\ref{tb:variation} gives a direct proof from Fubini's theorem.

For the finite maximum used here, the coarea formula can be proved
directly by partitioning the space according to the maximizing face.
This proof also identifies the exceptional levels and explains why the
independent Gamma shift permits differentiation at every fixed threshold.

\begin{fact}[First variation of a polyhedral sublevel integral]\label{tb:variation}
Let \(u_1,\ldots,u_m\) be unit vectors in \(\R^n\) (\(m\ge1\)), \(\kappa>0\),
\(b_1,\ldots,b_m\in\R\), and
\[
 \psi(g)\coloneqq \max_{j\le m}\{\kappa\,u_j\cdot g-b_j\},\qquad Q_s\coloneqq \{g:\psi(g)\le s\}.
\]
(a) For \(f\ge0\) measurable, or \(f\) integrable, and \(s\in\R\),
\[
 \int_{Q_s}f\,dg=\kappa^{-1}\int_{-\infty}^s\Bigl(\int_{\{\psi=s'\}}f\,d\Haus^{n-1}
 \Bigr)ds'.
\]
(b) If \(f\) is integrable, then \(\Lambda(s)\coloneqq \int_{Q_s}f\,dg\) is absolutely
continuous, and for almost every \(s\)
\[
 \Lambda'(s)=\kappa^{-1}\int_{\{\psi=s\}}f\,d\Haus^{n-1},\qquad
 |\Lambda'(s)|\le\kappa^{-1}\sum_{j=1}^m\int_{\{\kappa u_j\cdot g-b_j=s\}}|f|\,d\Haus^{n-1}.
\]
(c) If \(s>\inf\psi\), then \(Q_s\) is a proper, full-dimensional polyhedron
with face normals \(u_j\) and offsets \((s+b_j)/\kappa\), \(\{\psi=s\}=\partial Q_s\),
and \(\int_{\{\psi=s\}}f\,d\Haus^{n-1}\) is the surface integral of \(f\) over
\(\partial Q_s\) of Subsection~\ref{tb:sec:surface}. For \(s<\inf\psi\) the set
\(Q_s\) is empty, and \(s=\inf\psi\) is the only other level at which \(Q_s\)
may fail to be full-dimensional.

(d) For \(\Lambda\) from (b), assume in addition that it is
Lipschitz. If \(S\) is a real random variable with a density, then \(s\mapsto\E\Lambda(s+S)\) is differentiable with derivative
\(\E\Lambda'(s+S)\), where \(\Lambda'\) is defined almost everywhere.
\end{fact}

\begin{proof}
(a) Let \(C_j\) be the set where \(j\) is the smallest index attaining the
maximum defining \(\psi\); the Borel sets \(C_1,\ldots,C_m\) partition
\(\R^n\). On \(C_j\), \(\psi(g)=\kappa v-b_j\) in the coordinates
\(g=v u_j+y\), \(y\perp u_j\). By Fact~\ref{tb:fubini}, in these coordinates,
\[
 \int_{C_j\cap Q_s}f\,dg=\int_{u_j^\perp}\int_{-\infty}^{(s+b_j)/\kappa}
 (f\ind_{C_j})(v u_j+y)\,dv\,dy .
\]
The substitution \(s'=\kappa v-b_j\) turns the inner integral into
\(\kappa^{-1}\int_{-\infty}^s(f\ind_{C_j})(v u_j+y)\,ds'\) with
\(v=(s'+b_j)/\kappa\), and a further application of Fubini's theorem
exchanges the integrals. For fixed \(s'\), the set of \(v u_j+y\) in \(C_j\)
with this \(v\) is the slice \(C_j\cap\{\psi=s'\}\), and on the hyperplane
containing it \(\Haus^{n-1}\) is Lebesgue measure \(dy\)
(Fact~\ref{tb:hausdorff}(a)). Hence
\(\int_{C_j\cap Q_s}f\,dg=\kappa^{-1}\int_{-\infty}^s\int_{C_j\cap\{\psi=s'\}}
f\,d\Haus^{n-1}\,ds'\). Summing over \(j\) gives (a), since the slices
\(C_j\cap\{\psi=s'\}\) partition \(\{\psi=s'\}\).

(b) Applying (a) to \(|f|\) and letting \(s\to\infty\) shows that
\(I(s')\coloneqq \kappa^{-1}\int_{\{\psi=s'\}}f\,d\Haus^{n-1}\) is integrable on
\(\R\), and (a) says \(\Lambda(s)=\int_{-\infty}^sI(s')\,ds'\). Hence \(\Lambda\) is
absolutely continuous, and \(\Lambda'=I\) almost everywhere by Lebesgue's
differentiation theorem. The bound follows because
\(\{\psi=s\}\subseteq\bigcup_j\{\kappa u_j\cdot g-b_j=s\}\), a union of \(m\)
hyperplanes.

(c) The set \(Q_s\) is the polyhedron \(\bigcap_j\{u_j\cdot g\le(s+b_j)/\kappa\}\),
and \(\psi\) is convex and continuous. If \(\psi(g)=s>\inf\psi\), choose \(g'\)
with \(\psi(g')<s\); for \(\tau>0\), \(g=(g_\tau+\tau g')/(1+\tau)\) with
\(g_\tau\coloneqq g+\tau(g-g')\), so convexity gives
\(\psi(g_\tau)\ge s+\tau(s-\psi(g'))>s\). Thus \(g\) is a limit of points
outside \(Q_s\) and \(g\in\partial Q_s\); the open set \(\{\psi<s\}\) is
nonempty and lies in \(Q_s\), which is therefore full-dimensional, and it
is a proper subset because \(\psi(g)\ge\kappa u_1\cdot g-b_1\to\infty\) along
\(u_1\). A point with \(\psi<s\) is interior, so \(\partial Q_s=\{\psi=s\}\). After
the redundant and repeated inequalities are deleted, \(\partial Q_s=\{\psi=s\}\)
is the union of the sets \(Q_s\cap\{\kappa u_j\cdot g-b_j=s\}\) of the
remaining inequalities, which are faces of \(Q_s\) or empty. A face of
dimension \(n-1\) is a facet (Fact~\ref{tb:facets}(a)) and any other face is
\(\Haus^{n-1}\)-null, so the integral over \(\{\psi=s\}\) is the sum over the
facets of the irredundant description.

(d) The function \(\Lambda\) is bounded by \(\int|f|\), so the expectation
exists. A Lipschitz function is absolutely continuous and differentiable outside a
Lebesgue null set \(\mathcal N\). The random variable \(s+S\) avoids
\(\mathcal N\) almost surely, because \(S\) has a density, so the difference
quotients of \(\Lambda(s+S)\) converge almost surely and are bounded by the
Lipschitz constant. Dominated convergence gives the claim.
\end{proof}

\usedin{the first-variation paragraph in the proof of
\PR{Lemma}{appx-an:lem:matrix}{C.4}, immediately after
\PE{appx-an:gamma-mixture}{C.25}. Here
\(\psi(g)=\max_j\{u_j\cdot(V_0+\sigma g)-a_j\}\),
\(\kappa=\sigma\), \(b_j=a_j-u_j\cdot V_0\), and
\(f(g)=(\norm g^2-r)\varphi_r(g)\). Parts (a)--(c) give
\(\sigma^{-1}\surf((Q(a+s\mathbf1)-V_0)/\sigma)\) for almost every
\(s\). The independent shift \(-h\log T_\ell\), with
\(T_\ell\sim\operatorname{Gamma}(\ell,1)\), has a density, so part (d)
permits averaging these derivatives. A uniform bound
\(Cm(r+1)/\sigma\) follows from (b) and the Gaussian integral over
each supporting hyperplane; it justifies the interchange and does
not enter the final dimension-free bound.}

\clearpage
\subsection{Where each fact is used}\label{tb:sec:index}

The table follows the order of the proof. Each row identifies the
operation for which the corresponding facts are needed.

\medskip
\begin{center}
\small
\renewcommand{\arraystretch}{1.18}
\begin{tabular}{@{}p{0.27\textwidth}p{0.68\textwidth}@{}}
Facts~\ref{tb:fubini}, \ref{tb:taylor} &
Conditioning, Taylor remainders, and exchange of expectations with
convergent series in Appendices B--C.\\
Facts~\ref{tb:bernoulli}, \ref{tb:heat} &
Mask differentiation, the heat equation, and the signed Hessian
comparison in Subsection~\PN{appx-an:sec:path}{C.2}.\\
Facts~\ref{tb:faa}, \ref{tb:logistic}, \ref{tb:covariance} &
Derivative envelopes in Appendix~\PN{appx-sec:smoothing}{B}; early
bounds in \PN{appx-an:sec:early}{C.3} and coefficient matrices in
\PR{Lemma}{appx-an:lem:restoration}{C.2}.\\
Facts~\ref{tb:stein}, \ref{tb:nuclear} &
Two Gaussian integrations by parts and matrix contraction in
\PR{Lemma}{appx-an:lem:restoration}{C.2}.\\
Facts~\ref{tb:gauss-coord}, \ref{tb:logsum} &
Reference-face mean bounds and summation of reference masses in
\PR{Lemmas}{appx-an:lem:mean}{C.3} and~\PN{appx-an:lem:matrix}{C.4}.\\
Facts~\ref{tb:prekopa}--\ref{tb:lc-weights} &
Logistic-tilt covariance contraction in
\PR{Lemma}{appx-an:lem:matrix}{C.4}.\\
Facts~\ref{tb:facets}--\ref{tb:divergence} &
Surface measure and the facet--ridge identity in
\PR{Lemma}{appx-an:lem:surface}{C.5}.\\
Facts~\ref{tb:tubes}, \ref{tb:projection-moments} &
Adaptive normal tubes and Gaussian projection moments in
\PR{Lemma}{appx-an:lem:surface}{C.5}.\\
Facts~\ref{tb:coarea}, \ref{tb:variation} &
First variation of the Gamma mixture in
\PR{Lemma}{appx-an:lem:matrix}{C.4}.\\
\end{tabular}
\end{center}

\subsection{Sources}\label{tb:sec:sources}

\begin{description}\footnotesize\setlength{\itemsep}{0pt}\setlength{\parskip}{0pt}
\item[{[BL76]}] H.~J. Brascamp and E.~H. Lieb (1976). On extensions of the
Brunn--Minkowski and Pr\'ekopa--Leindler theorems, including inequalities
for log concave functions, and with an application to the diffusion
equation. \emph{Journal of Functional Analysis} \textbf{22}, 366--389.
\item[{[BV04]}] S. Boyd and L. Vandenberghe (2004).
\emph{Convex Optimization}. Cambridge University Press, Cambridge.
\item[{[CCK17N]}] V. Chernozhukov, D. Chetverikov, and K. Kato (2017).
Detailed proof of Nazarov's inequality. arXiv:1711.10696.
\item[{[Com74]}] L. Comtet (1974). \emph{Advanced Combinatorics}. Reidel,
Dordrecht.
\item[{[CT06]}] T.~M. Cover and J.~A. Thomas (2006). \emph{Elements of
Information Theory}, 2nd edition. Wiley.
\item[{[EG15]}] L.~C. Evans and R.~F. Gariepy (2015). \emph{Measure Theory
and Fine Properties of Functions}, revised edition. CRC Press.
\item[{[Fol99]}] G.~B. Folland (1999). \emph{Real Analysis: Modern
Techniques and Their Applications}, 2nd edition. Wiley.
\item[{[Gar02]}] R.~J. Gardner (2002). The Brunn--Minkowski inequality.
\emph{Bulletin of the American Mathematical Society} \textbf{39}, 355--405.
\item[{[10]}] Z. Gao and W. Zhang (2026). Cubic-root Gaussian approximation
under unrestricted covariance. arXiv:2608.30221v1.
\item[{[Lei72]}] L. Leindler (1972). On a certain converse of H\"older's
inequality II. \emph{Acta Scientiarum Mathematicarum (Szeged)} \textbf{33},
217--223.
\item[{[Naz03]}] F. Nazarov (2003). On the maximal perimeter of a convex
set in \(\R^n\) with respect to a Gaussian measure. In V.~D. Milman
and G. Schechtman (eds.), \emph{Geometric Aspects of Functional Analysis},
Lecture Notes in Mathematics \textbf{1807}, 169--187. Springer.
\item[{[Pre71]}] A. Pr\'ekopa (1971). Logarithmic concave measures with
application to stochastic programming. \emph{Acta Scientiarum Mathematicarum
(Szeged)} \textbf{32}, 301--316.
\item[{[Pre73]}] A. Pr\'ekopa (1973). On logarithmic concave measures and
functions. \emph{Acta Scientiarum Mathematicarum (Szeged)} \textbf{34},
335--343.
\item[{[Roc70]}] R.~T. Rockafellar (1970). \emph{Convex Analysis}. Princeton
University Press.
\item[{[RW98]}] R.~T. Rockafellar and R.~J.-B. Wets (1998). \emph{Variational
Analysis}. Springer, Berlin.
\item[{[Ste81]}] C.~M. Stein (1981). Estimation of the mean of a multivariate
normal distribution. \emph{The Annals of Statistics} \textbf{9}, 1135--1151.
\item[{[Zie95]}] G.~M. Ziegler (1995). \emph{Lectures on Polytopes}.
Springer, New York.
\end{description}


